% Le tourisme au Cameroun — Géographie Tle E — Cours signé OnBuch+
% Programme officiel MINESEC. Compiler avec : tectonic N10-tourisme-cameroun.tex
\newcommand{\DOCMATIERE}{Géographie}
\newcommand{\DOCNIVEAU}{Tle E}
\newcommand{\DOCMODULE}{Module 1 : Le Cameroun}
\newcommand{\DOCLECON}{N10}
\newcommand{\DOCTITRE}{Le tourisme au Cameroun}
% OnBuch+ — préambule « cours signés » ENRICHI (compilé avec tectonic / XeLaTeX)
% Système visuel « Pop » d'OnBuch+ : fond crème (jamais blanc), bordures encre
% épaisses, ombres « dures » décalées, titres Archivo Black en capitales.
% Version MATHÉMATIQUES : graphes (pgfplots/tikz), boîtes pédagogiques
% (définition, propriété, méthode, attention, le savais-tu, exercice résolu,
% exercice, TP, bilan), page de garde et signature OnBuch+ ancrée en bas.
%
% Macros attendues dans main.tex AVANT \input{preamble} :
%   \DOCMATIERE  \DOCNIVEAU  \DOCMODULE  \DOCLECON (numéro)  \DOCTITRE
\documentclass[11pt]{article}

\usepackage{fontspec}
\usepackage[french]{babel}
\usepackage[a4paper,margin=2cm,top=3.4cm,bottom=2.8cm,headheight=1.4cm,footskip=1.3cm]{geometry}
\usepackage{amsmath,amssymb,amsfonts}
\usepackage{mathtools} % \xmapsto, \coloneqq... souvent utilisés par les modèles
\usepackage{cancel} % simplifications barrées (bilans de forces)
% \abs{z} (module, valeur absolue) et \norm{u} (norme), utilisés par les modèles
\providecommand{\abs}{}\let\abs\relax\DeclarePairedDelimiter{\abs}{\lvert}{\rvert}
\providecommand{\norm}{}\let\norm\relax\DeclarePairedDelimiter{\norm}{\lVert}{\rVert}
\usepackage[table]{xcolor} % [table] : fournit aussi \rowcolors (alternance de lignes)
\usepackage{pagecolor}
\usepackage{tikz}
\usetikzlibrary{arrows.meta,positioning,calc,shapes.geometric,shapes.misc,shapes.arrows,decorations.pathmorphing,decorations.pathreplacing,decorations.markings,patterns,shadows,mindmap,trees,fit,backgrounds,matrix,intersections,angles,quotes}
\usepackage{pgfplots}
\usepackage[version=4]{mhchem} % équations nucléaires \ce{^{235}_{92}U}
\usepackage[european,siunitx]{circuitikz} % schémas électriques
\pgfplotsset{compat=1.18}
\usepgfplotslibrary{fillbetween,groupplots} % aires entre courbes (calcul intégral)
\usepackage{enumitem}
\usepackage{fancyhdr}
% [most] charge toutes les bibliothèques tcolorbox (listings, minted, raster,
% documentation...) et gonfle inutilement la mémoire TeX sur un document long
% (~300 boîtes) : on ne charge que ce qui est réellement utilisé.
\usepackage[skins,breakable]{tcolorbox}
\usepackage{graphicx}
\graphicspath{{/home/user/t-t/geo-onbuch/images/}}
\usepackage{colortbl}
\usepackage{booktabs}
\usepackage{tabularx}
\usepackage{multirow}
\usepackage{diagbox} % case d'en-tête diagonale des tableaux à double entrée
% Colonnes extensibles centrée / gauche / droite (C, L, R), souvent utilisées par les modèles
\newcolumntype{C}{>{\centering\arraybackslash}X}
\newcolumntype{L}{>{\raggedright\arraybackslash}X}
\newcolumntype{R}{>{\raggedleft\arraybackslash}X}
\usepackage{array}
\usepackage{multicol}
\usepackage{siunitx}
% parse-numbers=false : accepte aussi \SI{2,5 \times 10^{-3}}{...}
\sisetup{locale=FR,output-decimal-marker={,},group-minimum-digits=4,parse-numbers=false}
\DeclareSIUnit\ohm{\ensuremath{\symup{\Omega}}} % U+2126 absent de la police
\DeclareSIUnit\division{div} % réglages d'oscilloscope (ms/div, V/div)
\DeclareSIUnit\tr{tr} % tours (vitesse de rotation, ex. \SI{1500}{\tr\per\minute})
\DeclareSIUnit\molar{M} % concentration molaire (ex. \SI{3}{\molar}, \SI{1}{\micro\molar})
\DeclareSIUnit\an{an} % année (le modèle confond parfois avec \year, un compteur TeX) — ex. \SI{5}{\kilo\meter\per\an}
\DeclareSIUnit\FCFA{FCFA} % franc CFA (ex. \SI{500}{\FCFA\per\kilo\gram})
\DeclareSIUnit\degreeBrix{\textdegree Brix} % teneur en sucre d'un sirop/jus (réfractométrie)
\DeclareSIUnit\month{mois} % le modèle utilise parfois \month, un compteur TeX, comme unité de durée
\DeclareSIUnit\microliter{\micro\liter} % le modèle écrit parfois \microliter en un seul token
\DeclareSIUnit\million{millions} % dénombrement (ex. \SI{15}{\million\per\mL} spermatozoïdes)
\usepackage{qrcode}
% \degree : commande fréquemment utilisée par les modèles pour un angle ou une
% température en dehors de \SI{}{} (non fournie par les paquets chargés).
\providecommand{\degree}{^{\circ}}
% \qed (fin de démonstration), utilisé par les modèles sans amsthm
\providecommand{\qed}{\hfill$\square$}
% \barre : double barre des valeurs interdites dans un tableau de variations (tkz-tab)
\providecommand{\barre}{\ensuremath{\|}}
% environnement proof (amsthm non chargé) utilisé par les modèles
\ifdefined\proof\else\newenvironment{proof}[1][Démonstration]{\par\noindent\textit{#1.}\ }{\qed\par}\fi
\usepackage{lastpage}
\usepackage{hyperref}
\hypersetup{hidelinks}

% ============================================================
% Palette « Pop » OnBuch+ — mêmes hex que la classe P (plus_ui.dart)
% ============================================================
\definecolor{poporange}{HTML}{F59321}
\definecolor{popdark}{HTML}{E07A0C}
\definecolor{popcream}{HTML}{FFF6E8}
\definecolor{popink}{HTML}{1C1714}
\definecolor{poppurple}{HTML}{7A5AE0}
\definecolor{popgreen}{HTML}{2E9E6A}
\definecolor{poppink}{HTML}{E0457B}
\definecolor{popblue}{HTML}{2F7DE1}
\definecolor{popgold}{HTML}{C98A12}
\definecolor{popmuted}{HTML}{8A7F76}
\definecolor{popline}{HTML}{EADFD2}
\definecolor{popgray}{HTML}{8A7F76}
\colorlet{popgrey}{popgray}
% Alias tolérants (le modèle nomme parfois les couleurs différemment)
\colorlet{popwhite}{white}
\colorlet{popblack}{popink}
\colorlet{popred}{poppink}
\colorlet{popyellow}{popgold}
% Teintes claires (fonds de tableaux/encadrés) — le modèle nomme parfois
% les couleurs avec un suffixe L (ex. popblueL) sans les avoir définies.
\colorlet{poporangeL}{poporange!20}
\colorlet{popdarkL}{popdark!20}
\colorlet{poppurpleL}{poppurple!20}
\colorlet{popgreenL}{popgreen!20}
\colorlet{poppinkL}{poppink!20}
\colorlet{popblueL}{popblue!20}
\colorlet{popgoldL}{popgold!20}
\colorlet{popredL}{popred!20}
\colorlet{popgrayL}{popgray!20}
% Teintes claires pour fonds de tableaux / zones
\colorlet{popblueL}{popblue!10!white}
\colorlet{poporangeL}{poporange!14!white}
\colorlet{popgreenL}{popgreen!12!white}
\colorlet{poppurpleL}{poppurple!12!white}
\colorlet{poppinkL}{poppink!12!white}
\colorlet{popgoldL}{popgold!14!white}

\pagecolor{popcream}

% ============================================================
% Typographie
% ============================================================
\defaultfontfeatures{Ligatures=TeX}
\setmainfont{PlusJakartaSans-Medium.ttf}[Path=../../../fonts/,BoldFont=PlusJakartaSans-Bold.ttf]
% Maths en sans-serif (Fira Math) avec chiffres et lettres droites de Plus
% Jakarta, pour que les formules se fondent dans le texte.
\usepackage[math-style=ISO,bold-style=ISO]{unicode-math}
\setmathfont{FiraMath-Regular.otf}[Path=../../../fonts/]
\setmathfont{PlusJakartaSans-Medium.ttf}[Path=../../../fonts/,range={up/num,up/{Latin,latin},\mathup}]
\usepackage{icomma} % virgule décimale sans espace en mode math
% Caractères mathématiques Unicode absents des polices de texte (Plus Jakarta,
% Archivo Black) : rendus par leur équivalent mathématique, en texte comme en
% formule — sinon ils s'affichent en carré vide (ex. « Arithmétique dans ℤ »).
\usepackage{newunicodechar}
\newunicodechar{ℝ}{\ensuremath{\mathbb{R}}}
\newunicodechar{ℤ}{\ensuremath{\mathbb{Z}}}
\newunicodechar{ℂ}{\ensuremath{\mathbb{C}}}
\newunicodechar{ℕ}{\ensuremath{\mathbb{N}}}
\newunicodechar{ℚ}{\ensuremath{\mathbb{Q}}}
\newunicodechar{𝔻}{\ensuremath{\mathbb{D}}}
\newunicodechar{α}{\ensuremath{\alpha}}
\newunicodechar{β}{\ensuremath{\beta}}
\newunicodechar{γ}{\ensuremath{\gamma}}
\newunicodechar{δ}{\ensuremath{\delta}}
\newunicodechar{ε}{\ensuremath{\varepsilon}}
\newunicodechar{θ}{\ensuremath{\theta}}
\newunicodechar{λ}{\ensuremath{\lambda}}
\newunicodechar{μ}{\ensuremath{\mu}}
\newunicodechar{σ}{\ensuremath{\sigma}}
\newunicodechar{τ}{\ensuremath{\tau}}
\newunicodechar{φ}{\ensuremath{\varphi}}
\newunicodechar{ψ}{\ensuremath{\psi}}
\newunicodechar{ω}{\ensuremath{\omega}}
\newunicodechar{Σ}{\ensuremath{\Sigma}}
% majuscules produites par \MakeUppercase dans les titres (α-aminés -> Α)
\newunicodechar{Α}{\ensuremath{\alpha}}
\newunicodechar{Β}{\ensuremath{\beta}}
\newunicodechar{Γ}{\ensuremath{\Gamma}}
\newunicodechar{Δ}{\ensuremath{\Delta}}
\newunicodechar{Ε}{\ensuremath{\varepsilon}}
\newunicodechar{Θ}{\ensuremath{\Theta}}
\newunicodechar{Λ}{\ensuremath{\Lambda}}
\newunicodechar{Μ}{\ensuremath{\mu}}
\newunicodechar{Τ}{\ensuremath{\tau}}
\newunicodechar{Φ}{\ensuremath{\Phi}}
\newunicodechar{Ψ}{\ensuremath{\Psi}}
\newunicodechar{Ω}{\ensuremath{\Omega}}
\newunicodechar{↦}{\ensuremath{\mapsto}}
\newunicodechar{∈}{\ensuremath{\in}}
\newunicodechar{∉}{\ensuremath{\notin}}
\newunicodechar{∧}{\ensuremath{\wedge}}
\newunicodechar{⇔}{\ensuremath{\Leftrightarrow}}
\newunicodechar{⟺}{\ensuremath{\Longleftrightarrow}}
\newunicodechar{⇒}{\ensuremath{\Rightarrow}}
\newunicodechar{⟹}{\ensuremath{\Longrightarrow}}
\newunicodechar{∘}{\ensuremath{\circ}}
\newunicodechar{∪}{\ensuremath{\cup}}
\newunicodechar{∩}{\ensuremath{\cap}}
\newunicodechar{≡}{\ensuremath{\equiv}}
\newunicodechar{⊂}{\ensuremath{\subset}}
\newunicodechar{⊥}{\ensuremath{\perp}}
\newunicodechar{∃}{\ensuremath{\exists}}
\newunicodechar{∀}{\ensuremath{\forall}}
\newunicodechar{∛}{\ensuremath{\sqrt[3]{\,}}}
\newunicodechar{∜}{\ensuremath{\sqrt[4]{\,}}}
\newunicodechar{⃗}{\ensuremath{{}^{\to}}}
\newunicodechar{ⁿ}{\ifmmode^{n}\else\textsuperscript{\ensuremath{n}}\fi}
\newunicodechar{ˣ}{\ifmmode^{x}\else\textsuperscript{\ensuremath{x}}\fi}
\newunicodechar{ᵇ}{\ifmmode^{b}\else\textsuperscript{\ensuremath{b}}\fi}
\newunicodechar{ʳ}{\ifmmode^{r}\else\textsuperscript{\ensuremath{r}}\fi}
\newunicodechar{ᵘ}{\ifmmode^{u}\else\textsuperscript{\ensuremath{u}}\fi}
\newunicodechar{ʸ}{\ifmmode^{y}\else\textsuperscript{\ensuremath{y}}\fi}
\newunicodechar{ᵃ}{\ifmmode^{a}\else\textsuperscript{\ensuremath{a}}\fi}
\newunicodechar{ᵉ}{\ifmmode^{e}\else\textsuperscript{\ensuremath{e}}\fi}
\newunicodechar{ᵏ}{\ifmmode^{k}\else\textsuperscript{\ensuremath{k}}\fi}
\newunicodechar{ᵖ}{\ifmmode^{p}\else\textsuperscript{\ensuremath{p}}\fi}
\newunicodechar{ᴺ}{\ifmmode^{N}\else\textsuperscript{\ensuremath{N}}\fi}
\newunicodechar{ᶜ}{\ifmmode^{c}\else\textsuperscript{\ensuremath{c}}\fi}
\newunicodechar{ʰ}{\ifmmode^{h}\else\textsuperscript{\ensuremath{h}}\fi}
\newunicodechar{ᵠ}{\ifmmode^{q}\else\textsuperscript{\ensuremath{q}}\fi}
\newunicodechar{ₙ}{\ifmmode_{n}\else\textsubscript{\ensuremath{n}}\fi}
\newunicodechar{ₐ}{\ifmmode_{a}\else\textsubscript{\ensuremath{a}}\fi}
\newunicodechar{ᵢ}{\ifmmode_{i}\else\textsubscript{\ensuremath{i}}\fi}
\newunicodechar{ᵣ}{\ifmmode_{r}\else\textsubscript{\ensuremath{r}}\fi}
\newunicodechar{ₖ}{\ifmmode_{k}\else\textsubscript{\ensuremath{k}}\fi}
\newunicodechar{ᵦ}{\ifmmode_{\beta}\else\textsubscript{\ensuremath{\beta}}\fi}
\newunicodechar{ₓ}{\ifmmode_{x}\else\textsubscript{\ensuremath{x}}\fi}
\newfontfamily\popdisplay{ArchivoBlack-Regular.ttf}[Path=../../../fonts/]
\newfontfamily\poplogo{SpaceGrotesk-Bold.ttf}[Path=../../../fonts/]
\newfontfamily\popsemi{PlusJakartaSans-SemiBold.ttf}[Path=../../../fonts/]
\newfontfamily\popxbold{PlusJakartaSans-ExtraBold.ttf}[Path=../../../fonts/]
\newfontfamily\popxboldls{PlusJakartaSans-ExtraBold.ttf}[Path=../../../fonts/,LetterSpace=3.0]

\setlist[enumerate]{leftmargin=1.7em,itemsep=5pt,topsep=4pt}
\setlist[itemize]{leftmargin=1.7em,itemsep=4pt,topsep=4pt,label={\color{poporange}\textbullet}}
\setlength{\parindent}{0pt}
\setlength{\parskip}{5pt}
\renewcommand{\arraystretch}{1.25}

% pgfplots : style Pop par défaut
\pgfplotsset{
  every axis/.append style={axis line style={popink,line width=1pt},tick style={popink},
    grid style={popline,line width=0.5pt},label style={font=\small},tick label style={font=\footnotesize}},
  popcurve/.style={poporange,line width=1.8pt,no markers,smooth},
}
% Même style disponible dans un \draw TikZ simple (hors pgfplots)
\tikzset{popcurve/.style={poporange,line width=1.8pt}}
\tikzset{popgrid/.style={popline,very thin}}
\colorlet{popcurve}{poporange} % le nom du style est parfois utilisé comme couleur

% ============================================================
% Pill / squiggle
% ============================================================
\newcommand{\poppill}[3]{%
  \tikz[baseline=-0.55ex]\node[draw=popink,line width=1.2pt,rounded corners=2.6pt,%
    fill=#2,inner xsep=6.5pt,inner ysep=3pt]{%
    \popxboldls\fontsize{7}{7}\selectfont\color{#3}\MakeUppercase{#1}};%
}
\newcommand{\popsquiggle}[1]{%
  \tikz[baseline=-0.4ex]{
    \draw[#1,line width=2.6pt,line cap=round]
      plot [smooth] coordinates {(0,0.09) (0.34,0.01) (0.68,0.21) (1.02,0.01) (1.36,0.10)};
  }%
}
% Mot-clé surligné (marqueur orange sous le texte)
\newcommand{\cle}[1]{{\popxbold\color{popdark}#1}}

% ============================================================
% En-tête / pied
% ============================================================
\pagestyle{fancy}
\fancyhf{}
\renewcommand{\headrulewidth}{0pt}
\renewcommand{\footrulewidth}{0pt}
\lhead{\raisebox{-0.15ex}{\poplogo\fontsize{13}{13}\selectfont\color{popink}OnBuch\color{poporange}+}}
\rhead{\poppill{\DOCMATIERE\ \textbullet\ \DOCNIVEAU\ \textbullet\ Leçon \DOCLECON}{white}{popink}}
\lfoot{\popxbold\fontsize{7.6}{7.6}\selectfont\color{popmuted}\MakeUppercase{Cours signé OnBuch+}\ \textbullet\ {\popsemi\DOCTITRE}}
\rfoot{\tikz[baseline=-0.6ex]\node[rounded corners=8.5pt,fill=popink,minimum height=17pt,minimum width=17pt,inner xsep=5pt,inner sep=0]{\popxbold\fontsize{8}{8}\selectfont\color{popcream}\thepage\,/\,\pageref*{LastPage}};}

% ============================================================
% Titres
% ============================================================
\newcommand{\chaptitre}[2]{%
  \begin{tcolorbox}[enhanced,colback=poporange,colframe=popink,boxrule=2.6pt,arc=11pt,
      drop shadow=popink,left=15pt,right=15pt,top=12pt,bottom=11pt,before skip=3pt,after skip=18pt]
    \poppill{#2}{popink}{popcream}\par\vspace{9pt}
    {\raggedright\hyphenpenalty=10000\exhyphenpenalty=10000\popdisplay\fontsize{22}{24}\selectfont\color{popcream}\MakeUppercase{#1}\par}
  \end{tcolorbox}
}
% Section numérotée : pastille orange avec le numéro + titre Archivo
\newcounter{popsec}
\newcommand{\coursec}[1]{%
  \stepcounter{popsec}\setcounter{popsub}{0}%
  \par\vspace{16pt}\noindent
  \begin{minipage}{\linewidth}
  \tikz[baseline=-0.7ex]\node[draw=popink,line width=1.6pt,fill=poporange,rounded corners=4pt,minimum size=22pt,inner sep=0]{\popdisplay\fontsize{12}{12}\selectfont\color{popcream}\thepopsec};%
  \hspace{8pt}{\popdisplay\fontsize{15}{17}\selectfont\color{popink}\MakeUppercase{#1}}\par
  \vspace{4pt}\noindent{\color{poporange}\rule{36pt}{3.6pt}}
  \end{minipage}\par\nopagebreak\vspace{8pt}
}
\newcounter{popsub}
\newcommand{\courssub}[1]{%
  \stepcounter{popsub}%
  \par\vspace{9pt}\noindent{\popxbold\fontsize{11.5}{13}\selectfont\color{popink}{\color{poporange}\thepopsec.\thepopsub}\hspace{6pt}#1}\par\nopagebreak\vspace{4pt}
}

% ============================================================
% Boîtes pédagogiques — même recette : fond blanc, bordure épaisse colorée,
% ombre dure encre, badge plein. Titre optionnel [#1].
% ============================================================
\tcbset{popbase/.style={breakable,enhanced,colback=white,boxrule=2.3pt,arc=9pt,
  shadow={3pt}{-3pt}{0pt}{popink},left=14pt,right=14pt,top=11pt,bottom=11pt,before skip=11pt,after skip=15pt,
  fontupper=\popsemi\fontsize{10.6}{14.5}\selectfont\color{popink},
  fontlower=\popsemi\fontsize{10.6}{14.5}\selectfont\color{popink}}}
% Coche dessinée (le glyphe ✓ n'existe pas dans les polices du document)
\newcommand{\popcheck}[1]{\tikz[baseline=-0.5ex]\draw[#1,line width=1.6pt,line cap=round,line join=round](-0.13,0.0)--(-0.04,-0.09)--(0.13,0.1);}
\providecommand{\coche}{\popcheck{popgreen}}
\providecommand{\resultat}[1]{\par\smallskip\noindent\fcolorbox{popgreen}{popgreenL}{\parbox{\dimexpr\linewidth-2\fboxsep-2\fboxrule\relax}{\textbf{Résultat.} #1}}\par\smallskip}
\newcommand{\popboxhead}[3]{\poppill{#1}{#2}{white}\ifx\relax#3\relax\else\hspace{6pt}{\popxbold\fontsize{10.6}{12}\selectfont\color{popink}#3}\fi\par\vspace{7pt}}

\newtcolorbox{aretenir}[1][]{popbase,colframe=popgreen,before upper={\popboxhead{À retenir}{popgreen}{#1}}}
\newtcolorbox{exemplebox}[1][]{popbase,colframe=poporange,before upper={\popboxhead{Exemple}{poporange}{#1}}}
\newtcolorbox{definition}[1][]{popbase,colframe=popblue,before upper={\popboxhead{Définition}{popblue}{#1}}}
\newtcolorbox{propriete}[1][]{popbase,colframe=poppurple,before upper={\popboxhead{Propriété}{poppurple}{#1}}}
\newtcolorbox{methode}[1][]{popbase,colframe=popgold,before upper={\popboxhead{Méthode}{popgold}{#1}}}
\newtcolorbox{attention}[1][]{popbase,colframe=poppink,before upper={\popboxhead{Attention}{poppink}{#1}}}
\newtcolorbox{experience}[1][]{popbase,colframe=popblue,colback=popblueL,before upper={\popboxhead{Expérience}{popblue}{#1}}}
% « Le savais-tu ? » : carte violette pleine (contexte, culture, Cameroun)
\newtcolorbox{savaistu}[1][]{popbase,colback=poppurple,colframe=popink,
  fontupper=\popsemi\fontsize{10.4}{14.2}\selectfont\color{white},
  before upper={\poppill{Le savais-tu\,?}{white}{poppurple}\ifx\relax#1\relax\else\hspace{6pt}{\popxbold\color{white}#1}\fi\par\vspace{7pt}}}
% Exercice résolu : énoncé en haut, \tcblower puis la solution sur fond crème
\newcounter{popexo}\newcounter{popexr}
\newtcolorbox{exoresolu}[1][]{popbase,colframe=poporange,colbacklower=popcream,
  segmentation style={popink,line width=1.2pt,dashed},
  before upper={\stepcounter{popexr}\popboxhead{Exercice résolu \thepopexr}{poporange}{#1}},
  before lower={\poppill{Solution}{popink}{popcream}\par\vspace{6pt}}}
% Exercice (non résolu en place) — la correction est dans la partie Corrigés
\newtcolorbox{exercice}[1][]{popbase,colframe=popink,
  before upper={\stepcounter{popexo}\popboxhead{Exercice \thepopexo}{popink}{#1}}}
% Corrigé
\newtcolorbox{corrige}[1][]{popbase,colframe=popgreen,colback=popgreenL,
  before upper={\popboxhead{Corrigé}{popgreen}{#1}}}
% Objectifs / prérequis / bilan
\newtcolorbox{objectifs}{popbase,colframe=popink,colback=poporangeL,
  before upper={\popboxhead{Objectifs de la leçon}{popink}{}}}
\newtcolorbox{prerequis}{popbase,colframe=popink,colback=popblueL,
  before upper={\popboxhead{Prérequis}{popblue}{}}}
\newtcolorbox{bilan}{popbase,colframe=popink,colback=white,boxrule=2.8pt,
  before upper={\popboxhead{Fiche bilan}{poporange}{L'essentiel à connaître pour l'examen}}}
% Figure encadrée avec légende
\newtcolorbox{popfigure}[1][]{enhanced,breakable=false,colback=white,colframe=popline,boxrule=1.4pt,arc=8pt,
  left=10pt,right=10pt,top=10pt,bottom=8pt,before skip=10pt,after skip=12pt,halign=center,
  lower separated=false,halign lower=center,
  fontlower=\popsemi\fontsize{9}{11.5}\selectfont\color{popmuted},
  #1}
\newcommand{\legende}[1]{\par\vspace{6pt}{\popsemi\fontsize{9}{11.5}\selectfont\color{popmuted}#1}}

% ============================================================
% Signature OnBuch+
% ============================================================
\newcommand{\poplogoinline}[1]{{\poplogo\fontsize{#1}{#1}\selectfont\color{popink}OnBuch\color{poporange}+}}
\newcommand{\signatureonbuch}{%
  \begin{tcolorbox}[enhanced,colback=popink,colframe=popink,boxrule=2.2pt,arc=10pt,
      drop shadow=poporange,left=14pt,right=14pt,top=10pt,bottom=10pt,before skip=0pt,after skip=0pt]
    \tikz[baseline=-0.7ex]\node[circle,fill=popgreen,draw=popcream,line width=1.3pt,minimum size=22pt,inner sep=0]{\popcheck{white}};%
    \hspace{9pt}\begin{minipage}[c]{0.62\linewidth}
      {\popxbold\fontsize{12}{13}\selectfont\color{popcream}Signé\ \textbullet\ {\poplogo OnBuch\color{poporange}+}}\par\vspace{2pt}
      {\raggedright\popsemi\fontsize{8}{10.5}\selectfont\color{popline}\MakeUppercase{Équipe pédagogique OnBuch+}\\\MakeUppercase{Conforme au programme officiel MINESEC}\par}
    \end{minipage}\hfill
    {\poplogo\fontsize{22}{22}\selectfont\color{popcream}OnBuch\color{poporange}+}
  \end{tcolorbox}%
}

% ============================================================
% Page de garde
% #1 titre · #2 matière · #3 niveau · #4 module · #5 n° leçon · #6 durée ·
% #7 description / objectifs (texte ou itemize)
% ============================================================
\newcommand{\pagegarde}[7]{%
  \thispagestyle{empty}
  \newgeometry{a4paper,margin=1.7cm,top=1.6cm,bottom=1.4cm}%
  \begin{tikzpicture}[remember picture,overlay]
    \fill[poporange!18] ([xshift=-3.2cm,yshift=-3.6cm]current page.north east) circle (4.6cm);
    \fill[poppurple!14] ([xshift=2.4cm,yshift=7.6cm]current page.south west) circle (3.8cm);
  \end{tikzpicture}%
  {\poplogo\fontsize{30}{30}\selectfont\color{popink}OnBuch\color{poporange}+}\hfill
  \raisebox{0.6ex}{\poppill{Cours signé \textbullet\ Premium}{popink}{popcream}}\par
  \vspace{1pt}\popsquiggle{poppurple}\par
  \vspace{8pt}
  \begin{tcolorbox}[enhanced,colback=poporange,colframe=popink,boxrule=2.8pt,arc=13pt,
      drop shadow=popink,left=18pt,right=18pt,top=12pt,bottom=13pt,after skip=10pt]
    \poppill{#2 \textbullet\ #3}{popink}{popcream}\hspace{5pt}\poppill{#4}{white}{popink}\par\vspace{8pt}
    {\popdisplay\fontsize{14}{14}\selectfont\color{popink}LEÇON #5}\par\vspace{4pt}
    {\raggedright\hyphenpenalty=10000\exhyphenpenalty=10000\popdisplay\fontsize{25}{27}\selectfont\color{popcream}\MakeUppercase{#1}\par}
    \vspace{8pt}
    {\popxbold\fontsize{9.5}{9.5}\selectfont\color{popink}\MakeUppercase{Durée conseillée : #6}}
  \end{tcolorbox}
  \begin{tcolorbox}[enhanced,colback=white,colframe=popink,boxrule=2.2pt,arc=10pt,
      drop shadow=popink,left=16pt,right=16pt,top=10pt,bottom=10pt,after skip=8pt]
    \poppill{Dans cette leçon}{poppurple}{white}\par\vspace{5pt}
    \popsemi\fontsize{9.3}{12.6}\selectfont\color{popink}#7
  \end{tcolorbox}
  \noindent\begin{minipage}[t]{0.487\linewidth}
    \begin{tcolorbox}[enhanced,colback=popgreen,colframe=popink,boxrule=2.2pt,arc=10pt,
        drop shadow=popink,left=11pt,right=11pt,top=10pt,bottom=10pt,height=3.1cm,valign=center]
      \tcbox[colback=white,colframe=popink,boxrule=1.3pt,arc=5pt,boxsep=0pt,nobeforeafter,
        left=3pt,right=3pt,top=3pt,bottom=3pt,valign=center]{\qrcode[height=1.9cm]{https://play.google.com/store/apps/details?id=com.luvvix.onbuch&hl=fr}}%
      \hspace{8pt}\begin{minipage}[c]{0.5\linewidth}
        {\popdisplay\fontsize{11}{12.5}\selectfont\color{white}\MakeUppercase{Télécharge OnBuch}}\par\vspace{4pt}
        {\raggedright\popsemi\fontsize{8.4}{11}\selectfont\color{white}Gratuit \textbullet\ Google Play\\Tuteur IA, annales, concours, résultats\par}
      \end{minipage}
    \end{tcolorbox}
  \end{minipage}\hfill
  \begin{minipage}[t]{0.487\linewidth}
    \begin{tcolorbox}[enhanced,colback=poppurple,colframe=popink,boxrule=2.2pt,arc=10pt,
        drop shadow=popink,left=11pt,right=11pt,top=10pt,bottom=10pt,height=3.1cm,valign=center]
      \tikz[baseline=-0.7ex]\node[circle,fill=white,draw=popink,line width=1.3pt,minimum size=1.6cm,inner sep=0]{\popdisplay\fontsize{24}{24}\selectfont\color{poppurple}+};%
      \hspace{8pt}\begin{minipage}[c]{0.6\linewidth}
        {\popdisplay\fontsize{11}{12.5}\selectfont\color{white}\MakeUppercase{Passe à OnBuch+}}\par\vspace{4pt}
        {\raggedright\popsemi\fontsize{8.4}{11}\selectfont\color{white}Tous les cours signés, Tuteur IA illimité, fiches PDF — onglet OnBuch+ de l'app\par}
      \end{minipage}
    \end{tcolorbox}
  \end{minipage}
  \vspace{8pt}
  \popcontenu
  \vfill
  \signatureonbuch
  \restoregeometry
  \newpage
}

% Rangée « Ce cours contient » (page de garde)
\newcommand{\poptile}[3]{% #1 couleur #2 gros texte #3 légende
  \begin{tcolorbox}[enhanced,colback=#1,colframe=popink,boxrule=2pt,arc=9pt,shadow={2.5pt}{-2.5pt}{0pt}{popink},
      left=6pt,right=6pt,top=9pt,bottom=9pt,halign=center,valign=center,height=1.75cm,width=\linewidth,nobeforeafter]
    {\popdisplay\fontsize{12.5}{13}\selectfont\color{white}\MakeUppercase{#2}}\par\vspace{3pt}
    {\centering\popsemi\fontsize{7.8}{9.5}\selectfont\color{white}#3\par}
  \end{tcolorbox}}
\newcommand{\popcontenu}{%
  \par\noindent{\popxboldls\fontsize{7.5}{7.5}\selectfont\color{popmuted}CE COURS SIGNÉ CONTIENT}\par\vspace{7pt}
  \noindent\begin{minipage}[t]{0.235\linewidth}\poptile{popblue}{Cours}{complet, illustré, pas à pas}\end{minipage}\hfill
  \begin{minipage}[t]{0.235\linewidth}\poptile{poporange}{Méthodes}{et exercices résolus}\end{minipage}\hfill
  \begin{minipage}[t]{0.235\linewidth}\poptile{poppink}{Exercices}{type examen + corrigés}\end{minipage}\hfill
  \begin{minipage}[t]{0.235\linewidth}\poptile{popgreen}{Fiche bilan}{l'essentiel à retenir}\end{minipage}\par}

% Sommaire
\newtcolorbox{sommaire}{enhanced,colback=white,colframe=popink,boxrule=2.2pt,arc=10pt,
  drop shadow=popink,left=18pt,right=18pt,top=13pt,bottom=13pt,before skip=6pt,after skip=20pt,
  before upper={\poppill{Sommaire}{poppurple}{white}\par\vspace{9pt}\popsemi\fontsize{10.6}{14.5}\selectfont\color{popink}}}

% Signature de fin de cours — poussée tout en bas de la dernière page
\newcommand{\findecours}{%
  \par\vfill\nopagebreak
  \begin{center}\popsquiggle{poporange}\end{center}\vspace{4pt}
  \signatureonbuch
}
\AtBeginDocument{\def\llbracket{\mathopen{[\mkern-3.5mu[}}\def\rrbracket{\mathclose{]\mkern-3.5mu]}}} % intervalles d'entiers (glyphe ⟦ absent de la police)
% Glyphes absents de Fira Math : redéfinis à partir de glyphes disponibles
\AtBeginDocument{%
  \def\setminus{\mathbin{\backslash}}\let\smallsetminus\setminus
  \def\vdots{\mathord{\vbox{\baselineskip3.2pt\lineskiplimit0pt\kern4pt\hbox{.}\hbox{.}\hbox{.}}}}%
  \def\ddots{\mathinner{\mkern1mu\raise6.5pt\hbox{.}\mkern2mu\raise3.6pt\hbox{.}\mkern2mu\raise0.7pt\hbox{.}\mkern1mu}}%
  \def\triangle{\mathord{\scalebox{0.9}{$\symup{\Delta}$}}}%
  \def\nabla{\mathord{\raisebox{\depth}{\scalebox{1}[-1]{$\symup{\Delta}$}}}}%
  \def\checkmark{\popcheck{popdark}}%
  \def\star{\mathbin{\ast}}\def\bigstar{\mathord{\ast}}%
  \def\diamond{\mathbin{\scalebox{0.6}{$\Diamond$}}}%
  \def\bigcup{\mathop{\mathchoice{\raisebox{-0.3em}{\scalebox{1.7}{$\cup$}}}{\scalebox{1.25}{$\cup$}}{\cup}{\cup}}\displaylimits}%
  \def\bigcap{\mathop{\mathchoice{\raisebox{-0.3em}{\scalebox{1.7}{$\cap$}}}{\scalebox{1.25}{$\cap$}}{\cap}{\cap}}\displaylimits}%
  \def\lesseqqgtr{\mathrel{\lessgtr}}\def\bigoplus{\mathop{\oplus}\displaylimits}%
}
\providecommand{\ssi}{si et seulement si} % abréviation fréquente
\AtBeginDocument{\providecommand{\wideparen}{\overparen}} % arc de cercle

%% FIGFIX-BEGIN v5 — correctifs globaux des figures (docs/onbuch-generation/tools/figfix)
\usepackage{adjustbox}\usepackage{etoolbox}\tcbuselibrary{raster}
% toute figure TikZ/pgfplots plus large que la ligne est réduite à la largeur disponible
\BeforeBeginEnvironment{tikzpicture}{\begin{adjustbox}{max width=\linewidth}}
\AfterEndEnvironment{tikzpicture}{\end{adjustbox}}
% légendes pgfplots lisibles par-dessus les courbes
\pgfplotsset{every axis/.append style={legend style={fill=white,fill opacity=0.92,text opacity=1,draw=black!15,inner sep=3pt}}}
% étiquettes d'arêtes (syntaxe edge["..."]) avec fond blanc
\tikzset{every edge quotes/.append style={fill=white,inner sep=1.2pt}}
%% FIGFIX-END
\begin{document}

\pagegarde{Le tourisme au Cameroun}{Géographie}{Tle E}{Module 1 : Le Cameroun}{N10}{2 h}{Cette leçon étudie le tourisme au Cameroun : ses potentialités exceptionnelles réparties sur tout le territoire, son importance économique, politique et sociale, ainsi que les problèmes qui freinent son développement. Elle s'appuie sur l'analyse de cartes, de photographies et de données chiffrées, et sur la construction de croquis et de diagrammes.
\vspace{4pt}
\begin{multicols}{2}\small
\begin{itemize}
\item À la fin de cette leçon, je sais définir les notions de site touristique et de tourisme d'affaires
\item À la fin de cette leçon, je sais localiser sur une carte les principales zones et sites touristiques du Cameroun
\item À la fin de cette leçon, je sais décrire et classer les potentialités touristiques du Cameroun (balnéaire, safari, culturel, écotourisme)
\item À la fin de cette leçon, je sais expliquer l'importance économique, politique et sociale du tourisme au Cameroun
\item À la fin de cette leçon, je sais identifier et hiérarchiser les problèmes qui freinent le secteur touristique camerounais
\item À la fin de cette leçon, je sais légender une carte des sites touristiques du Cameroun
\end{itemize}\end{multicols}}

\begin{sommaire}
\begin{multicols}{2}
\begin{enumerate}
\item Définitions et cadre général du tourisme au Cameroun
\item Les potentialités touristiques du Cameroun
\item L'importance du tourisme au Cameroun
\item Les problèmes du secteur touristique camerounais
\item Perspectives et valorisation du tourisme pour l'intégration nationale
\item Activité d'intégration
\item Méthodes et exercices résolus
\item Exercices
\item Corrigés des exercices
\item Fiche bilan
\end{enumerate}
\end{multicols}
\end{sommaire}

\begin{objectifs}
\begin{itemize}
\item À la fin de cette leçon, je sais définir les notions de site touristique et de tourisme d'affaires
\item À la fin de cette leçon, je sais localiser sur une carte les principales zones et sites touristiques du Cameroun
\item À la fin de cette leçon, je sais décrire et classer les potentialités touristiques du Cameroun (balnéaire, safari, culturel, écotourisme)
\item À la fin de cette leçon, je sais expliquer l'importance économique, politique et sociale du tourisme au Cameroun
\item À la fin de cette leçon, je sais identifier et hiérarchiser les problèmes qui freinent le secteur touristique camerounais
\item À la fin de cette leçon, je sais légender une carte des sites touristiques du Cameroun
\item À la fin de cette leçon, je sais construire un croquis ou un diagramme à partir de données touristiques chiffrées
\item À la fin de cette leçon, je sais inventorier, classer et interpréter des documents (cartes, photographies, tableaux statistiques) relatifs au tourisme
\end{itemize}
\end{objectifs}

\begin{prerequis}
\begin{itemize}
\item Les grandes régions naturelles et humaines du Cameroun (leçons de Première et de début d'année)
\item La localisation des parcs nationaux et des grandes villes du Cameroun
\item La notion de développement et les indicateurs économiques (PIB, devises, emploi)
\item Les techniques de base de lecture et de légendage d'une carte
\item La notion d'intégration nationale (famille de situations du module)
\end{itemize}
\end{prerequis}

\newpage

\coursec{Définitions et cadre général du tourisme au Cameroun}

Chaque année, des milliers de visiteurs étrangers découvrent les plages de sable fin de Kribi, observent les éléphants du parc national de Waza ou s'émerveillent devant l'architecture des chefferies de l'Ouest, tandis que des hommes d'affaires se pressent dans les grands hôtels de Yaoundé et de Douala. Pourtant, le Cameroun, surnommé « l'Afrique en miniature », accueille moins d'un million de touristes internationaux par an (ordre de grandeur : quelques centaines de milliers de visiteurs, selon les années), loin derrière le Maroc ou l'Afrique du Sud qui en comptent plusieurs millions. Comment expliquer ce paradoxe entre des potentialités touristiques énormes et une fréquentation modeste ? Quelle place le tourisme occupe-t-il réellement dans l'économie et l'intégration nationale du Cameroun ? C'est à ces questions que cette leçon se propose de répondre. Mais avant de mesurer les richesses et les faiblesses du secteur, il faut d'abord s'entendre sur les mots : qu'est-ce que le tourisme ? Qu'est-ce qu'un site touristique ? Qui organise cette activité au Cameroun ?

\courssub{Qu'est-ce que le tourisme ?}

\textbf{Une intuition pour commencer.} Imagine un habitant de Bafoussam qui, pendant les vacances, part passer une semaine à Limbé pour profiter de la plage et visiter le jardin botanique. Imagine aussi une entrepreneuse nigériane qui vient passer trois jours à Douala pour un salon professionnel, logeant à l'hôtel et dînant au restaurant. Dans les deux cas, ces personnes ont quitté temporairement leur lieu de résidence habituelle, ont dépensé de l'argent sur place (transport, hébergement, restauration, visites) et retourneront chez elles. Toutes deux font du \cle{tourisme}, même si leurs motivations sont très différentes.

\begin{definition}[Le tourisme]
Le \cle{tourisme} est l'ensemble des déplacements temporaires de personnes hors de leur lieu de résidence habituelle, à des fins de loisir, de découverte ou d'affaires, ainsi que l'ensemble des activités et des équipements (transports, hôtels, restaurants, agences, sites aménagés) liés à ces déplacements.
\end{definition}

Trois conditions caractérisent donc le tourisme :

\begin{itemize}
\item un \textbf{déplacement} : on quitte son lieu de résidence habituelle ;
\item un caractère \textbf{temporaire} : le touriste a l'intention de revenir (on exclut donc les migrations définitives) ;
\item une \textbf{motivation non lucrative sur place} : le touriste ne vient pas travailler durablement dans le lieu visité ; il vient se reposer, découvrir, se divertir ou accomplir une mission professionnelle ponctuelle.
\end{itemize}

\begin{attention}[Piège fréquent]
Ne confonds pas le tourisme avec la \textbf{migration}. Un étudiant qui quitte Ngaoundéré pour s'installer durablement à Yaoundé afin d'y faire ses études n'est pas un touriste : il migre. De même, les déplacements quotidiens domicile-travail (migrations alternantes) ne relèvent pas du tourisme. Le touriste se déplace \emph{temporairement} et \emph{hors de son cadre de vie habituel}.
\end{attention}

\courssub{Tourisme interne et tourisme international}

On distingue classiquement deux grandes formes de tourisme selon l'origine des visiteurs.

\begin{definition}[Tourisme interne et tourisme international]
Le \cle{tourisme interne} (ou domestique) désigne les déplacements touristiques des résidents d'un pays à l'intérieur de leur propre pays. Le \cle{tourisme international} désigne les déplacements de visiteurs étrangers entrant dans un pays (tourisme récepteur) ou de résidents partant à l'étranger (tourisme émetteur).
\end{definition}

Au Cameroun, le tourisme interne concerne par exemple les familles camerounaises qui descendent à Kribi pendant les fêtes de fin d'année, les élèves en excursion au monument de la Réunification à Yaoundé, ou les fidèles qui visitent des lieux culturels lors des festivals traditionnels (Ngondo à Douala, Nguon à Foumban). Ce tourisme est encore peu développé, faute de revenus suffisants et d'une culture des loisirs organisés, mais il progresse avec la croissance des classes moyennes urbaines.

Le tourisme international, lui, est surtout récepteur : le Cameroun reçoit des visiteurs venus principalement d'Europe (France en tête), d'Afrique et, de plus en plus, d'Asie. Une part importante de ces visiteurs est constituée de \textbf{ressortissants camerounais de la diaspora} qui reviennent au pays pour les vacances ou les cérémonies familiales : on parle parfois de « tourisme de retour ».

\begin{exemplebox}[Deux touristes, deux logiques]
\textbf{Cas 1.} Mme Etoundi, fonctionnaire à Ebolowa, emmène ses enfants voir les chutes de la Lobé pendant les vacances scolaires : c'est du \emph{tourisme interne de découverte}.\\
\textbf{Cas 2.} Un couple de retraités français effectue un circuit de douze jours organisé par une agence : Kribi, Foumban, les monts Mandara et le parc de Waza : c'est du \emph{tourisme international récepteur}.\\
Dans les deux cas, les dépenses réalisées (carburant, hôtels, restaurants, guides, artisanat) irriguent l'économie locale.
\end{exemplebox}

\courssub{Le site touristique : le cœur de l'attraction}

Tout touriste se déplace \emph{vers quelque chose}. Ce « quelque chose » est le site touristique.

\begin{definition}[Site touristique]
Un \cle{site touristique} est un lieu qui, par ses attraits naturels (paysage, plage, faune, relief, climat) ou culturels (monument, musée, chefferie, festival, savoir-faire artisanal), attire des visiteurs et fait l'objet d'une fréquentation organisée (accès, accueil, encadrement, mise en valeur).
\end{definition}

Deux éléments de cette définition méritent d'être soulignés. D'une part, le site possède un \textbf{attrait} : une caractéristique remarquable qui le distingue (la chute d'eau, l'éléphant, le palais du sultan). D'autre part, il fait l'objet d'une \textbf{fréquentation organisée} : un espace magnifique mais totalement inaccessible et inconnu n'est pas encore un site touristique ; il ne le devient que lorsqu'il est aménagé (pistes, auberges, panneaux, guides) et fréquenté. C'est la différence entre une \emph{potentialité} touristique et un \emph{site} touristique effectif.

\begin{methode}[Analyser une photographie de site touristique]
Face à une photographie de site touristique (exercice classique au Baccalauréat), procède en trois étapes :
\begin{enumerate}
\item \textbf{Décrire} ce que l'on voit, de l'avant-plan vers l'arrière-plan, sans interpréter : éléments naturels (mer, rochers, végétation), éléments humains (touristes, pirogues, constructions), équipements (hôtels, routes, panneaux).
\item \textbf{Identifier l'attrait} : qu'est-ce qui attire le visiteur sur ce site ? Le paysage ? La faune ? Un monument ? Une activité ?
\item \textbf{Dégager le type de tourisme} concerné (balnéaire, safari, culturel, écotourisme, affaires) et, si possible, formuler une hypothèse sur les retombées ou les problèmes du site.
\end{enumerate}
\end{methode}

\begin{exoresolu}[Commentaire d'une photographie : les chutes de la Lobé (Kribi)]
\textbf{Énoncé.} Une photographie montre, au premier plan, une plage de sable clair bordée de cocotiers ; au centre, une série de chutes d'eau larges et peu hautes qui se jettent directement dans l'océan Atlantique ; à l'arrière-plan, une forêt dense ; quelques pirogues de pêcheurs et de petits restaurants de fortune en bordure de plage ; des visiteurs se baignent au pied des chutes. Décris cette photographie, identifie l'attrait du site et précise le type de tourisme concerné.
\tcblower
\textbf{Solution détaillée.}
\begin{enumerate}
\item \textbf{Description.} L'avant-plan est occupé par une plage de sable clair plantée de cocotiers. Au centre de l'image, la rivière Lobé franchit un dernier ressaut rocheux et forme de larges chutes qui se déversent directement dans l'océan : c'est un phénomène rare, une rivière qui se jette en cascade dans la mer. L'arrière-plan est couvert d'une forêt dense. On observe enfin des pirogues, de petits restaurants sommaires et des baigneurs.
\item \textbf{Attrait.} L'attrait principal est \emph{paysager} : la rencontre exceptionnelle entre une chute d'eau et l'océan, dans un cadre de plage tropicale. L'attrait secondaire est la possibilité de se baigner et de déguster du poisson frais.
\item \textbf{Type de tourisme.} Il s'agit d'un \emph{tourisme balnéaire} (mer, plage, détente) combiné à un \emph{tourisme de découverte paysagère}, voire à un début d'\emph{écotourisme} (cadre naturel préservé). On remarque toutefois la simplicité des équipements d'accueil (restaurants de fortune), ce qui illustre un problème récurrent : l'insuffisance des structures d'accueil autour des sites camerounais.
\end{enumerate}
\end{exoresolu}

\courssub{Le tourisme d'affaires : une forme méconnue mais essentielle}

Quand on pense « touriste », on imagine un vacancier en short avec un appareil photo. Or, au Cameroun, une part considérable de la clientèle des grands hôtels n'est pas là pour se reposer.

\begin{definition}[Tourisme d'affaires]
Le \cle{tourisme d'affaires} désigne les déplacements temporaires motivés par des activités professionnelles : congrès, séminaires, salons, foires, missions, négociations commerciales. Il constitue une part importante de la clientèle hôtelière des grandes villes camerounaises, en particulier Yaoundé (capitale politique et administrative) et Douala (capitale économique et portuaire).
\end{definition}

Le tourisme d'affaires présente des caractéristiques propres : il est \textbf{régulier} (il ne dépend pas des saisons de vacances), il concerne une clientèle au \textbf{pouvoir d'achat élevé} (hôtels de standing, restaurants, location de salles et de véhicules), et il est \textbf{concentré dans les grandes villes}. À Yaoundé, les institutions nationales (ministères, ambassades, sièges d'organisations internationales) génèrent congrès et séminaires ; le Palais des Congrès accueille de grandes manifestations. À Douala, le port, les banques et les sièges d'entreprises attirent hommes d'affaires et expatriés. Les hôtels de ces deux villes (Hilton, Mont Fébé, La Falaise à Yaoundé ; Pullman, Sawa, Akwa Palace à Douala) doivent une large part de leur taux d'occupation à cette clientèle.

\begin{attention}[Ne pas sous-estimer le tourisme d'affaires]
Dans les copies, beaucoup d'élèves réduisent le tourisme au seul tourisme de loisirs. Or, au Cameroun, le tourisme d'affaires est souvent \emph{plus rentable et plus stable} que le tourisme de vacances. L'oublier, c'est passer à côté d'une réalité majeure du secteur.
\end{attention}

\courssub{La typologie du tourisme au Cameroun}

À partir des attraits du pays, on peut classer le tourisme camerounais en cinq grands types, que le schéma suivant synthétise et qui seront étudiés en détail dans la section consacrée aux potentialités.

\begin{popfigure}
\begin{center}
\begin{tikzpicture}[
  font=\small,
  root/.style={rectangle, rounded corners, draw=popdark, fill=popblue, text=white, align=center, minimum width=4.6cm, minimum height=1cm},
  br/.style={rectangle, rounded corners, draw=popdark, fill=popblueL, align=center, minimum width=4.4cm, minimum height=1.3cm},
  ex/.style={rectangle, rounded corners, draw=popline, fill=popcream, align=center, minimum width=4.4cm, minimum height=1.1cm, font=\footnotesize}
]
\node[root] (T) {Tourisme au Cameroun};
\node[br, below left=1.6cm and 3.4cm of T, align=center] (B){\textbf{Tourisme balnéaire}\\ (mer et plages)};
\node[br, below left=1.6cm and -3.4cm of T, align=center] (S){\textbf{Tourisme safari}\\ (faune et parcs)};
\node[br, below right=1.6cm and -3.4cm of T, align=center] (C){\textbf{Tourisme culturel}\\ (traditions, monuments)};
\node[br, below right=1.6cm and 3.4cm of T, align=center] (E){\textbf{Écotourisme}\\ (forêt, nature)};
\node[br, below=4.6cm of T, align=center] (A){\textbf{Tourisme d'affaires}\\ (congrès, séminaires, missions)};
\node[ex, below=0.25cm of B, align=center] (Be){Kribi, Limbé,\\ chutes de la Lobé};
\node[ex, below=0.25cm of S, align=center] (Se){Parcs de Waza, Bénoué,\\ Bouba Ndjida (Septentrion)};
\node[ex, below=0.25cm of C, align=center] (Ce){Chefferies de l'Ouest,\\ Foumban, festivals (Ngondo, Nguon)};
\node[ex, below=0.25cm of E, align=center] (Ee){Forêts du Sud et de l'Est,\\ réserve du Dja, mont Cameroun};
\node[ex, below=0.25cm of A, align=center] (Ae){Yaoundé et Douala\\ (hôtels, Palais des Congrès)};
\draw[-{Stealth}, thick, popdark] (T) -- (B);
\draw[-{Stealth}, thick, popdark] (T) -- (S);
\draw[-{Stealth}, thick, popdark] (T) -- (C);
\draw[-{Stealth}, thick, popdark] (T) -- (E);
\draw[-{Stealth}, thick, popdark] (T) -- (A);
\draw[-, popline] (B) -- (Be);
\draw[-, popline] (S) -- (Se);
\draw[-, popline] (C) -- (Ce);
\draw[-, popline] (E) -- (Ee);
\draw[-, popline] (A) -- (Ae);
\end{tikzpicture}
\end{center}
\legende{Schéma 1 : typologie du tourisme au Cameroun. 1 : tourisme balnéaire (littoral atlantique) ; 2 : tourisme safari (parcs nationaux du septentrion) ; 3 : tourisme culturel (hauts plateaux de l'Ouest et traditions) ; 4 : écotourisme (plateau forestier du Sud) ; 5 : tourisme d'affaires (Yaoundé et Douala).}
\end{popfigure}

Ce schéma montre déjà l'essentiel : le Cameroun possède une gamme \emph{complète} de produits touristiques, ce qui justifie son surnom.

\begin{savaistu}[Pourquoi dit-on « l'Afrique en miniature » ?]
Le Cameroun est surnommé \cle{l'Afrique en miniature} parce qu'il réunit sur son territoire presque tous les paysages et tous les climats du continent africain : la forêt équatoriale humide du Sud, les savanes herbeuses de l'Adamaoua, les savanes sèches et les steppes du Grand Nord, les hautes terres fraîches de l'Ouest, un volcan actif (le mont Cameroun, environ \num{4100} m d'altitude, point culminant de l'Afrique de l'Ouest et centrale), des plages atlantiques et une mosaïque de plus de 250 groupes ethniques aux cultures variées. En un seul pays, le visiteur peut découvrir ce que l'Afrique offre de plus divers.
\end{savaistu}

\courssub{Le cadre institutionnel du tourisme au Cameroun}

Le tourisme n'est pas une activité spontanée : il est organisé, encadré et promu par des institutions et des professionnels.

\textbf{Le ministère du Tourisme et des Loisirs (MINTOUL).} C'est l'organe de l'État chargé de définir et de mettre en œuvre la politique touristique nationale : promotion du pays à l'étranger (campagnes, salons internationaux, slogan « Cameroon, All of Africa in one country »), réglementation des professions touristiques, classement des hôtels, aménagement des sites, formation. Le tourisme a longtemps été rattaché à d'autres ministères avant de disposer d'un département ministériel à part entière, signe d'une prise de conscience progressive de son importance.

\textbf{Les agences de voyage et les tours-opérateurs.} Elles conçoivent et vendent les circuits (transport, hébergement, visites), principalement à Douala et Yaoundé. Ce sont elles qui transforment une potentialité en produit touristique consommable.

\textbf{Les guides touristiques.} Professionnels de l'accueil et de l'interprétation, ils accompagnent les visiteurs dans les parcs (guides et pisteurs à Waza), sur les sites culturels ou en montagne. Leur formation, encore insuffisante, est un enjeu majeur que nous retrouverons parmi les problèmes du secteur.

\textbf{Les autres acteurs.} S'ajoutent les hôteliers et restaurateurs, les transporteurs (compagnies aériennes, dont Camair-Co, la compagnie nationale ; agences de location de véhicules ; cars interurbains), les artisans (vendeurs d'objets d'art), les collectivités locales et les communautés villageoises qui gèrent certains sites (écotourisme communautaire autour de la réserve du Dja, par exemple).

\begin{aretenir}[L'essentiel de la section]
\begin{itemize}
\item Le \cle{tourisme} : déplacement temporaire hors du lieu de résidence habituelle, à des fins de loisir, de découverte ou d'affaires.
\item On distingue le \cle{tourisme interne} (résidents visitant leur pays) du \cle{tourisme international} (visiteurs étrangers).
\item Un \cle{site touristique} est un lieu dont les attraits naturels ou culturels attirent des visiteurs et qui fait l'objet d'une fréquentation organisée.
\item Le \cle{tourisme d'affaires} (congrès, séminaires, missions) est concentré à Yaoundé et Douala et nourrit une large part de l'hôtellerie urbaine.
\item Le tourisme camerounais se décline en cinq types : balnéaire, safari, culturel, écotourisme et affaires.
\item Le secteur est encadré par le ministère du Tourisme et des Loisirs, les agences de voyage, les guides, les hôteliers et les transporteurs.
\end{itemize}
\end{aretenir}

\begin{exercice}[Pour vérifier tes acquis]
\begin{enumerate}
\item Définis les termes suivants : tourisme ; site touristique ; tourisme d'affaires.
\item Un commerçant de Maroua se rend chaque semaine à N'Djamena (Tchad) pour y acheter des marchandises qu'il revend au Cameroun. S'agit-il de tourisme ? Justifie ta réponse.
\item Pourquoi le tourisme d'affaires est-il particulièrement développé à Yaoundé et à Douala ?
\item À partir du schéma 1, cite pour chaque type de tourisme un exemple de site ou de lieu camerounais.
\end{enumerate}
\end{exercice}

\begin{corrige}[Exercice « Pour vérifier tes acquis »]
\begin{enumerate}
\item Voir les boîtes de définition : le tourisme est un déplacement temporaire hors du lieu de résidence habituelle à des fins de loisir, de découverte ou d'affaires ; un site touristique est un lieu dont les attraits naturels ou culturels attirent des visiteurs et font l'objet d'une fréquentation organisée ; le tourisme d'affaires regroupe les déplacements motivés par des activités professionnelles (congrès, séminaires, missions).
\item Non. Le déplacement est bien temporaire, mais sa motivation est \emph{lucrative et régulière} : il s'agit d'une activité commerciale habituelle (commerce transfrontalier), non d'un séjour de loisir, de découverte ou d'une mission ponctuelle. Ce n'est donc pas du tourisme.
\item Yaoundé est la capitale politique et administrative (ministères, ambassades, organisations internationales, Palais des Congrès) ; Douala est la capitale économique (port, banques, sièges d'entreprises). Ces fonctions génèrent en permanence congrès, séminaires et missions, qui remplissent les hôtels des deux villes.
\item Balnéaire : plages de Kribi ou de Limbé ; safari : parc national de Waza ; culturel : chefferie de Foumban (palais des sultans Bamoun) ; écotourisme : réserve de faune du Dja ; affaires : hôtels et Palais des Congrès de Yaoundé.
\end{enumerate}
\end{corrige}

Maintenant que le vocabulaire et le cadre institutionnel sont posés, nous pouvons aborder la section suivante : l'inventaire détaillé des \textbf{potentialités touristiques} du Cameroun, de l'attrait de la mer à l'écotourisme forestier, en passant par les safaris du septentrion et le dynamisme culturel des hauts plateaux.

\coursec{Les potentialités touristiques du Cameroun}

Dans la section précédente, nous avons défini le \cle{tourisme} et situé le secteur touristique camerounais dans son cadre général : un pays que l'on surnomme « l'Afrique en miniature ». Ce surnom n'est pas une formule publicitaire : il traduit une réalité géographique fondamentale. En moins de \SI{475 000}{km^{2}}, le Cameroun concentre la mer, la savane, la forêt équatoriale, les hautes terres volcaniques et une mosaïque de cultures. Cette section répond à une question simple : \emph{que possède réellement le Cameroun pour attirer les visiteurs ?} Autrement dit, quelles sont ses \cle{potentialités touristiques} ?

\begin{definition}[Potentialité touristique]
Une \cle{potentialité touristique} est l'ensemble des ressources naturelles (relief, climat, faune, flore, littoral) et culturelles (monuments, traditions, artisanat, gastronomie) d'un territoire susceptibles d'attirer des visiteurs et d'être aménagées en \cle{sites touristiques}. On distingue les potentialités \emph{naturelles} (mer, parcs, montagnes, forêts) et les potentialités \emph{culturelles} (chefferies, danses, musées, artisanat).
\end{definition}

L'idée directrice de cette section est la suivante : \textbf{à chaque grande région naturelle du Cameroun correspond un type de tourisme dominant}. Nous allons donc parcourir le pays du littoral vers le septentrion, puis des hauts plateaux de l'Ouest vers le plateau forestier du Sud, en identifiant pour chaque espace ses atouts propres.

\courssub{L'attrait de la mer : le tourisme balnéaire sur la côte atlantique}

Imagine un touriste européen qui débarque à l'aéroport de Douala en décembre, fuyant l'hiver. À quelques heures de route, deux destinations s'offrent à lui : \textbf{Kribi} au sud et \textbf{Limbé} à l'ouest. C'est le premier grand atout du Cameroun : son ouverture sur l'océan Atlantique, avec environ \SI{400}{km} de côtes sur le golfe de Guinée.

\begin{itemize}
\item \textbf{Les plages de Kribi} (région du Sud) : sable blond et fin, mer relativement calme, cocotiers en bordure d'une forêt dense. Kribi est la station balnéaire la plus réputée du pays. Son développement est renforcé par la proximité du \textbf{Port Autonome de Kribi (PAK)}, port en eau profonde qui draine aussi un \cle{tourisme d'affaires}.
\item \textbf{Les plages de Limbé} (région du Sud-Ouest) : sable noir d'origine volcanique, héritage du mont Cameroun tout proche. Limbé possède en outre un jardin botanique créé à l'époque allemande (1892) et un centre de protection de la faune (\emph{Limbe Wildlife Centre}), ce qui combine tourisme balnéaire et découverte de la nature.
\item \textbf{Les chutes de la Lobé}, près de Kribi : un site exceptionnel où la rivière Lobé franchit une dernière cascade avant de se jeter directement dans l'océan Atlantique, au milieu d'une forêt de cocotiers.
\end{itemize}

\begin{savaistu}[Un phénomène rare au monde]
Les chutes de la Lobé, à quelques kilomètres de Kribi, comptent parmi les très rares endroits au monde où une cascade se jette directement dans l'océan. On peut y observer les pêcheurs qui franchissent les rapides en pirogue. Ce site, facilement accessible, est devenu l'emblème du tourisme balnéaire camerounais.
\end{savaistu}

\courssub{Le safari dans les parcs nationaux du septentrion}

Le nord du Cameroun (régions du Nord, de l'Adamaoua et de l'Extrême-Nord) est un domaine de savane soudanienne et sahélienne : herbes hautes, arbres épars, saisons bien marquées. C'est le décor classique du \cle{safari} photographique, comparable aux grands parcs d'Afrique de l'Est. Trois parcs nationaux dominent :

\begin{itemize}
\item \textbf{Le parc national de Waza} (Extrême-Nord, créé en 1934 comme réserve, érigé en parc national en 1968) : environ \SI{1 700}{km^{2}} de plaines inondables et de savanes. C'est \emph{historiquement} le parc le plus visité du Cameroun. On y observe éléphants, lions, girafes, antilopes (damalisques, cobes, gazelles), hyènes, ainsi qu'une avifaune très riche (autruches, aigles, marabouts).
\item \textbf{Le parc national de la Bénoué} (région du Nord) : environ \SI{1 800}{km^{2}} le long de la rivière Bénoué ; hippopotames, buffles, lions, antilopes. Il est classé réserve de biosphère par l'UNESCO (1981).
\item \textbf{Le parc national de Bouba Ndjida} (région du Nord) : environ \SI{2 200}{km^{2}} ; il est notamment connu pour abriter l'éland de Derby, la plus grande antilope du monde.
\end{itemize}

\begin{exemplebox}[Le parc national de Waza, vitrine du safari camerounais]
Situé à environ \SI{120}{km} au nord de Maroua, le parc de Waza est le parc le plus connu du Cameroun et abrite l'une des faunes les plus riches d'Afrique centrale : plusieurs centaines d'éléphants, des lions, des girafes du Kordofan (sous-espèce menacée) et plus de 350 espèces d'oiseaux (ordre de grandeur). La meilleure période d'observation est la saison sèche (novembre à mai), lorsque les animaux se concentrent autour des points d'eau. Le parc a cependant souffert de l'insécurité liée à Boko Haram dans les années 2010, ce qui illustre la fragilité du secteur (voir section 4).
\end{exemplebox}

\courssub{Le dynamisme culturel des hauts plateaux de l'Ouest}

Les régions de l'Ouest et du Nord-Ouest, domaine des hauts plateaux (1 000 à \SI{2 000}{m} d'altitude, climat tempéré), constituent le cœur du \cle{tourisme culturel} camerounais. Ici, l'attraction principale n'est ni la mer ni la faune, mais des civilisations vivantes :

\begin{itemize}
\item \textbf{Les chefferies et palais royaux} : la chefferie de \textbf{Foumban}, siège du sultanat Bamoun, avec son palais, son musée et l'écriture \emph{shümom} inventée par le sultan Njoya ; la chefferie de \textbf{Bafut} (Nord-Ouest) et son palais classé ; les chefferies bamiléké autour de \textbf{Bafoussam} et Bandjoun, remarquables par leur architecture (grandes cases au toit conique).
\item \textbf{Les danses traditionnelles} : danses guerrières et cérémonielles (Toupouri et Kirdi au Nord, danses bamiléké et bamoun à l'Ouest, danses des sociétés secrètes au Nord-Ouest), souvent liées aux funérailles et aux fêtes annuelles, comme la fête du Nguon à Foumban (tous les deux ans).
\item \textbf{L'artisanat} : sculptures sur bois, masques, bronzes de Foumban, poteries, tissus (le \emph{ndop}), tabourets perlés. Les marchés artisanaux de Foumban et de Bamenda attirent collectionneurs et touristes.
\end{itemize}

Ce dynamisme culturel est un atout majeur : il permet un tourisme toute l'année (il ne dépend pas des saisons comme le safari) et il valorise l'identité nationale, donc l'\cle{intégration nationale}.

\courssub{L'écotourisme dans le plateau forestier du Sud}

Le sud du pays (régions du Sud et de l'Est, sud du Centre) est couvert par la forêt équatoriale dense et humide, l'une des plus anciennes et des plus riches du monde. C'est le domaine de l'\cle{écotourisme}, c'est-à-dire d'un tourisme fondé sur l'observation de la nature dans le respect des écosystèmes et des populations locales.

\begin{itemize}
\item \textbf{Le parc national de Campo-Ma'an} (région du Sud, environ \SI{2 600}{km^{2}}) : gorilles des plaines de l'Ouest, chimpanzés, éléphants de forêt, mandrills, et une diversité végétale exceptionnelle.
\item \textbf{Le parc national de Korup} (région du Sud-Ouest, environ \SI{1 260}{km^{2}}) : l'une des forêts tropicales les plus anciennes d'Afrique (refuge pléistocène) ; on y trouve des primates rares (drills, talapoïns) et des centaines d'espèces d'oiseaux et de papillons.
\item \textbf{L'observation des gorilles} : le Cameroun fait partie des rares pays où l'on peut observer les gorilles de plaine en milieu naturel (Campo-Ma'an, mais aussi le sanctuaire de la Mefou près de Yaoundé, géré par l'ONG Ape Action Africa).
\end{itemize}

\begin{attention}[Parc national ou réserve de chasse ?]
Erreur fréquente au Bac : confondre \cle{parc national} et \cle{réserve de chasse}. Dans un parc national (Waza, Bénoué, Korup, Campo-Ma'an), la protection est stricte : la chasse y est interdite, seul le safari photographique est permis. Les réserves de chasse (ou zones d'intérêt cynégétique), nombreuses dans l'Est et le Nord, sont au contraire des espaces où la chasse sportive est autorisée et encadrée, moyennant des permis payants : c'est le \cle{tourisme de chasse}, pratiqué surtout dans l'Est (région forestière giboyeuse : bongos, sitatungas, buffles de forêt). Ne jamais écrire que « les touristes chassent dans le parc de Waza ».
\end{attention}

\courssub{Compléments utiles au Bac : montagne et chasse}

Deux formes de tourisme complètent ce tableau :

\begin{itemize}
\item \textbf{Le tourisme de montagne} : le mont Cameroun (\SI{4 070}{m}), point culminant de l'Afrique centrale et de l'Afrique de l'Ouest, est un volcan actif dont l'ascension attire randonneurs et sportifs. La \emph{course de l'espoir} (\emph{Mount Cameroon Race of Hope}), organisée chaque année à Buéa depuis 1995 (édition internationale), est une course internationale d'ascension qui draine athlètes et spectateurs : bel exemple de tourisme sportif.
\item \textbf{Le tourisme de chasse} : dans l'Est et le Nord, des zones cynégétiques aménagées accueillent des chasseurs sportifs étrangers, source de devises mais aussi de débats sur la conservation.
\end{itemize}

\courssub{Répartition spatiale : un type de tourisme par région naturelle}

L'observation des documents (carte des parcs, photographies des sites) permet de dégager une loi de répartition : \textbf{chaque grande région naturelle du Cameroun possède un type de tourisme dominant}. C'est cette complémentarité qui fonde le slogan « l'Afrique en miniature ».

\begin{center}
\begin{tabularx}{\linewidth}{|l|X|X|}
\hline
\rowcolor{popblueL} \textbf{Région naturelle} & \textbf{Type de tourisme dominant} & \textbf{Sites et attraits principaux} \\
\hline
Littoral atlantique (Sud-Ouest, Littoral, Sud) & Tourisme balnéaire et d'affaires & Plages de Kribi et de Limbé, chutes de la Lobé, Douala, Port de Kribi \\
\hline
Savanes du septentrion (Nord, Adamaoua, Extrême-Nord) & Safari (tourisme animalier) & Parcs de Waza, de la Bénoué, de Bouba Ndjida ; éléphants, lions, girafes, antilopes \\
\hline
Hauts plateaux de l'Ouest (Ouest, Nord-Ouest) & Tourisme culturel & Chefferies de Foumban, Bafut, Bafoussam ; danses, artisanat, palais des sultans \\
\hline
Plateau forestier du Sud (Sud, Est, Sud-Ouest) & Écotourisme & Parcs de Campo-Ma'an et de Korup, gorilles, forêt équatoriale \\
\hline
Montagne occidentale (Sud-Ouest) & Tourisme de montagne et sportif & Ascension du mont Cameroun, course de l'espoir \\
\hline
\end{tabularx}
\end{center}

\begin{popfigure}
\begin{center}
\includegraphics[width=\linewidth,height=11cm,keepaspectratio]{N01/cmr_parcs.jpg}
\end{center}
\legende{Les grandes zones touristiques sur fond de relief : au sud-ouest, le littoral balnéaire (Kribi, Limbé) et le mont Cameroun ; au nord, les savanes et les parcs du septentrion (Waza, Bénoué, Bouba Ndjida), lieux du tourisme de safari ; à l'ouest, les hauts plateaux (tourisme culturel) ; au sud et à l'est, la forêt (Korup, Dja, Campo-Ma'an, Boumba Bek, Nki), terrain d'écotourisme. La carte localise les parcs nationaux du pays ; titre en espagnol, noms de parcs lisibles. \\ Source : Wikimedia Commons, Teogomez, CC0 (légende des parcs en espagnol pour le titre ; noms des parcs lisibles).}
\end{popfigure}

\courssub{TP guidé : localiser et légender les sites touristiques sur une carte muette}

\begin{methode}[Légender une carte des sites touristiques]
Pour réussir le TP de localisation (savoir-faire exigible au Bac), procède par étapes :
\begin{enumerate}
\item \textbf{Trace le fond de carte} : contour du Cameroun, limites des grandes régions naturelles, flèche du nord, titre et échelle indicative.
\item \textbf{Choisis des symboles distincts par type de tourisme} : par exemple un point bleu pour les plages, un triangle vert pour les parcs animaliers, un carré violet pour les sites culturels, un cercle vert foncé pour les forêts d'écotourisme.
\item \textbf{Place les sites avec soin} : Kribi et Limbé sur la côte ; Waza au nord de Maroua, Bénoué près de Garoua, Bouba Ndjida à l'est ; Foumban, Bafoussam, Bafut sur les hauts plateaux ; Campo-Ma'an et Korup dans le Sud forestier ; le mont Cameroun près de Buéa.
\item \textbf{Utilise une échelle de couleurs cohérente} : une couleur par type de tourisme, jamais deux couleurs pour le même type.
\item \textbf{Rédige la légende} dans l'ordre logique (du littoral vers l'intérieur), avec des traits fins reliant chaque symbole à sa signification.
\end{enumerate}
\end{methode}

\begin{exoresolu}[Associer un site à son type de tourisme]
Énoncé : classe les sites suivants selon le type de tourisme dominant (balnéaire, safari, culturel, écotourisme, montagne) : Limbé, Waza, Foumban, Korup, mont Cameroun, Kribi, Bafut, Campo-Ma'an, Bouba Ndjida.
\tcblower
Solution détaillée :
\begin{itemize}
\item \textbf{Tourisme balnéaire} : Limbé (plages de sable noir volcanique) et Kribi (plages de sable fin, chutes de la Lobé).
\item \textbf{Safari (tourisme animalier)} : Waza (éléphants, lions, girafes) et Bouba Ndjida (éland de Derby), tous deux dans les savanes du septentrion.
\item \textbf{Tourisme culturel} : Foumban (palais du sultan Bamoun, musée, bronzes) et Bafut (chefferie du Nord-Ouest).
\item \textbf{Écotourisme} : Korup (forêt ancienne, primates rares) et Campo-Ma'an (gorilles, éléphants de forêt).
\item \textbf{Tourisme de montagne} : mont Cameroun (ascension, course de l'espoir).
\end{itemize}
Vérification : chaque site est associé à la région naturelle qui le porte, ce qui confirme la loi de répartition spatiale dégagée dans cette leçon.
\end{exoresolu}

\begin{aretenir}[L'essentiel de la section]
Le Cameroun possède des potentialités touristiques \emph{énormes} et \emph{complémentaires} : la mer (Kribi, Limbé, chutes de la Lobé), le safari dans les parcs du septentrion (Waza, Bénoué, Bouba Ndjida), le dynamisme culturel des hauts plateaux de l'Ouest (Foumban, Bafoussam, Bafut), l'écotourisme dans la forêt du Sud (Campo-Ma'an, Korup, gorilles), auxquels s'ajoutent le tourisme de montagne (mont Cameroun, 4 070 m) et le tourisme de chasse (Est). À chaque région naturelle correspond un type de tourisme dominant : c'est ce qui fait du Cameroun « l'Afrique en miniature ». Mais une potentialité n'est pas encore une richesse : sa valorisation suppose des infrastructures et une sécurité, comme nous le verrons dans les sections suivantes.
\end{aretenir}

\coursec{L'importance du tourisme au Cameroun}

Dans la section précédente, nous avons dressé l'inventaire des \cle{potentialités touristiques} du Cameroun : plages de Kribi et de Limbé, safaris dans les parcs du septentrion (Waza, Bénoué, Bouba Ndjida), dynamisme culturel des hauts plateaux de l'Ouest (chefferies, danses, palais des sultans de Foumban), écotourisme dans la forêt du Sud et de l'Est. Mais à quoi servent ces richesses si elles ne profitent ni à l'État, ni aux entreprises, ni aux populations ? C'est toute la question de cette section : montrer que le tourisme n'est pas un luxe réservé aux étrangers fortunés, mais un \cle{secteur stratégique} qui rapporte des devises, crée des emplois, fait rayonner le pays et tisse des liens entre les régions. Pour le comprendre, partons d'une situation concrète.

\begin{exemplebox}[Une nuit à Kribi : qui gagne quoi ?]
Un couple de touristes européens passe une semaine à Kribi. Il paie son billet d'avion, une chambre d'hôtel en bord de mer, des repas de poisson braisé, une excursion aux chutes de la Lobé en pirogue, des sculptures achetées à un artisan, et les services d'un guide local. Suivons l'argent : l'hôtelier encaisse et paie ses salariés (réceptionnistes, cuisiniers, femmes de chambre) ; le restaurateur achète le poisson au pêcheur du village voisin ; le piroguier, le guide et l'artisan reçoivent directement un revenu ; l'État prélève des taxes et la TVA ; la compagnie aérienne et l'agence de voyage touchent leur part. Une seule dépense touristique déclenche donc une \cle{cascade de revenus} : c'est ce que les économistes appellent l'\emph{effet multiplicateur} du tourisme.
\end{exemplebox}

\courssub{L'importance économique du tourisme}

\textbf{Une source de devises étrangères.}\par
Pour un pays comme le Cameroun, qui importe des médicaments, des machines, des véhicules et une partie de sa nourriture, il faut disposer de \cle{devises} (euros, dollars) pour payer ces importations. Or le touriste international arrive avec des devises qu'il dépense sur place : chaque nuit d'hôtel payée par un visiteur étranger est, économiquement, l'équivalent d'une exportation — on « exporte » un service sans rien expédier hors du territoire. C'est pourquoi on parle d'\emph{exportation invisible}. Le tourisme figure ainsi parmi les premiers pourvoyeurs de devises du Cameroun \emph{hors pétrole}, aux côtés du cacao, du café, du coton et du bois.

\textbf{Recettes hôtelières et contribution au PIB.}\par
Les recettes du secteur proviennent de l'hôtellerie (le Cameroun compte plusieurs centaines d'hôtels classés, concentrés à Yaoundé, Douala, Kribi, Limbé, Garoua et Maroua), de la restauration, du transport et des agences de voyage. La contribution directe et indirecte du tourisme au \cle{PIB} camerounais reste modeste : de l'ordre de \textbf{quelques pourcents} (les estimations varient selon les sources et les années, généralement entre 2 \% et 5 \% si l'on inclut les effets indirects — ordre de grandeur à retenir avec prudence). Ce poids limité ne doit pas masquer l'essentiel : c'est un secteur \emph{en croissance} et surtout un secteur \emph{diversifiant}.

\begin{exemplebox}[Un poids modeste mais des devises précieuses]
Le tourisme ne représente que quelques pourcents du PIB camerounais, loin derrière l'agriculture, le commerce ou l'industrie pétrolière. Pourtant, il est l'un des premiers pourvoyeurs de \cle{devises} hors hydrocarbures. Retenons la nuance : un secteur peut être « petit » en part de PIB et « grand » en utilité stratégique, car les devises qu'il rapporte servent à financer les importations vitales du pays.
\end{exemplebox}

\textbf{La création d'emplois directs et indirects.}\par
C'est sans doute la contribution la plus visible du tourisme. Distinguons soigneusement deux catégories :

\begin{itemize}
\item Les \cle{emplois directs} : ceux qui existent \emph{parce que} le touriste est là — réceptionnistes et femmes de chambre des hôtels, serveurs et cuisiniers des restaurants, guides touristiques et guides de parc (écogardes assurant les safaris à Waza), chauffeurs de taxi et transporteurs, personnel des agences de voyage et des compagnies aériennes (Camair-Co).
\item Les \cle{emplois indirects} : ceux qui sont stimulés par la demande touristique sans contact direct avec le visiteur — l'agriculteur qui fournit légumes et fruits aux hôtels de Kribi, le pêcheur de Limbé, l'éleveur de l'Adamaoua qui approvisionne les restaurants de Ngaoundéré, l'artisan sculpteur de Foumban ou potier de Bamessing (près de Ndop), le commerçant qui vend tissus et souvenirs, le maçon qui construit les hôtels.
\end{itemize}

Le tourisme est une activité à forte intensité de main-d'œuvre : contrairement à l'extraction pétrolière, qui rapporte beaucoup mais emploie peu, chaque tranche de recettes touristiques soutient de nombreux emplois, souvent accessibles aux jeunes et aux femmes, et répartis sur tout le territoire — y compris dans des régions rurales où les autres employeurs sont rares.

\textbf{Un instrument de diversification de l'économie.}\par
L'économie camerounaise a longtemps reposé sur le pétrole (dont la production décline depuis la fin des années 1980) et sur quelques produits de base (cacao, café, coton, bois, aluminium) dont les prix fluctuent brutalement sur le marché mondial. Dépendre de peu de produits, c'est être vulnérable. Le tourisme offre une voie de \cle{diversification} : il mobilise des ressources \emph{renouvelables} (les paysages, la faune, les cultures ne s'épuisent pas si on les gère bien), il fait appel à des compétences locales et il injecte des revenus dans des régions éloignées des zones industrielles. C'est pourquoi les documents de stratégie nationale (comme la Stratégie de croissance et de l'emploi, puis la SND30) comptent le tourisme parmi les secteurs porteurs.

\courssub{L'importance politique du tourisme}

\textbf{Le rayonnement international et l'image de marque.}\par
Un pays visité est un pays connu, respecté, attractif. Le Cameroun a construit son image de marque sur un slogan célèbre : \cle{« l'Afrique en miniature »} — l'idée que le pays condense, sur son territoire, toutes les ambiances du continent (mer, savane, forêt, montagne, désert sahélien, diversité de plus de 250 groupes ethniques). Ce slogan, repris par le ministère du Tourisme et des Loisirs (MINTOUL), sert dans les salons internationaux du tourisme pour attirer visiteurs et investisseurs. Le tourisme devient ainsi un outil de \emph{diplomatie douce} (\emph{soft power}) : il façonne la perception du pays à l'étranger, bien plus efficacement qu'un discours officiel.

\textbf{La diplomatie par les grands événements.}\par
Les compétitions sportives et les conférences internationales fonctionnent comme des vitrines. L'organisation de la \cle{Coupe d'Afrique des Nations} (CAN 2021, disputée en janvier-février 2022) a placé le Cameroun sous les projecteurs du continent entier : des centaines de journalistes, des délégations, des supporters ont découvert le pays, et des téléspectateurs du monde entier ont vu ses stades et ses villes. De même, les sommets et conférences internationales accueillis à Yaoundé (au Palais des Congrès) ou à Douala nourrissent le \cle{tourisme d'affaires} : déplacements motivés par le travail (réunions, congrès, missions), qui remplissent les grands hôtels en semaine, hors des saisons de vacances.

\begin{savaistu}[La CAN 2021, un accélérateur touristique]
La CAN 2021 (jouée en 2022 à cause de la pandémie de Covid-19) a entraîné la construction du stade d'Olembe à Yaoundé et du stade de Japoma à Douala, la rénovation des stades Ahmadou Ahidjo (Yaoundé) et Roumdé Adjia (Garoua), ainsi que la construction ou la réhabilitation de nombreux hôtels et de voies d'accès. Ces infrastructures, financées pour le sport, servent durablement au \cle{tourisme événementiel} et d'affaires : congrès, concerts, compétitions. C'est un bel exemple d'investissement dont les retombées dépassent l'événement initial.
\end{savaistu}

\courssub{L'importance sociale du tourisme}

\textbf{La valorisation des cultures locales.}\par
Quand des visiteurs paient pour assister à une danse traditionnelle, visiter la chefferie de Bafut ou le palais des sultans de Foumban (haut lieu culturel du pays, dont le musée conserve les objets royaux bamoun), acheter un masque ou un tissu ndop, ils transforment la culture en \cle{ressource économique}. Cette valorisation a un effet social profond : elle encourage les jeunes générations à préserver et à transmettre des savoir-faire (sculpture, vannerie, musique, gastronomie) qui risquaient de disparaître. Le tourisme culturel des hauts plateaux de l'Ouest et du Nord-Ouest illustre parfaitement ce mécanisme.

\textbf{Des revenus pour les populations riveraines des parcs.}\par
Autour des parcs nationaux (Waza, Bénoué, Korup, Campo-Ma'an), le tourisme et l'écotourisme génèrent des emplois de guides, de pisteurs, de personnel de lodges et de campements. Une partie des droits d'entrée et des taxes de visite est reversée aux communautés locales, qui financent puits, écoles ou dispensaires. Ce partage est décisif : lorsque la faune « rapporte » aux villageois, ceux-ci deviennent les premiers protecteurs des animaux au lieu de les braconner. Le tourisme devient alors un outil de \cle{conservation} autant que de développement.

\textbf{La découverte du pays par les Camerounais eux-mêmes.}\par
On oublie souvent le \cle{tourisme interne} : les Camerounais qui visitent leur propre pays pendant les vacances, les voyages scolaires, les excursions d'entreprise, les pèlerinages et les retours au village. Ce tourisme-là a une vertu sociale majeure : un élève de Douala qui découvre les chutes d'Ekom Nkam, un fonctionnaire de Yaoundé qui visite le lamidat de Ngaoundéré, apprennent à connaître et à aimer la diversité de leur pays. On ne peut pas se sentir citoyen d'un territoire que l'on n'a jamais parcouru.

\courssub{Le tourisme, facteur d'intégration nationale}

C'est le point de convergence avec notre famille de situations. Les flux touristiques — nationaux et internationaux — \emph{relient} les régions : la route du tourisme va de Douala à Kribi, de Yaoundé à Foumban et Bafoussam, de Ngaoundéré à Waza et Maroua. Pour accueillir les visiteurs, l'État et les opérateurs sont incités à construire routes, aéroports et hôtels, infrastructures qui profitent ensuite à tous les habitants et désenclavent des régions entières. Surtout, le tourisme fabrique du \cle{sentiment national} : en circulant entre la forêt du Sud, les hautes terres de l'Ouest et les savanes du Nord, en consommant les produits et en découvrant les cultures de toutes les régions, les Camerounais éprouvent concrètement leur appartenance à \emph{une même nation}. Le slogan « Afrique en miniature » devient alors un projet politique : faire de la diversité une richesse partagée.

\begin{popfigure}
\begin{tikzpicture}[scale=1]
\node[draw=popdark, fill=popblue, text=white, rounded corners, font=\bfseries, minimum width=3.4cm, minimum height=1.1cm, align=center] (c) at (0,0) {Importance\\ du tourisme};
\node[draw=poporange, fill=poporangeL, rounded corners, align=center, minimum width=3.6cm, font=\small] (eco) at (-5.2,2.6) {\textbf{Économique}\\ devises, recettes hôtelières\\ emplois directs et indirects\\ diversification hors pétrole};
\node[draw=poppurple, fill=poppurpleL, rounded corners, align=center, minimum width=3.6cm, font=\small] (pol) at (5.2,2.6) {\textbf{Politique}\\ image « Afrique en miniature »\\ rayonnement international\\ CAN 2021, conférences};
\node[draw=popgreen, fill=popgreenL, rounded corners, align=center, minimum width=3.6cm, font=\small] (soc) at (0,-3.2) {\textbf{Sociale}\\ cultures valorisées, artisanat\\ revenus des riverains des parcs\\ intégration nationale};
\draw[-{Stealth}, very thick, poporange] (c) -- (eco);
\draw[-{Stealth}, very thick, poppurple] (c) -- (pol);
\draw[-{Stealth}, very thick, popgreen] (c) -- (soc);
\end{tikzpicture}
\legende{Schéma en étoile : les trois dimensions de l'importance du tourisme au Cameroun. 1 : dimension économique (devises, emplois, diversification) ; 2 : dimension politique (image de marque, diplomatie événementielle) ; 3 : dimension sociale (cultures, revenus locaux, intégration nationale).}
\end{popfigure}

\courssub{Analyse guidée d'un tableau statistique : arrivées et recettes touristiques}

Passons à la démarche du géographe : confronter les idées aux chiffres. Le tableau ci-dessous présente l'évolution des arrivées de touristes internationaux au Cameroun (ordres de grandeur arrondis, à retenir comme indicatifs).

\begin{center}
\begin{tabularx}{\linewidth}{|l|X|X|X|X|X|}
\hline
\rowcolor{popblueL} \textbf{Année} & \textbf{2015} & \textbf{2017} & \textbf{2019} & \textbf{2020} & \textbf{2022} \\
\hline
Arrivées internationales (milliers) & 800 & 950 & 1 000 & 350 & 700 \\
\hline
Recettes estimées (centaines de milliards de FCFA) & 300 & 350 & 400 & 150 & 300 \\
\hline
\end{tabularx}
\end{center}

\begin{popfigure}
\begin{tikzpicture}
\begin{axis}[
ybar, bar width=14pt,
width=\linewidth, height=6.5cm,
xlabel={Année}, ylabel={Arrivées (milliers de touristes)},
symbolic x coords={2015,2017,2019,2020,2022},
xtick=data,
ymin=0, ymax=1200,
nodes near coords, nodes near coords style={font=\small},
ymajorgrids, grid style={popline},
]
\addplot[fill=popblue, draw=popdark] coordinates {(2015,800) (2017,950) (2019,1000) (2020,350) (2022,700)};
\end{axis}
\end{tikzpicture}
\legende{Diagramme en bâtons : évolution estimée des arrivées de touristes internationaux au Cameroun (ordres de grandeur arrondis). On observe une hausse jusqu'en 2019, un effondrement en 2020 (pandémie de Covid-19), puis une reprise à partir de 2022 (année de la CAN).}
\end{popfigure}

\begin{methode}[Interpréter un tableau statistique touristique]
\begin{enumerate}
\item \textbf{Lire le titre, l'unité et la source} : de quoi parle-t-on (arrivées ? recettes ?), en quelle unité (milliers de personnes, milliards de FCFA) ?
\item \textbf{Repérer la tendance générale} : hausse, baisse ou stagnation entre la première et la dernière date.
\item \textbf{Calculer des variations} : taux de variation $= \dfrac{\text{valeur finale} - \text{valeur initiale}}{\text{valeur initiale}} \times 100$.
\item \textbf{Identifier les ruptures} : années de chute ou de bond exceptionnel.
\item \textbf{Relier les chiffres aux événements} : crises (pandémie, insécurité), grands événements (CAN), politiques de promotion.
\item \textbf{Conclure avec nuance} : dégager ce que les chiffres montrent... et ce qu'ils ne montrent pas.
\end{enumerate}
\end{methode}

\begin{exoresolu}[Calculer et interpréter une variation touristique]
Énoncé : à partir du tableau ci-dessus, calcule le taux de variation des arrivées entre 2015 et 2019, puis entre 2019 et 2020. Interprète ces deux résultats.
\tcblower
Solution détaillée :
\begin{align*}
\text{Variation 2015-2019} &= \frac{1\,000 - 800}{800} \times 100 = +25\,\% \\
\text{Variation 2019-2020} &= \frac{350 - 1\,000}{1\,000} \times 100 = -65\,\%
\end{align*}
Interprétation : entre 2015 et 2019, les arrivées progressent d'un quart (+25 \%), traduisant une dynamique favorable (promotion du pays, amélioration de l'offre hôtelière). Entre 2019 et 2020, elles s'effondrent des deux tiers (-65 \%) : cette rupture brutale ne s'explique pas par une perte d'attractivité du Cameroun, mais par la pandémie de Covid-19, qui a fermé les frontières et cloué les avions au sol dans le monde entier. La reprise de 2022 (environ 700 000 arrivées) coïncide avec la levée des restrictions et avec l'effet d'attraction de la CAN. Conclusion : le secteur est dynamique mais très \emph{sensible aux chocs extérieurs}.
\end{exoresolu}

\begin{attention}[Ne surestime pas le poids du tourisme camerounais]
Erreur fréquente dans les copies : présenter le tourisme comme « le moteur de l'économie camerounaise ». C'est faux. Le tourisme ne pèse que quelques pourcents du PIB et le Cameroun accueille environ un million de visiteurs par an, contre plus de 10 millions pour le Maroc et près de 2 millions pour le Kenya (ordres de grandeur). Le Cameroun reste un pays touristique \emph{en développement} : ses potentialités sont énormes, mais leur exploitation est encore limitée. Au Baccalauréat, la nuance est une qualité : écris « secteur prometteur mais encore marginal », jamais « secteur dominant ».
\end{attention}

\begin{exercice}[À toi de jouer]
1) Calcule le taux de variation des arrivées entre 2020 et 2022 à partir du tableau. Que traduit ce rebond ? \par
2) Pourquoi dit-on que le tourisme camerounais est « vulnérable aux chocs extérieurs » ? Appuie ta réponse sur deux exemples tirés du tableau. \par
3) Rédige un commentaire de cinq lignes du diagramme en bâtons ci-dessus, en suivant la méthode (tendance, rupture, explication, nuance).
\end{exercice}

\begin{corrige}[Exercice]
1) Variation 2020-2022 $= \dfrac{700 - 350}{350} \times 100 = +100\,\%$ : les arrivées doublent. Ce rebond traduit la reprise post-Covid et l'effet d'attraction de la CAN 2021 (jouée début 2022). \par
2) Le secteur dépend de facteurs qu'il ne contrôle pas : la pandémie de 2020 (-65 \% d'arrivées) et les crises sécuritaires (insécurité dans l'Extrême-Nord, crise anglophone) qui détournent les visiteurs malgré l'attractivité des sites. \par
3) Exemple de commentaire : « De 2015 à 2019, les arrivées internationales au Cameroun croissent régulièrement, passant d'environ 800 000 à 1 million (+25 \%). L'année 2020 marque une rupture brutale (-65 \%), conséquence de la pandémie de Covid-19. La reprise s'amorce en 2022 avec environ 700 000 arrivées, stimulée par la CAN. Toutefois, le niveau de 2019 n'est pas encore retrouvé, ce qui confirme que le tourisme camerounais, bien que dynamique, demeure fragile. »
\end{corrige}

\begin{aretenir}[L'essentiel de la section]
\begin{itemize}
\item Le tourisme rapporte des \cle{devises} (l'un des premiers pourvoyeurs hors pétrole) et des recettes hôtelières, mais ne représente que \emph{quelques pourcents} du PIB.
\item Il crée des emplois \cle{directs} (hôtellerie, restauration, guides, transport) et \cle{indirects} (artisanat, agriculture, commerce), et contribue à la \cle{diversification} de l'économie.
\item Sur le plan \cle{politique}, il fait rayonner l'image « Afrique en miniature » et sert la diplomatie par les grands événements (CAN 2021).
\item Sur le plan \cle{social}, il valorise les cultures, rémunère les riverains des parcs et permet aux Camerounais de découvrir leur pays, renforçant ainsi l'\cle{intégration nationale}.
\item L'analyse statistique montre un secteur dynamique mais vulnérable aux chocs (Covid-19 en 2020).
\end{itemize}
\end{aretenir}

Ainsi, le tourisme camerounais, sans être le géant économique que certains imaginent, joue un rôle stratégique : il rapporte des devises, emploie, fait rayonner et unit le pays. Reste une question décisive : si les potentialités sont si grandes et les bénéfices si réels, pourquoi le secteur reste-t-il modeste ? C'est précisément l'objet de la prochaine section, consacrée aux problèmes du tourisme camerounais.

\coursec{Les problèmes du secteur touristique camerounais}

Nous venons de voir que le tourisme constitue pour le Cameroun un secteur d'importance économique, politique et sociale : il rapporte des devises, crée des emplois et valorise l'image du pays à l'étranger. Pourtant, un paradoxe saute aux yeux : le pays qui se présente fièrement comme « l'Afrique en miniature » n'accueille qu'un nombre modeste de touristes internationaux — de l'ordre de quelques centaines de milliers de visiteurs par an (ordre de grandeur, les statistiques variant selon les sources et les années), loin derrière le Maroc, l'Afrique du Sud ou même le Kenya. Comment expliquer cet écart entre des potentialités énormes et une fréquentation faible ? La réponse tient en un mot : les \cle{problèmes} qui freinent le secteur. Cette section les analyse un à un — déficit de formation des acteurs, insuffisance des structures d'accueil, insuffisance des voies de communication, insécurité — puis montre leurs conséquences, avant de te guider dans la construction d'un croquis de synthèse.

\begin{definition}[Problème du secteur touristique]
On appelle \cle{problème du secteur touristique} tout facteur — humain, matériel, sécuritaire ou organisationnel — qui limite l'accueil des visiteurs, dégrade leur expérience ou décourage leur venue, et empêche ainsi la valorisation des \cle{sites touristiques} (lieux remarquables par leur paysage, leur faune, leur culture ou leur histoire, aménagés ou aménageables pour la visite).
\end{definition}

\courssub{Le déficit de formation des acteurs}

\textbf{Intuition.} Imagine un touriste allemand qui débarque au parc national de la Bénoué. Il veut comprendre la vie des lions, des hippopotames et des antilopes. Si son guide ne parle ni anglais ni allemand, ne connaît pas les comportements des animaux et ne sait pas organiser un safari en sécurité, la visite sera décevante — et le touriste le dira à son retour. La qualité d'un secteur touristique repose d'abord sur la qualité de ses \cle{hommes}.

\textbf{Constat rigoureux.} Le Cameroun souffre d'un déficit de formation professionnelle dans les métiers du tourisme :

\begin{itemize}
\item \cle{Guides professionnels} insuffisants et souvent autodidactes : beaucoup d'accompagnateurs dans les parcs du septentrion (Waza, Bénoué, Bouba Ndjida) ou sur les hauts plateaux de l'Ouest n'ont reçu aucune formation certifiée en accueil, en langues étrangères, en interprétation du patrimoine naturel et culturel ou en secourisme.
\item \cle{Personnel hôtelier qualifié} en nombre limité : réceptionnistes, chefs cuisiniers, gouvernantes et managers formés aux standards internationaux se concentrent dans les grands hôtels de Yaoundé et de Douala ; ailleurs, l'hôtellerie repose sur un personnel formé « sur le tas ».
\item \cle{Écoles de tourisme} trop peu nombreuses : les filières hôtelières et touristiques existent dans quelques lycées techniques, instituts privés et établissements universitaires (à Douala, Yaoundé, Buea notamment), mais leur capacité d'accueil reste très inférieure aux besoins, et les formations sont parfois mal adaptées aux réalités du terrain (écotourisme, guidage de safari, gestion de sites).
\end{itemize}

\begin{attention}[Piège fréquent]
Ne confonds pas « manque de personnel » et « manque de personnel \emph{qualifié} ». Le Cameroun ne manque pas de jeunes désireux de travailler dans le tourisme ; il manque de \textbf{structures de formation} capables de leur donner les compétences exigées par une clientèle internationale. Au Bac, précise toujours la nature exacte du déficit.
\end{attention}

\courssub{L'insuffisance des structures d'accueil}

\textbf{Intuition.} Un touriste qui a marché toute la journée dans la réserve de faune du Dja veut le soir une chambre propre, de l'eau courante, de l'électricité et une connexion téléphonique. Si le seul hébergement disponible est une case sans confort, il écourtera son séjour — ou ne viendra pas du tout.

\textbf{Constat rigoureux.} L'offre d'hébergement camerounaise est à la fois \textbf{insuffisante} et \textbf{très inégalement répartie} :

\begin{itemize}
\item Les \cle{hôtels de standing} (catégorie internationale, 3 à 5 étoiles) se concentrent massivement à \textbf{Yaoundé}, \textbf{Douala} et, dans une moindre mesure, \textbf{Kribi} — villes qui doivent cette concentration au \cle{tourisme d'affaires} (clientèle venue pour des congrès, des missions commerciales ou des séminaires, et non pour les loisirs).
\item À l'inverse, les grandes zones touristiques du \textbf{Nord} (parcs de Waza, de la Bénoué, de Bouba Ndjida, paysages kapsiki de Rhumsiki) et de l'\textbf{Est} (réserves du Dja, de Lobéké) ne disposent que de quelques campements et lodges, souvent vétustes ou fermés, très en deçà des standards attendus.
\item Les équipements complémentaires (restaurants, centres d'information touristique, stations-service, dispensaires, guichets bancaires) suivent la même logique de concentration urbaine.
\end{itemize}

\begin{exemplebox}[Un contraste révélateur]
L'essentiel des chambres d'hôtels de catégorie internationale du Cameroun se trouve à Yaoundé et à Douala, les deux métropoles politique et économique. Or les parcs animaliers du septentrion — première vitrine du tourisme camerounais — se situent à plus de \SI{1000}{km} de là et manquent cruellement d'hébergements de standing. Résultat : le touriste de safari doit accepter un confort sommaire, tandis que les belles chambres climatisées accueillent surtout des hommes d'affaires qui ne visiteront jamais Waza.
\end{exemplebox}

\courssub{L'insuffisance des voies de communication}

\textbf{Intuition.} Un site magnifique mais inaccessible est un site sans touristes. La distance ne se mesure pas en kilomètres mais en \emph{heures de route} et en \emph{difficulté du trajet}.

\textbf{Constat rigoureux.} Le réseau de transport camerounais pénalise l'accès aux sites :

\begin{itemize}
\item \textbf{Routes dégradées ou inachevées} : les axes menant aux grandes zones touristiques sont en mauvais état ou en cours de bitumage. La route Yaoundé--Ngaoundéré, porte d'entrée du septentrion, est longue et fatigante ; les liaisons vers l'Est (vers la réserve du Dja ou la Boumba-et-Ngoko) comptent parmi les plus difficiles du pays.
\item \textbf{Pistes impraticables en saison des pluies} : à l'intérieur et autour des parcs, les pistes en latérite deviennent boueuses ou coupées par les crues ; certaines zones de l'Est et du Sud sont quasi isolées plusieurs mois par an — paradoxalement pendant les vacances d'été européennes.
\item \textbf{Faiblesse du transport aérien intérieur} : les lignes régulières desservent surtout Yaoundé, Douala, Garoua, Maroua et Ngaoundéré ; les fréquences sont limitées, les tarifs élevés, et aucune desserte ne touche les petits sites. Le chemin de fer (Transcamerounais, Douala--Yaoundé--Ngaoundéré) soulage l'axe du septentrion mais reste lent et vétuste.
\end{itemize}

\courssub{L'insécurité}

\textbf{Intuition.} Aucun voyageur ne choisit une destination où il craint pour sa vie. L'insécurité, même localisée, ternit l'image de tout un pays, car les médias internationaux retiennent le nom « Cameroun » sans nuance géographique.

\textbf{Constat rigoureux.} L'insécurité touristique au Cameroun est \textbf{réelle mais localisée} :

\begin{itemize}
\item \textbf{Extrême-Nord} : les incursions de la secte \cle{Boko Haram} (enlèvements, attentats) depuis 2014 ont rendu dangereuse la frange frontalière du Nigeria et du lac Tchad, affectant l'accès au parc de Waza, aux monts Mandara et à Rhumsiki.
\item \textbf{Nord-Ouest et Sud-Ouest} : la crise anglophone (depuis 2016-2017), avec ses affrontements armés, ses « villes mortes » et ses blocages de routes, a pratiquement fermé au tourisme des sites comme le mont Cameroun, les ring roads de Bamenda ou les chutes de Menchum.
\item \textbf{Frontière orientale} : les exactions de groupes armés venus de Centrafrique et les coupeurs de route sur certains axes de l'Est et de l'Adamaoua découragent les voyages terrestres.
\item \textbf{Délinquance ordinaire} : vols et agressions sur certains axes routiers et dans les grandes villes ajoutent au sentiment d'insécurité.
\end{itemize}

\begin{attention}[Erreur fréquente à l'examen]
Ne présente \textbf{jamais} l'insécurité comme généralisée à tout le Cameroun : c'est factuellement faux et cela coûte des points. Elle touche des \textbf{zones précises} (Extrême-Nord, Nord-Ouest, Sud-Ouest, frange orientale). L'Ouest, le Littoral, le Centre et le Sud restent globalement sûrs et accueillent l'essentiel des visiteurs. Raisonne toujours en géographe : localise les faits sur une carte mentale.
\end{attention}

\begin{savaistu}[Waza, joyau en souffrance]
Le parc national de Waza, créé en 1934 comme réserve, érigé en parc national en 1968 puis classé réserve de biosphère par l'UNESCO en 1979, était dans les années 1970-1980 l'une des destinations de safari les plus réputées d'Afrique centrale : lions, éléphants, girafes, antilopes et des centaines d'espèces d'oiseaux y attiraient des milliers de visiteurs. Les attaques de Boko Haram dans l'Extrême-Nord à partir de 2014 ont fait chuter drastiquement sa fréquentation et conduit à la fermeture temporaire de plusieurs campements. La reprise progressive observée ces dernières années montre que la sécurité est la condition première du tourisme.
\end{savaistu}

\courssub{Les autres freins et les conséquences}

\textbf{Autres freins (complément utile au Bac).}

\begin{itemize}
\item \textbf{Promotion insuffisante} : le Cameroun investit peu dans la publicité touristique à l'étranger ; sa présence dans les grands salons internationaux du tourisme est irrégulière, et l'image « Afrique en miniature » reste mal vendue face au Kenya ou à l'Afrique du Sud.
\item \textbf{Coûts élevés} : le prix du visa, le coût des billets d'avion vers Yaoundé ou Douala et le tarif des prestations sur place rendent la destination chère par rapport à sa concurrence.
\item \textbf{Dégradation des sites} : braconnage dans les parcs, déforestation, déchets sur les plages, manque d'entretien des monuments (palais des sultans, cases obus des monts Mandara) altèrent le produit touristique lui-même.
\item \textbf{Lourdeurs administratives} : tracasseries aux frontières et contrôles multiples sur les routes fatiguent les visiteurs.
\end{itemize}

\textbf{Conséquences.} L'addition de ces problèmes produit trois effets majeurs :

\begin{enumerate}
\item une \textbf{faible fréquentation internationale} : le Cameroun accueille de l'ordre de quelques centaines de milliers de visiteurs par an (ordre de grandeur à manier avec prudence, les statistiques variant selon les sources et les années), dont une forte part de voyageurs d'affaires et de la diaspora, et non de touristes de loisirs ;
\item une \textbf{concentration spatiale} du tourisme sur le triangle Yaoundé--Douala--Kribi, au détriment du septentrion et de l'Est ;
\item une \textbf{sous-exploitation des potentialités} : l'« Afrique en miniature » reste un slogan plus qu'une réalité économique.
\end{enumerate}

\begin{exoresolu}[Analyser un tableau statistique]
\textbf{Énoncé.} Un document indique que sur 100 chambres d'hôtels de catégorie internationale recensées au Cameroun, environ 80 se trouvent à Yaoundé et Douala, 8 à Kribi et dans le reste du Littoral et du Sud, et 12 dans l'ensemble des sept autres régions. 1) Calcule la part des sept autres régions (qui incluent le septentrion et l'Est) dans l'hébergement de standing. 2) Dégage un problème majeur du secteur touristique camerounais.
\tcblower
\textbf{Solution.} 1) Les sept autres régions (qui incluent le Nord, l'Extrême-Nord, l'Adamaoua et l'Est, zones des grands parcs) ne disposent que de $\dfrac{12}{100} = 12\,\%$ de l'offre, alors que Yaoundé et Douala en concentrent $\dfrac{80}{100} = 80\,\%$, soit près de sept fois plus. 2) Ce tableau révèle l'\textbf{insuffisance et l'inégale répartition des structures d'accueil} : l'hébergement de standing suit la logique du tourisme d'affaires (métropoles) et non celle des potentialités touristiques (parcs, plages, sites culturels). Cette contradiction explique la sous-exploitation des atouts du septentrion et de l'Est.
\end{exoresolu}

\begin{popfigure}
\begin{center}
\includegraphics[width=\linewidth,height=9.5cm,keepaspectratio]{N01/cmr_map.jpg}
\end{center}
\legende{Fond de carte du Cameroun (régions, villes, routes, voie ferrée) servant à représenter l'inégale répartition des infrastructures touristiques. On y repère Douala et Yaoundé, où se concentrent les hôtels de standing ; le Grand Nord (Waza, au nord de Maroua) et l'Adamaoua (Bénoué, près de Garoua) loin de ces métropoles ; un réseau routier mal développé vers l'Est (Dja) ; et les régions affectées par l'insécurité (Extrême-Nord, Nord-Ouest, Sud-Ouest). Les hôtels, les parcs et les zones d'insécurité ne figurent pas sur l'image : ils sont à reporter par l'élève. \\ Source : Wikimedia Commons, CIA, domaine public.}
\end{popfigure}

\begin{popfigure}
\begin{center}
\begin{tikzpicture}[scale=1.0]
% Diagramme circulaire (parts indicatives)
\fill[popblue!70] (0,0) -- (0:2.2) arc (0:108:2.2) -- cycle;
\fill[poporange!75] (0,0) -- (108:2.2) arc (108:198:2.2) -- cycle;
\fill[popgreen!70] (0,0) -- (198:2.2) arc (198:288:2.2) -- cycle;
\fill[poppink!70] (0,0) -- (288:2.2) arc (288:360:2.2) -- cycle;
\node[font=\small\bfseries, white] at (54:1.3) {30\,\%};
\node[font=\small\bfseries, white] at (153:1.3) {25\,\%};
\node[font=\small\bfseries, white] at (243:1.3) {25\,\%};
\node[font=\small\bfseries, white] at (324:1.3) {20\,\%};
% Légende
\node[anchor=west, font=\small] at (2.8,1.5) {\textcolor{popblue!70}{\rule{0.25cm}{0.25cm}} Structures d'accueil insuffisantes};
\node[anchor=west, font=\small] at (2.8,0.9) {\textcolor{poporange!75}{\rule{0.25cm}{0.25cm}} Voies de communication défaillantes};
\node[anchor=west, font=\small] at (2.8,0.3) {\textcolor{popgreen!70}{\rule{0.25cm}{0.25cm}} Déficit de formation des acteurs};
\node[anchor=west, font=\small] at (2.8,-0.3) {\textcolor{poppink!70}{\rule{0.25cm}{0.25cm}} Insécurité (zones localisées)};
\end{tikzpicture}
\end{center}
\legende{Hiérarchisation indicative des principaux problèmes du secteur touristique camerounais (parts estimées à titre pédagogique, sans valeur statistique officielle). L'insuffisance de l'accueil et des communications constitue le frein structurel majeur, devant la formation et l'insécurité.}
\end{popfigure}

\courssub{TP guidé : construire un croquis de l'inégale répartition des infrastructures touristiques}

\begin{methode}[Construire un croquis géographique en cinq étapes]
\begin{enumerate}
\item \textbf{Simplifier les contours} : trace le contour du Cameroun sous forme d'un polygone simplifié (une dizaine de segments suffisent) ; ajoute la flèche du nord et un titre précis annonçant le sujet et l'espace étudié.
\item \textbf{Choisir les figurés} : retiens 4 à 6 figurés maximum, \emph{proportionnels à l'importance des faits} : points ou carrés nombreux là où les hôtels se concentrent, rares ailleurs ; traits pleins pour les axes en bon état, pointillés pour les routes dégradées ; hachures pour les zones d'insécurité.
\item \textbf{Placer les faits} : localise d'abord les repères (Douala, Yaoundé, Kribi, Maroua, Garoua, Ngaoundéré), puis les parcs (Waza, Bénoué, Dja), puis les infrastructures et les contraintes.
\item \textbf{Rédiger la légende} : ordonne-la du plus important au moins important, numérote les figurés, et fais correspondre \emph{exactement} chaque figuré de la légende à ceux de la carte.
\item \textbf{Vérifier la cohérence} : le croquis doit « parler » seul — un lecteur doit voir immédiatement la concentration sud-ouest des infrastructures et l'isolement des zones touristiques du septentrion et de l'Est.
\end{enumerate}
\end{methode}

\begin{exercice}[À toi de jouer]
À partir du fond de carte proposé ci-dessus et de tes connaissances : 1) Reproduis le contour simplifié du Cameroun et place-y les trois métropoles de l'hébergement de standing. 2) Ajoute les trois grands parcs étudiés et les axes routiers qui les desservent, en distinguant par tes figurés les axes praticables des axes dégradés. 3) Rédige une légende complète et ordonnée, puis formule en deux phrases le problème géographique que ton croquis met en évidence.
\end{exercice}

\begin{corrige}[Exercice 1]
1) Le contour est tracé en polygone simplifié ; Douala est placée sur la côte, Yaoundé légèrement à l'est et à l'intérieur, Kribi au sud sur l'océan. 2) Waza figure à l'extrême-nord, la Bénoué au nord-centre, le Dja au sud-est ; l'axe Yaoundé--Ngaoundéré--Garoua--Maroua et l'axe vers l'Est sont dessinés en pointillés (routes dégradées), tandis que la liaison Douala--Yaoundé est en trait plein. 3) La légende ordonne : hôtels de standing, parcs, routes en bon état, routes dégradées, zones d'insécurité. Problème dégagé : les infrastructures d'accueil se concentrent dans les métropoles du sud-ouest, tandis que les grandes potentialités touristiques du septentrion et de l'Est, éloignées et mal desservies, restent sous-valorisées.
\end{corrige}

\begin{aretenir}[L'essentiel de la section]
Les problèmes du tourisme camerounais sont de quatre ordres : \cle{déficit de formation} des acteurs (guides, hôteliers, écoles), \cle{insuffisance des structures d'accueil} (concentrées à Yaoundé, Douala, Kribi), \cle{insuffisance des voies de communication} (routes dégradées, pistes saisonnières, aérien faible) et \cle{insécurité localisée} (Extrême-Nord, Nord-Ouest, Sud-Ouest, Est). S'y ajoutent la faible promotion, les coûts élevés et la dégradation des sites. Conséquences : faible fréquentation internationale, concentration spatiale du tourisme et sous-exploitation des potentialités — autant de défis que la valorisation du secteur, objet de la prochaine section, devra relever.
\end{aretenir}

\coursec{Perspectives et valorisation du tourisme pour l'intégration nationale}

Nous venons de le voir : malgré des potentialités exceptionnelles, le tourisme camerounais reste freiné par le déficit de formation des acteurs, l'insuffisance des structures d'accueil et des voies de communication, ainsi que par l'insécurité dans certaines régions. Faut-il pour autant désespérer ? Non. Le diagnostic posé, il reste à dessiner les \cle{perspectives} : comment transformer ces atouts endormis en un véritable moteur de développement et d'\cle{intégration nationale} ? C'est l'objet de cette dernière partie, qui bouclera notre étude du tourisme au Cameroun.

\courssub{Les pistes de relance du secteur touristique}

\courssub{Former les acteurs et améliorer les infrastructures}

Commençons par une intuition simple : un touriste déçu par l'accueil, la propreté d'un hôtel ou la qualité d'un guide ne reviendra pas, et surtout, il le dira autour de lui. La \cle{formation des acteurs} est donc la première condition de la relance. Le Cameroun dispose déjà d'établissements spécialisés : des \cle{écoles hôtelières} et des instituts de formation aux métiers du tourisme (à Douala, Yaoundé, Buea notamment), qui forment réceptionnistes, cuisiniers, guides et gestionnaires d'hébergements. L'enjeu est de renforcer ces filières, de les adapter aux normes internationales et de professionnaliser les guides des sites comme Waza ou Foumban.

La deuxième piste concerne les \cle{infrastructures} :
\begin{itemize}
\item les \textbf{infrastructures routières} : désenclaver les sites touristiques exige de poursuivre le bitumage des axes qui y mènent (route de l'Unité vers le septentrion, accès à Kribi et à la côte, routes des hauts plateaux de l'Ouest) ;
\item les \textbf{infrastructures hôtelières} : augmenter la capacité d'accueil (hôtels, auberges, campements, écolodges) et en améliorer la qualité, tant dans les grandes villes qu'à proximité des parcs nationaux ;
\item les \textbf{infrastructures aériennes et numériques} : moderniser les aéroports de Douala et Yaoundé-Nsimalen, développer la réservation en ligne et la couverture internet des zones touristiques.
\end{itemize}

\begin{exemplebox}[Un objectif chiffré : le tourisme dans la vision « Cameroun émergent 2035 »]
Dans sa stratégie de développement (Document de Stratégie pour la Croissance et l'Emploi, puis Stratégie Nationale de Développement 2020-2030), l'État camerounais affiche l'ambition de faire du tourisme un \textbf{secteur clé} de la vision « \cle{Cameroun émergent 2035} ». L'objectif est d'accueillir plusieurs centaines de milliers, puis à terme plus d'un million de visiteurs internationaux par an (ordre de grandeur indicatif : les arrivées internationales tournaient autour de 500 000 à 900 000 par an dans les années 2010, loin du million). Cet objectif suppose des investissements massifs dans la formation, l'hôtellerie et les transports.
\end{exemplebox}

\courssub{Promouvoir le pays à l'étranger}

Un produit, même excellent, ne se vend pas tout seul : il faut le faire connaître. Le Cameroun mène donc une politique de \cle{promotion touristique} à l'étranger :
\begin{itemize}
\item la campagne de communication « \cle{Afrique en miniature} », slogan qui résume en trois mots l'argument maître du pays : toute la diversité de l'Afrique (mer, savane, forêt, cultures) concentrée en un seul État ;
\item la participation régulière aux \textbf{salons internationaux du tourisme} (par exemple les grandes foires du tourisme de Berlin \emph{ITB} ou de Londres \emph{WTM}), où les opérateurs camerounais présentent leurs offres aux agences de voyage européennes, américaines et asiatiques ;
\item l'organisation d'événements culturels et sportifs internationaux (festivals, compétitions comme la Coupe d'Afrique des Nations de football accueillie en 2022), vitrines médiatiques du pays.
\end{itemize}

\courssub{Développer l'écotourisme et le tourisme communautaire}

Troisième grande piste : miser sur ce que le Cameroun possède de plus rare, sa nature, tout en faisant participer les populations locales. L'\cle{écotourisme} est un tourisme responsable, centré sur l'observation de la nature et le respect des écosystèmes. Le \cle{tourisme communautaire} va plus loin : les villages riverains des parcs participent à l'accueil des visiteurs (guides locaux, campements villageois, artisanat) et perçoivent une \textbf{part des revenus} générés (droits d'entrée des parcs, taxes touristiques reversées).

\begin{savaistu}[Concilier conservation et développement local]
Le tourisme communautaire permet aux villages proches des parcs de tirer des revenus de la \textbf{protection} de la nature : autour de parcs comme Korup (Sud-Ouest) ou Campo-Ma'an (Sud), une partie des recettes touristiques finance des puits, des écoles ou des dispensaires villageois. Résultat : les populations deviennent les premières gardiennes de la faune, car un gorille vivant rapporte plus, sur la durée, qu'un gorille braconné. Conservation et développement local cessent d'être ennemis.
\end{savaistu}

\courssub{Sécuriser les zones touristiques : la condition du retour des visiteurs}

Aucune promotion ne peut compenser la peur. La \cle{sécurisation} des zones touristiques est la condition préalable du retour des visiteurs, en particulier :
\begin{itemize}
\item dans l'\textbf{Extrême-Nord}, où les exactions de Boko Haram ont fortement réduit la fréquentation du parc de Waza et des paysages des monts Mandara ;
\item dans le \textbf{Nord-Ouest} et le \textbf{Sud-Ouest}, où la crise anglophone affecte l'accès aux sites (Korup, mont Cameroun) ;
\item sur certains axes où la criminalité routière décourage les déplacements.
\end{itemize}
Le retour progressif de la sécurité, la présence d'escortes et la stabilisation politique de ces régions sont donc des investissements touristiques autant que sécuritaires.

\courssub{Le tourisme, facteur d'intégration nationale}

Venons-en au cœur de la famille de situations de cette leçon : l'\cle{intégration nationale}. En quoi le tourisme peut-il cimenter l'unité du Cameroun ?

\begin{propriete}[Le tourisme comme levier d'intégration nationale]
Le tourisme contribue à l'intégration nationale par trois mécanismes complémentaires :
\begin{enumerate}
\item la \textbf{mise en valeur de toutes les régions} : chaque région possède un atout (mer au Littoral et au Sud-Ouest, safari au septentrion, cultures des hauts plateaux de l'Ouest et du Nord-Ouest, forêt du plateau du Sud et de l'Est) ; le développement touristique réduit les déséquilibres régionaux en créant des emplois et des revenus hors des deux métropoles ;
\item le \textbf{brassage des populations} : le tourisme intérieur (Camerounais visitant leur propre pays) et les flux de travailleurs du secteur font circuler les personnes, les langues et les habitudes, tissant du lien national ;
\item la \textbf{fierté nationale} : la valorisation du patrimoine naturel et culturel (palais des sultanats, parcs classés, gastronomie) nourrit un sentiment d'appartenance commune à un pays exceptionnel.
\end{enumerate}
\end{propriete}

\begin{exemplebox}[Des exemples localisés d'intégration par le tourisme]
Un hôtel à Kribi emploie des Camerounais de toutes les régions ; le festival des arts et de la culture de Foumban attire des visiteurs de Yaoundé, Douala ou Maroua ; la route de l'Unité, construite pour relier le Nord au Sud, est aussi l'axe des safaris vers Waza. Tourisme et unité nationale avancent ensemble.
\end{exemplebox}

\courssub{Synthèse de la leçon et bilan des savoir-faire}

\begin{aretenir}[Synthèse : l'essentiel de la leçon]
Le Cameroun possède des \textbf{potentialités touristiques énormes} : l'attrait de la mer (Kribi, Limbé), le safari dans les parcs nationaux du septentrion (Waza, Bénoué, Bouba Ndjida), le dynamisme culturel des hauts plateaux (Foumban, Bafoussam, chefferies bamiléké et bamoun) et l'écotourisme dans le plateau forestier (Korup, Campo-Ma'an, Dja). Le tourisme a une \textbf{importance réelle} : économique (devises, emplois), politique (rayonnement international) et sociale (revenus, fierté). Mais le secteur est \textbf{freiné} par de nombreux problèmes : déficit de formation des acteurs, insuffisance des structures d'accueil et des voies de communication, insécurité. Les \textbf{perspectives} passent par la formation, les infrastructures, la promotion (« Afrique en miniature »), l'écotourisme communautaire et la sécurisation — autant de leviers d'intégration nationale.
\end{aretenir}

\begin{popfigure}
\begin{center}\tcbox[colback=popink,colframe=popink,coltext=white,arc=9pt,boxsep=4pt,left=6pt,right=6pt]{\bfseries\small Tourisme au Cameroun}\end{center}
\vspace{-2pt}
\begin{tcbraster}[raster columns=2,raster equal height=rows,raster column skip=6pt,raster row skip=6pt,enhanced,blankest]
\begin{tcolorbox}[enhanced,colback=white,colframe=popgreen,boxrule=0.9pt,arc=3pt,title={Potentialités},fonttitle=\bfseries\scriptsize,coltitle=white,colbacktitle=popgreen,left=4pt,right=4pt,top=2pt,bottom=2pt,boxsep=1pt,toptitle=1.5pt,bottomtitle=1.5pt,halign=left]
\begin{itemize}[nosep,leftmargin=*,itemsep=1pt,font=\scriptsize]
\item Mer : Kribi, Limbé
\item Safari : Waza, Bénoué
\item Cultures : Foumban, hauts plateaux
\item Écotourisme : Korup, Dja
\end{itemize}
\end{tcolorbox}
\begin{tcolorbox}[enhanced,colback=white,colframe=popgold,boxrule=0.9pt,arc=3pt,title={Importance},fonttitle=\bfseries\scriptsize,coltitle=white,colbacktitle=popgold,left=4pt,right=4pt,top=2pt,bottom=2pt,boxsep=1pt,toptitle=1.5pt,bottomtitle=1.5pt,halign=left]
\begin{itemize}[nosep,leftmargin=*,itemsep=1pt,font=\scriptsize]
\item Économique : devises, emplois
\item Politique : rayonnement
\item Sociale : revenus, fierté
\end{itemize}
\end{tcolorbox}
\begin{tcolorbox}[enhanced,colback=white,colframe=poporange,boxrule=0.9pt,arc=3pt,title={Problèmes},fonttitle=\bfseries\scriptsize,coltitle=white,colbacktitle=poporange,left=4pt,right=4pt,top=2pt,bottom=2pt,boxsep=1pt,toptitle=1.5pt,bottomtitle=1.5pt,halign=left]
\begin{itemize}[nosep,leftmargin=*,itemsep=1pt,font=\scriptsize]
\item Formation insuffisante
\item Routes et hôtels déficients
\item Insécurité
\end{itemize}
\end{tcolorbox}
\begin{tcolorbox}[enhanced,colback=white,colframe=poppurple,boxrule=0.9pt,arc=3pt,title={Perspectives},fonttitle=\bfseries\scriptsize,coltitle=white,colbacktitle=poppurple,left=4pt,right=4pt,top=2pt,bottom=2pt,boxsep=1pt,toptitle=1.5pt,bottomtitle=1.5pt,halign=left]
\begin{itemize}[nosep,leftmargin=*,itemsep=1pt,font=\scriptsize]
\item Écoles hôtelières
\item « Afrique en miniature »
\item Tourisme communautaire
\item Sécurisation
\end{itemize}
\end{tcolorbox}
\end{tcbraster}

\legende{Schéma-bilan de la leçon : carte mentale du tourisme au Cameroun. Quatre branches relient les potentialités, l'importance, les problèmes et les perspectives : c'est exactement la structure d'un plan de dissertation sur le sujet.}
\end{popfigure}

\begin{methode}[Réussir une dissertation de géographie sur le tourisme au Cameroun]
Face à un sujet du type « Le tourisme au Cameroun : atouts, rôle et difficultés », procède en étapes :
\begin{enumerate}
\item \textbf{Analyse du sujet} : repère les termes clés (tourisme, Cameroun) et définis-les en introduction (site touristique, tourisme d'affaires) ; annonce un plan en trois parties.
\item \textbf{Première partie — les atouts} : présente les quatre grandes potentialités (mer, safari, cultures, forêt), chacune illustrée par un exemple localisé.
\item \textbf{Deuxième partie — le rôle} : montre l'importance économique (devises, emplois), politique (image du pays) et sociale (revenus, intégration nationale).
\item \textbf{Troisième partie — les difficultés} : expose les problèmes (formation, infrastructures, insécurité) et termine par les perspectives de relance.
\item \textbf{Conclusion} : fais le bilan (potentiel énorme, exploitation faible) et ouvre sur l'émergence 2035.
\item \textbf{Règle d'or} : chaque idée doit être appuyée par un exemple précis et localisé, et si possible par un croquis ou un chiffre.
\end{enumerate}
\end{methode}

\begin{attention}[Erreur fréquente au Bac : la copie sans exemples localisés]
Beaucoup de candidats écrivent des généralités (« le Cameroun a de beaux paysages », « il y a des animaux ») sans jamais citer un lieu précis. C'est une erreur rédhibitoire en géographie. Dans toute copie du Bac sur le tourisme, tu dois obligatoirement mobiliser des \textbf{exemples précis et localisés} : le parc national de \cle{Waza} (Extrême-Nord), les plages de \cle{Kribi} (Sud), le palais des rois bamoun à \cle{Foumban} (Ouest), la réserve de biosphère du \cle{Dja} (Sud-Est, classée au patrimoine mondial de l'UNESCO) ou le parc de \cle{Korup} (Sud-Ouest). Une idée sans exemple localisé est une idée à moitié notée.
\end{attention}

\begin{exoresolu}[Exercice résolu : commentaire d'un document statistique]
\textbf{Énoncé.} Document : « Au Cameroun, les arrivées de touristes internationaux sont estimées à quelques centaines de milliers par an dans les années 2010 (ordre de grandeur : environ 500 000 à 900 000 selon les années et les sources), pour une contribution du tourisme au PIB de l'ordre de quelques pourcents seulement, alors que le pays ambitionne, dans sa vision d'émergence à l'horizon 2035, de dépasser le million de visiteurs. » À partir du document et de tes connaissances, montre que le tourisme camerounais est un secteur sous-exploité, puis propose deux pistes de relance.
\tcblower
\textbf{Solution détaillée.}
\begin{enumerate}
\item \textbf{Lecture du document} : les arrivées internationales restent modestes (moins d'un million de visiteurs) et la part du tourisme dans le PIB est faible (quelques pourcents). Or le Cameroun concentre la mer, la savane, la forêt et une grande diversité culturelle : il y a donc un \textbf{écart entre le potentiel et la réalité}, ce qui définit un secteur sous-exploité.
\item \textbf{Explication par les connaissances} : ce sous-emploi s'explique par les problèmes étudiés : déficit de formation des acteurs, insuffisance des hôtels et des routes (sites enclavés comme Waza ou Korup), insécurité dans l'Extrême-Nord et les régions anglophones.
\item \textbf{Pistes de relance} (deux attendues, au choix) : former davantage de professionnels dans les écoles hôtelières ; améliorer les infrastructures routières et hôtelières ; promouvoir la destination à l'étranger (campagne « Afrique en miniature », salons internationaux) ; développer l'écotourisme communautaire autour des parcs ; sécuriser les zones touristiques.
\item \textbf{Conclusion} : atteindre l'objectif de la vision « Cameroun émergent 2035 » suppose de lever ces freins ; le tourisme deviendrait alors un secteur clé de l'économie et un facteur d'intégration nationale.
\end{enumerate}
\end{exoresolu}

\begin{popfigure}
\begin{center}
\includegraphics[width=\linewidth,height=11cm,keepaspectratio]{N01/cmr_physique.png}
\end{center}
\legende{Fond de carte muet du Cameroun (relief, fleuves et frontières, sans noms de lieux) à compléter par les élèves : « Le Cameroun, destination touristique ». Consigne : localise et nomme quatre sites touristiques (réponses attendues : Waza au nord ; Kribi sur le littoral ; Foumban dans les hauts plateaux de l'Ouest ; Korup ou Dja en forêt du Sud), puis rédige une légende organisée en trois rubriques (atouts naturels, atouts culturels, villes-relais). \\ Source : Wikimedia Commons, Urutseg, CC0.}
\end{popfigure}

\begin{exercice}[Pour s'entraîner : construire un diagramme et interpréter un document]
\begin{enumerate}
\item À partir des données suivantes (ordres de grandeur indicatifs) : mer et côte 30 \%, safari et savane 25 \%, tourisme culturel 25 \%, écotourisme forestier 20 \% de l'attractivité touristique camerounaise, construis un diagramme circulaire (ou en barres) légendé et intitulé.
\item Rédige un commentaire de 15 lignes du document suivant : « Le Cameroun, "Afrique en miniature", reçoit moins d'un million de touristes internationaux par an, quand le Maroc ou l'Afrique du Sud en accueillent plusieurs millions. » Tu montreras le contraste entre potentiel et fréquentation, tu en proposeras des explications et tu concluras sur les perspectives.
\item Sur la carte de synthèse ci-dessus, ajoute à main levée un flux (flèche) représentant l'axe touristique Douala--Kribi et un autre représentant l'axe des safaris Douala--Garoua--Waza.
\end{enumerate}
\end{exercice}

\begin{corrige}[Exercice 1]
\begin{enumerate}
\item Le diagramme circulaire comporte quatre secteurs : 30 \% correspondent à un angle de $0{,}30 \times 360 = 108^{\circ}$ (mer et côte), 25 \% à $90^{\circ}$ (safari), 25 \% à $90^{\circ}$ (culturel), 20 \% à $72^{\circ}$ (écotourisme). Vérifie que la somme des angles vaut bien $360^{\circ}$. Le titre pourrait être : « Répartition indicative de l'attractivité touristique du Cameroun par type d'atout ».
\item Plan du commentaire : (a) constat : faible fréquentation malgré un potentiel exceptionnel (mer, parcs du septentrion, hauts plateaux, forêt) ; (b) explications : déficit de formation, infrastructures insuffisantes, insécurité, promotion limitée ; (c) perspectives : formation, investissements, campagne « Afrique en miniature », écotourisme communautaire, sécurisation ; conclusion sur l'intégration nationale et l'émergence 2035.
\item Les flèches doivent partir de Douala : l'une vers le sud-est en direction de Kribi (côte), l'autre vers le nord via Garoua jusqu'à Waza (axe de la route de l'Unité). L'épaisseur des flèches peut symboliser l'importance relative des flux.
\end{enumerate}
\end{corrige}

\begin{aretenir}[Bilan des savoir-faire acquis dans cette leçon]
À l'issue de cette leçon, tu sais :
\begin{itemize}
\item \textbf{observer, localiser et décrire} les grands sites touristiques camerounais sur une carte ;
\item \textbf{représenter et cartographier} : légender une carte, construire un croquis de répartition des atouts touristiques ;
\item \textbf{construire un diagramme} (barres, circulaire) à partir de données statistiques et vérifier les calculs de pourcentages et d'angles ;
\item \textbf{inventorier, classer et interpréter} des documents variés (cartes, photographies, tableaux statistiques) pour rédiger un commentaire structuré ;
\item \textbf{mobiliser des exemples localisés} (Waza, Kribi, Foumban, Korup) dans une dissertation du Bac.
\end{itemize}
\end{aretenir}

\newpage

\coursec{Activité d'intégration}

\coursec{Le tourisme au Cameroun}

\courssub{Définitions et cadre général}

\begin{definition}[Tourisme]
Ensemble des activités des personnes qui voyagent et séjournent dans des lieux situés en dehors de leur environnement habituel pour une durée n'excédant pas une année consécutive, à des fins de loisirs, d'affaires ou autres.
\end{definition}

\begin{definition}[Site touristique]
Lieu présentant un intérêt particulier (naturel, culturel, historique, religieux) qui attire les visiteurs. Au Cameroun, on distingue les sites naturels (plages, parcs, montagnes, chutes) et les sites culturels (chefferies, musées, festivals, architecture).
\end{definition}

\begin{definition}[Tourisme d'affaires]
Segment du tourisme lié aux déplacements professionnels : congrès, conférences, séminaires, salons, voyages d'incitation. Il génère des retombées élevées (hôtellerie haut de gamme, restauration, transport). Au Cameroun, il se concentre à Yaoundé (Palais des Congrès, hôtels Mont Fébé, Hilton), Douala (Akwa Palace, Sawa Hotel, Centre International des Conférences) et Kribi.
\end{definition}

\begin{aretenir}[Cadre national]
Le tourisme relève du Ministère du Tourisme et des Loisirs (MINTOUR). La politique nationale vise à faire du Cameroun une \og destination de premier plan \fg{} en Afrique centrale, en s'appuyant sur le slogan \og l'Afrique en miniature \fg{}. Le secteur est inscrit dans la Stratégie Nationale de Développement (SND30) comme levier de croissance et d'intégration territoriale.
\end{aretenir}

\courssub{Les potentialités touristiques du Cameroun}

Le programme officiel identifie \textbf{quatre grandes potentialités}, réparties sur l'ensemble du territoire, ce qui fait du Cameroun une \og Afrique en miniature \fg{}.

\begin{propriete}[Les quatre potentialités majeures]
\begin{enumerate}
\item \textbf{L'attrait de la mer (tourisme balnéaire)} : Littoral atlantique (400 km). Plages de sable fin (Kribi, Limbé, Ebodjé, Campo), criques, chutes de la Lobé (Kribi) se jetant dans l'océan, pêche sportive, observation des tortues marines. Infrastructures hôtelières en développement (Kribi, Limbé).
\item \textbf{Le safari dans les parcs nationaux du Septentrion (tourisme cynégétique et de vision)} : Régions du Nord, de l'Adamaoua et de l'Extrême-Nord. Parcs majeurs : Waza (Extrême-Nord, faune sahélienne : éléphants, girafes, lions), Bénoué (Nord, hippopotames, buffles, derby de Buffon), Bouba Ndjida (Adamaoua, rhinocéros, éléphants), Faro (Nord). Activités : photo-safari, chasse sportive (zones cynégétiques gérées).
\item \textbf{Le dynamisme culturel des Hauts Plateaux (tourisme culturel)} : Région de l'Ouest (Bamileke, Bamoun) et Nord-Ouest. Chefferies traditionnelles classées UNESCO (Bafut, Bandjoun, Baham, Foumban), musées (Musée des Civilisations à Dschang, Palais des Rois Bamoun à Foumban), architecture traditionnelle (cases à toits coniques), festivals (Ngoun, Medumba, Nyem-Nyem), artisanat d'art (bronze, perles, tissus ndop, sculpture sur bois).
\item \textbf{L'écotourisme dans le plateau forestier (tourisme vert)} : Régions du Sud, de l'Est, du Centre. Réserve de biosphère du Dja (UNESCO, gorilles, chimpanzés, éléphants de forêt), Parc national de Campo-Ma'an (Sud, biodiversité côtière), Parc de Lobéké (Est, gorilles, clairières salines), chutes d'Ekom-Nkam (Moungo), Mont Cameroun (région Sud-Ouest, randonnée, ascension, course de l'Espoir).
\end{enumerate}
\end{propriete}

\begin{popfigure}
\begin{center}
\includegraphics[width=\linewidth,height=11cm,keepaspectratio]{N01/cmr_physique.png}
\end{center}
\legende{Fond de carte physique du Cameroun pour la répartition des potentialités touristiques : balnéaire (littoral atlantique, au sud-ouest), safari (septentrion, au nord), culturel (hauts plateaux de l'Ouest) et écotourisme (forêt et montagne, au sud et à l'est). Les zones sont à colorier par l'élève. \\ Source : Wikimedia Commons, Urutseg, CC0.}
\end{popfigure}

\courssub{L'importance du tourisme au Cameroun}

\begin{propriete}[Importance économique]
\begin{itemize}
\item \textbf{Recettes en devises :} 3\textsuperscript{e} source de devises après le pétrole et le bois (ordre de grandeur : 300--350 milliards FCFA/an avant crise).
\item \textbf{Emplois :} Environ 300 000 emplois directs (hôtellerie, restauration, transport, guides, artisanat) et 600 000 indirects (agriculture vivrière pour hôtels, BTP, commerce).
\item \textbf{Investissements :} Construction d'hôtels, réhabilitation des routes d'accès (ex: route Kribi-Lolodorf, route du Nord), aéroports (Yaoundé-Nsimalen, Douala, Maroua-Salak, Garoua).
\item \textbf{Développement local :} Revenus pour les communautés riveraines des parcs (redevances, emplois d'éco-gardes, vente artisanat).
\end{itemize}
\end{propriete}

\begin{propriete}[Importance politique et sociale]
\begin{itemize}
\item \textbf{Intégration nationale :} Le tourisme relie les 10 régions. Il oblige à désenclaver le Nord (routes, aéroports), valorise le patrimoine de l'Ouest (chefferies), protège les forêts du Sud/Est. Il crée une \cle{solidarité territoriale}.
\item \textbf{Image du pays :} Vecteur de diplomatie publique (\og soft power \fg{}). Le slogan \og Afrique en miniature \fg{} positionne le Cameroun comme résumé du continent.
\item \textbf{Cohésion sociale :} Festivals (Ngoun, Ngondo, Medumba) rassemblent populations locales et diaspora. Échanges interrégionaux (artisanat de l'Ouest vendu à Kribi, poisson du Littoral consommé au Nord).
\item \textbf{Protection de l'environnement :} L'écotourisme finance la conservation (parcs, réserves) et sensibilise aux changements climatiques.
\end{itemize}
\end{propriete}

\courssub{Les problèmes du secteur touristique camerounais}

\begin{aretenir}[Diagnostic structurel]
Malgré un potentiel exceptionnel, le Cameroun stagne autour de 800 000--1 000 000 visiteurs/an (contre >10 millions pour le Maroc, ~2 millions pour le Kenya). Quatre blocs de freins :
\end{aretenir}

\begin{enumerate}
\item \textbf{Déficit de formation et de professionnalisation des acteurs :} Pénurie de guides certifiés, personnel hôtelier non qualifié (langues, gestion, service), restaurateurs peu formés aux normes d'hygiène internationales. Absence de grandes écoles hôtelières publiques de niveau international (hors instituts privés).
\item \textbf{Insuffisance et mauvaise répartition des infrastructures d'accueil et de transport :}
    \begin{itemize}
    \item Hébergement : ~1 000 hôtels classés (majorité 1--2 étoiles), concentrés à Douala (30\%) et Yaoundé (25\%). Pénurie d'écolodges et de camps de brousse de standing dans les parcs.
    \item Transports : État dégradé des routes nationales vers les sites (RN1 Nord, RN10 Est, pistes d'accès aux parcs). Saison des pluies = coupure physique (ex: Mayo-Rey, Dja). Transport aérien intérieur cher et irrégulier (Camair-Co). Absence de liaison ferroviaire touristique.
    \end{itemize}
\item \textbf{Insécurité et instabilité :} Crise anglophone (Nord-Ouest, Sud-Ouest) depuis 2016 : fermetures d'hôtels, impossibilité d'accéder aux sites (Mont Cameroun, chutes Ekom-Nkam, chefferies Bafut). Insécurité Extrême-Nord (Boko Haram) : parc de Waza inaccessible par moments. Banditisme sur axes routiers.
\item \textbf{Faible promotion et gouvernance :} Budget promotion MINTOUR faible. Pas de représentation touristique permanente sur les grands marchés émetteurs (France, Belgique, Chine, USA, Nigeria). Procédures de visa lourdes (bien que e-visa lancé en 2023). Fiscalité lourde (TVA 19,25\%, taxes aéroportuaires).
\end{enumerate}

\courssub{Perspectives et valorisation du tourisme pour l'intégration nationale}

\begin{methode}[Leviers stratégiques pour l'intégration par le tourisme]
\begin{enumerate}
\item \textbf{Aménagement du territoire :} Programmer les infrastructures (routes, électricité, eau, numérique) en partant des sites touristiques (schéma directeur d'aménagement touristique). Ex: bitumage axe Ngaoundéré--Garoua--Maroua--Waza ; boucle de l'Ouest (Bafoussam--Foumban--Bamenda).
\item \textbf{Formation et certification :} Créer un \emph{Institut National du Tourisme et de l'Hôtellerie} (INTH) public, décentralisé (antennes Nord, Ouest, Sud). Certifier les guides locaux (connaissance biodiversité, langues, secourisme).
\item \textbf{Produits touristiques intégrés :} Concevoir des \og circuits intégrés \fg{} (ex: \emph{Yaoundé -- Dja -- Kribi -- Douala -- Ouest -- Bénoué -- Waza -- Maroua}) vus par des tour-opérateurs nationaux, liant balnéaire, culture, safari, écotourisme.
\item \textbf{Marketing ciblé et numérique :} Plateforme unique \og Visit Cameroon \fg{} multilingue. Participation aux salons (ITB Berlin, FITUR Madrid, AKWAABA Lagos). Cibler la diaspora (tourisme de racines) et la sous-région CEMAC (tourisme d'affaires, week-ends).
\item \textbf{Gouvernance et sécurité :} Mise en œuvre effective du \emph{Plan de Sécurité Touristique} (police touristique, numéros d'urgence, assurance voyage). Simplification visa (visa à l'arrivée pour pays CEMAC/CEEAC, e-visa généralisé).
\end{enumerate}
\end{methode}

\begin{savaistu}[Tourisme d'affaires : un levier méconnu]
Yaoundé et Douala accueillent des conférences internationales (UA, CEMAC, OAPI, OMS, sommets climatiques). Le tourisme d'affaires (MICE : Meetings, Incentives, Conferences, Exhibitions) génère un panier moyen journalier 3 à 4 fois supérieur au tourisme de loisirs. Le Cameroun dispose d'atouts : stabilité relative des capitales, hôtels 4-5 étoiles, centres de conférences. Son développement structure la demande hôtelière haut de gamme qui profite in fine au tourisme de loisirs (mutualisation des infrastructures).
\end{savaistu}

\courssub{Activité d'intégration : Diagnostic du tourisme camerounais}

\courssub{Objectifs de l'activité}

Cette activité d'intégration te place dans une situation professionnelle réaliste : celle d'un jeune géographe chargé de produire un diagnostic du tourisme camerounais. À la fin de l'activité, tu dois être capable de :

\begin{itemize}
\item \cle{observer} et \cle{décrire} des photographies de sites touristiques selon la méthode vue en classe (description, localisation, interprétation) ;
\item \cle{localiser} et \cle{cartographier} les principaux sites touristiques du Cameroun sur une carte muette, avec une légende organisée ;
\item \cle{construire} un diagramme en bâtons à partir d'un tableau statistique et \cle{interpréter} les tendances qu'il révèle ;
\item \cle{inventorier} et \cle{classer} les problèmes du secteur touristique à partir d'un extrait de rapport ;
\item \cle{argumenter} des propositions de développement et les relier à la famille de situations : \emph{l'intégration nationale}.
\end{itemize}

\courssub{Situation}

\textbf{Communiqué du lycée.} Dans le cadre de la Semaine nationale du tourisme, le Ministère du Tourisme et des Loisirs lance un concours inter-lycées intitulé \og Faire du Cameroun une destination touristique de premier plan \fg{}. Ton lycée de Yaoundé y participe. Chaque équipe doit remettre une \textbf{affiche de synthèse} comportant : une carte légendée des sites touristiques, un diagramme statistique commenté, et trois propositions concrètes de développement du secteur. Tu es désigné(e) rapporteur de ton groupe. Voici le dossier documentaire remis par le jury.

\textbf{Document 1 — Carte muette du Cameroun.} Une carte vierge du territoire national, avec les contours du pays, la position des dix chefs-lieux de région et la flèche du nord. Elle est à compléter.

\textbf{Document 2 — Deux photographies.}
\begin{itemize}
\item \emph{Photo A} : la plage de Kribi (région du Sud). On observe une longue plage de sable clair bordée de cocotiers, quelques pirogues de pêcheurs, des hôtels de moyenne gamme en arrière-plan, et des chalutiers au large. La route d'accès, récemment bitumée, est visible.
\item \emph{Photo B} : le parc national de Waza (région de l'Extrême-Nord). On observe une savane arborée en saison sèche, un point d'eau autour duquel se concentrent éléphants, girafes et antilopes ; au premier plan, une piste latéritique et un véhicule tout-terrain de touristes.
\end{itemize}

\textbf{Document 3 — Tableau statistique.} Arrivées de touristes internationaux et recettes touristiques du Cameroun (ordres de grandeur, sources : OMT et INS, valeurs arrondies).

\begin{center}
\begin{tabularx}{\linewidth}{|l|X|X|X|X|X|}
\hline
\rowcolor{popblueL} \textbf{Année} & \textbf{2015} & \textbf{2017} & \textbf{2019} & \textbf{2021} & \textbf{2023} \\
\hline Arrivées (milliers) & 800 & 900 & 1\,000 & 400 & 750 \\
\hline Recettes (milliards FCFA) & 250 & 300 & 350 & 130 & 280 \\
\hline
\end{tabularx}
\end{center}

\textbf{Document 4 — Extrait d'un rapport sur les difficultés du secteur.} \og Le Cameroun, surnommé \emph{l'Afrique en miniature}, n'accueille qu'environ un million de visiteurs par an, loin derrière le Maroc (plus de 10 millions) ou le Kenya (environ 2 millions). Le secteur souffre d'un déficit de formation des acteurs (guides, hôteliers, restaurateurs), d'une insuffisance des structures d'accueil (moins de 1 000 hôtels classés, concentrés à Douala et Yaoundé), du mauvais état de nombreuses voies de communication vers les sites (pistes du Nord et de l'Est impraticables en saison des pluies), de l'insécurité dans certaines zones (Extrême-Nord, Nord-Ouest, Sud-Ouest) et d'une faible promotion du pays à l'étranger. \fg{}

\courssub{Consignes}

\begin{enumerate}
\item \textbf{Observer et décrire.} Décris la photo A puis la photo B en trois étapes (ce que je vois / où cela se situe / ce que cela montre comme type de tourisme). Identifie pour chacune la \cle{potentialité} touristique correspondante du programme.
\item \textbf{Localiser et cartographier.} Sur la carte muette, place et légende au moins six éléments : Kribi et Limbé (tourisme balnéaire), les parcs de Waza et de la Bénoué (safari), les chefferies des hauts plateaux de l'Ouest (tourisme culturel), la réserve forestière du Dja (écotourisme). Ajoute un titre, la flèche du nord et une légende organisée par type de tourisme.
\item \textbf{Construire et interpréter.} À partir du document 3, construis un diagramme en bâtons des arrivées de touristes (2015--2023). Calcule le taux de variation des arrivées entre 2019 et 2021, puis entre 2021 et 2023. Interprète : que s'est-il passé ? La reprise est-elle complète ?
\item \textbf{Inventorier et classer.} À partir du document 4, classe les problèmes du secteur en trois catégories : problèmes humains (formation), problèmes matériels (infrastructures), problèmes de contexte (sécurité, promotion).
\item \textbf{Proposer et argumenter.} Formule trois propositions concrètes pour développer le tourisme, en précisant pour chacune le problème qu'elle résout.
\item \textbf{Relier à l'intégration nationale.} Rédige un court paragraphe (5 à 8 lignes) répondant à la question : en quoi le développement du tourisme renforce-t-il l'intégration nationale du Cameroun ?
\end{enumerate}

\courssub{Ressources à mobiliser}

\begin{aretenir}[Boîte à outils de l'activité]
\begin{itemize}
\item \textbf{Méthode d'analyse de photographie} : décrire (les faits visibles, sans interpréter) $\rightarrow$ localiser (site, région) $\rightarrow$ interpréter (quel type de tourisme ? quelle potentialité ? quels atouts et limites visibles ?).
\item \textbf{Règles du croquis cartographique} : titre en haut, légende dans un cartouche, flèche du nord, échelle indicative, peu de couleurs (une par type de tourisme), noms lisibles.
\item \textbf{Taux de variation} (en \%) :
\[ T = \frac{V_{\text{arrivée}} - V_{\text{départ}}}{V_{\text{départ}}} \times 100 \]
\item \textbf{Les quatre potentialités} : attrait de la mer (Kribi, Limbé) ; safari dans les parcs du septentrion (Waza, Bénoué, Bouba Ndjida) ; dynamisme culturel des hauts plateaux (chefferies de l'Ouest, festivals) ; écotourisme dans le plateau forestier (Dja, Campo-Ma'an).
\item \textbf{Les problèmes types} : déficit de formation des acteurs ; insuffisance des structures d'accueil et des voies de communication ; insécurité ; faible promotion.
\end{itemize}
\end{aretenir}

\courssub{Démarche et corrigé commenté}

\begin{exoresolu}[Corrigé commenté de l'affiche de synthèse]
\textbf{Étape 1 — Analyse des photographies.}

\emph{Photo A (Kribi).} Je vois une plage de sable clair, des cocotiers, des hôtels et des pirogues. Il s'agit du littoral atlantique de la région du Sud. Cette photo illustre le \cle{tourisme balnéaire}, fondé sur l'attrait de la mer. Atouts : plage accessible par une route bitumée, cadre naturel préservé. Limite visible : hôtellerie de moyenne gamme seulement, pêche artisanale et tourisme parfois en concurrence pour l'espace.

\emph{Photo B (Waza).} Je vois une savane sèche, un point d'eau fréquenté par des éléphants et des girafes, une piste et un véhicule de touristes. Il s'agit du parc national de Waza (Extrême-Nord). Cette photo illustre le \cle{safari} (tourisme de vision de la faune), potentialité majeure du septentrion camerounais. Limite : piste non bitumée, ce qui renvoie au problème des voies de communication ; l'insécurité régionale a longtemps détourné les visiteurs.

\tcblower

\textbf{Étape 2 — Carte légendée (croquis).}

\begin{popfigure}
\begin{center}
\includegraphics[width=\linewidth,height=11cm,keepaspectratio]{N01/cmr_parcs.jpg}
\end{center}
\legende{Fond de carte des parcs nationaux du Cameroun sur relief (titre en espagnol, noms lisibles) pour le croquis des grands sites touristiques. Le croquis attendu distingue : le tourisme balnéaire (littoral : Kribi, Limbé), le safari (parcs du septentrion : Waza, Bénoué, Bouba Ndjida), le tourisme culturel (hauts plateaux de l'Ouest) et l'écotourisme (plateau forestier : Dja, Campo-Ma'an, Korup). \\ Source : Wikimedia Commons, Teogomez, CC0 (légende des parcs en espagnol pour le titre ; noms des parcs lisibles).}
\end{popfigure}

\textbf{Étape 3 — Diagramme en bâtons et calculs.}

\begin{popfigure}
\begin{tikzpicture}
\begin{axis}[
  ybar, bar width=14pt,
  width=0.85\linewidth, height=5.5cm,
  ymin=0, ymax=1200,
  ylabel={Arrivées (milliers de touristes)},
  symbolic x coords={2015,2017,2019,2021,2023},
  xtick=data,
  nodes near coords,
  every node near coord/.append style={font=\footnotesize},
  axis lines=left,
  ymajorgrids=true, grid style={popline!40},
]
\addplot[fill=popblue,draw=popdark] coordinates {(2015,800) (2017,900) (2019,1000) (2021,400) (2023,750)};
\end{axis}
\end{tikzpicture}
\legende{Évolution des arrivées de touristes internationaux au Cameroun, 2015--2023 (ordres de grandeur arrondis, sources OMT/INS).}
\end{popfigure}

Calculs des taux de variation :
\begin{align*}
T_{2019\rightarrow 2021} &= \frac{400 - 1\,000}{1\,000}\times 100 = -60\ \% \\
T_{2021\rightarrow 2023} &= \frac{750 - 400}{400}\times 100 = +87{,}5\ \%
\end{align*}

\emph{Interprétation.} Les arrivées progressent régulièrement de 2015 à 2019 (de 800 000 à 1 million), puis s'effondrent de 60 \% en 2021, sous l'effet conjugué de la crise sanitaire mondiale (Covid-19) et de l'insécurité. La reprise de 2021 à 2023 est forte (+87,5 \%) mais \textbf{incomplète} : en 2023, le niveau de 2019 (1 million) n'est pas encore retrouvé. Les recettes suivent la même courbe (350 milliards en 2019, 130 en 2021, 280 en 2023), ce qui confirme le lien direct entre fréquentation et retombées économiques.

\textbf{Étape 4 — Classement des problèmes.}

\begin{center}
\begin{tabularx}{\linewidth}{|l|X|}
\hline
\rowcolor{popblueL} \textbf{Catégorie} & \textbf{Problèmes identifiés dans le rapport} \\
\hline Humains & Déficit de formation des guides, hôteliers et restaurateurs \\
\hline Matériels & Insuffisance des structures d'accueil (hôtels classés concentrés à Douala et Yaoundé) ; voies de communication dégradées vers les sites \\
\hline Contexte & Insécurité (Extrême-Nord, Nord-Ouest, Sud-Ouest) ; faible promotion du pays à l'étranger \\
\hline
\end{tabularx}
\end{center}

\textbf{Étape 5 — Trois propositions argumentées.}
\begin{enumerate}
\item Créer des \textbf{filières de formation professionnelle} au tourisme et à l'hôtellerie dans les lycées techniques et les universités (problème résolu : déficit de formation des acteurs).
\item \textbf{Bitumer les routes d'accès} aux grands sites (axe vers Waza, boucle des chefferies de l'Ouest) et développer des écolodges dans les zones forestières (problème résolu : insuffisance des infrastructures).
\item Mener une \textbf{campagne internationale de promotion} (\og l'Afrique en miniature \fg{}) et sécuriser les zones touristiques (problème résolu : image du pays et insécurité).
\end{enumerate}

\textbf{Étape 6 — Lien avec l'intégration nationale.} Le tourisme oblige à relier les régions entre elles : construire une route vers Waza ou vers le Dja, c'est désenclaver le septentrion et le plateau forestier et les raccorder au reste du pays. Il fait circuler les Camerounais et les étrangers entre littoral, hauts plateaux et savanes, créant des échanges économiques (emplois, marchés locaux, artisanat) et culturels (festivals, chefferies) qui tissent le sentiment d'appartenance à une même nation. En valorisant toutes les régions — chacune possède une potentialité — le tourisme réduit les déséquilibres territoriaux et fait du Cameroun un espace national intégré et solidaire.
\end{exoresolu}

\begin{attention}[Pièges fréquents dans ce type de dossier]
\begin{itemize}
\item Ne confonds pas \textbf{description} et \textbf{interprétation} d'une photographie : décris d'abord ce qui est visible, puis seulement explique.
\item Un taux de variation négatif n'est pas une erreur : $-60\ \%$ signifie une baisse de 60 \%.
\item Sur un croquis, une légende sans titre ni nord est une légende incomplète : des points sont perdus au Baccalauréat.
\item Évite les propositions vagues (\og développer le tourisme \fg{}) : chaque proposition doit cibler un problème précis du dossier.
\end{itemize}
\end{attention}

\courssub{Bilan de l'activité}

Cette activité t'a permis de mobiliser ensemble tous les savoir-faire de la leçon. L'\textbf{observation} des photographies a montré que les potentialités camerounaises sont réelles mais inégalement aménagées. La \textbf{cartographie} a révélé la répartition complémentaire des sites sur tout le territoire : mer au sud-ouest, safari au nord, culture à l'ouest, forêt au sud-est. La \textbf{construction du diagramme} et le calcul des taux de variation ont mis en évidence la fragilité du secteur face aux crises (sanitaire, sécuritaire) et sa capacité de rebond. Enfin, le classement des problèmes et l'élaboration de propositions t'ont entraîné(e) à raisonner comme un décideur. La conclusion est claire : avec environ un million de visiteurs par an, le Cameroun est loin de son potentiel d'\og Afrique en miniature \fg{} ; mais parce que le tourisme relie les régions, crée des emplois et valorise toutes les identités du pays, son développement est un puissant levier d'\cle{intégration nationale} — c'est précisément le sens de la famille de situations de ce module.

\courssub{Bilan de la leçon}

\begin{aretenir}[Synthèse structurée pour le Bac]
\begin{itemize}
\item \textbf{Potentialités :} 4 piliers (Mer, Safari, Culture, Éco) couvrant les 10 régions $\rightarrow$ atout d'intégration.
\item \textbf{Importance :} Devises (3\textsuperscript{e} rang), Emplois (directs/indirects), Intégration (désenclavement, cohésion), Image (soft power).
\item \textbf{Problèmes :} Formation (humain), Infrastructures (matériel), Insécurité/Promotion (contexte).
\item \textbf{Perspectives :} Aménagement territoire + Formation certifiante + Produits intégrés + Marketing numérique + Gouvernance/Sécurité.
\item \textbf{Mots-clés Bac :} \cle{Intégration nationale}, \cle{Afrique en miniature}, \cle{Potentialités}, \cle{Désenclavement}, \cle{Tourisme d'affaires}, \cle{Écotourisme}, \cle{Safari}, \cle{Chefferies}.
\end{itemize}
\end{aretenir}

\begin{exercice}[Entraînement Bac - Sujet type]
\begin{enumerate}
\item Présentez les quatre grandes potentialités touristiques du Cameroun en les localisant précisément (régions, sites emblématiques).
\item À l'aide d'exemples, montrez que le tourisme est un facteur d'intégration nationale.
\item Expliquez pourquoi le secteur peine à décoller malgré ses atouts. Proposez deux mesures prioritaires.
\end{enumerate}
\end{exercice}

\begin{corrige}[Exercice n°1]
\begin{enumerate}
\item \textbf{Quatre potentialités :} 1) Balnéaire : Kribi, Limbé, Campo (Sud, Sud-Ouest). 2) Safari : Waza (Extrême-Nord), Bénoué (Nord), Bouba Ndjida (Adamaoua). 3) Culturel : Chefferies Bafut, Bandjoun, Foumban (Ouest, Nord-Ouest). 4) Écotourisme : Réserve du Dja (Sud/Est), Campo-Ma'an (Sud), Mont Cameroun (Sud-Ouest).
\item \textbf{Facteur d'intégration :} Infrastructures routières reliant Nord/Sud (ex: bitumage RN1, RN10) ; flux de touristes et Camerounais entre régions (artisanat Ouest vendu à Kribi, poisson Littoral au Nord) ; valorisation patrimoine commun (festivals, parcs) créant fierté nationale ; emplois locaux réduisant exode rural et tensions.
\item \textbf{Freins :} Insécurité (NO, SO, EN), routes dégradées, hôtels concentrés Douala/Yaoundé, formation insuffisante, promotion faible. \textbf{Mesures prioritaires :} 1) Sécurisation zones touristiques + bitumage axes d'accès (Waza, Dja, Ouest). 2) Création Institut National Tourisme/Hôtellerie + campagne marketing ciblée diaspora/CEMAC.
\end{enumerate}
\end{corrige}

\newpage

\coursec{Méthodes et exercices résolus}

\courssub{Méthode 1 : Observer, localiser et décrire un site touristique}

\begin{methode}[Décrire et localiser un site touristique]
\begin{enumerate}
\item \textbf{Situer précisément} : nommer le site, la commune ou le département, la région, la grande zone naturelle (littoral, septentrion, hauts plateaux de l'Ouest, plateau forestier du Sud). Exemple : le parc national de Waza, département du Logone-et-Chari, région de l'Extrême-Nord.
\item \textbf{Identifier le type de tourisme} : balnéaire (mer), safari (parcs), culturel (chefferies, festivals), écotourisme (forêt), tourisme d'affaires (villes, congrès). Un site peut combiner plusieurs types.
\item \textbf{Décrire les attraits} : éléments naturels (faune, plages, reliefs, climat) et éléments culturels (musées, monuments, artisanat, danses).
\item \textbf{Décrire les équipements} : accès (route, aéroport, piste), hébergement (hôtels, campements), encadrement (guides, agences).
\item \textbf{Conclure} par un jugement nuancé : atouts réels mais contraintes (enclavement, saisonnalité, insécurité).
\end{enumerate}
\textbf{Réflexe} : une bonne description va toujours du général (localisation) au particulier (attraits, équipements), et cite au moins un fait concret (un animal, un monument, une ville-relais).
\end{methode}

\begin{exoresolu}[Décrire un site touristique : le parc national de Waza]
\textbf{Énoncé.} À partir de vos connaissances, décrivez le site touristique du parc national de Waza en suivant la démarche : localisation, type de tourisme, attraits, équipements, limites. (6 points)
\tcblower
\textbf{Solution détaillée.}

\emph{1. Localisation.} Le parc national de Waza se situe dans le département du Logone-et-Chari, région de l'Extrême-Nord du Cameroun, à environ 120 km au nord de Maroua, non loin de la frontière avec le Tchad et le Nigeria. Il appartient à la zone soudano-sahélienne du septentrion camerounais.

\emph{2. Type de tourisme.} C'est le site emblématique du \cle{tourisme de safari} au Cameroun, essentiellement le safari photographique et d'observation de la faune (la chasse sportive, elle, est pratiquée hors du parc, dans les zones d'intérêt cynégétique aménagées du septentrion).

\emph{3. Attraits.} Créé en 1934 comme réserve, érigé en parc national en 1968, il couvre environ \num{170000} hectares (ordre de grandeur) de savane. Il abrite une faune riche : éléphants, lions, girafes, antilopes, buffles, autruches et de nombreux oiseaux. Il est classé réserve de biosphère par l'UNESCO depuis 1979.

\emph{4. Équipements.} Le parc est accessible par la route nationale n°1 (Yaoundé--Maroua, puis route vers Waza) et par l'aéroport de Maroua-Salak. On y trouve un campement touristique (campement de Waza) et des pistes d'observation avec guides.

\emph{5. Limites.} La fréquentation est freinée par l'éloignement des grands centres émetteurs, la chaleur et la saisonnalité (visite surtout en saison sèche, de novembre à mai), et surtout par l'insécurité liée aux exactions de Boko Haram dans la région de l'Extrême-Nord depuis 2014.

\emph{Conclusion.} Waza illustre un site à fort potentiel mais sous-valorisé : la richesse faunique y est exceptionnelle, tandis que l'accessibilité et la sécurité restent des contraintes majeures.
\end{exoresolu}

\courssub{Méthode 2 : Légender une carte des potentialités touristiques}

\begin{methode}[Construire la légende d'une carte touristique du Cameroun]
\begin{enumerate}
\item \textbf{Choisir un titre précis} : « Les grandes potentialités touristiques du Cameroun », jamais un titre vague comme « Le tourisme ».
\item \textbf{Hiérarchiser la légende} en 3 ou 4 rubriques maximum :
\begin{itemize}
\item fonds de carte colorés ou figurés de surface pour les grandes zones (littoral balnéaire, parcs du septentrion, hauts plateaux culturels, plateau forestier d'écotourisme) ;
\item symboles ponctuels pour les sites précis (parcs, plages, chefferies, chutes) ;
\item traits et points nommés pour les villes-relais et les axes d'accès.
\end{itemize}
\item \textbf{Respecter les conventions} : une couleur = une seule signification ; figurés simples et lisibles ; pas plus de 6 ou 7 entrées de légende.
\item \textbf{Placer la légende} dans un cadre, en bas à droite ou à gauche, avec titre « Légende ».
\item \textbf{Ajouter} l'orientation (flèche du nord) et, si possible, une échelle indicative.
\end{enumerate}
\textbf{Réflexe} : au Bac, la légende vaut souvent à elle seule plusieurs points ; chaque figuré utilisé sur la carte doit figurer dans la légende, et réciproquement.
\end{methode}

\begin{exoresolu}[Légender une carte des potentialités touristiques]
\textbf{Énoncé.} On vous donne un fond de carte du Cameroun. Proposez une légende complète et ordonnée pour une carte intitulée « Les grandes potentialités touristiques du Cameroun », puis indiquez où placer quatre sites : Kribi, Waza, Foumban, la réserve du Dja. (8 points)
\tcblower
\textbf{Solution détaillée.}

\emph{Légende proposée (hiérarchisée) :}
\begin{enumerate}
\item \textbf{Zones de tourisme balnéaire} (figuré de surface, bleu) : côte atlantique de Limbé à Kribi et Campo.
\item \textbf{Parcs et réserves de safari} (figuré de surface, vert) : parcs nationaux de Waza, Bénoué, Bouba-Ndjida, Faro.
\item \textbf{Zones de tourisme culturel} (figuré de surface, orange) : hauts plateaux de l'Ouest et du Nord-Ouest (chefferies bamiléké, palais du sultan bamoun à Foumban, musées et artisanat).
\item \textbf{Zones d'écotourisme forestier} (figuré de surface, vert foncé hachuré) : plateau forestier du Sud et de l'Est (réserve du Dja, parc de Campo-Ma'an, parc de Lobéké).
\item \textbf{Sites majeurs} (symbole ponctuel : étoile) nommés.
\item \textbf{Villes-relais et accès} : points noirs nommés (Douala, Yaoundé, Maroua) et traits pour les axes routiers principaux.
\end{enumerate}

\emph{Placement des sites :}
\begin{itemize}
\item \textbf{Kribi} : sur la côte, au sud de Douala, dans la région du Sud (zone 1) ; célèbre pour ses plages et les chutes de la Lobé qui se jettent directement dans l'océan.
\item \textbf{Waza} : dans l'Extrême-Nord, au nord de Maroua (zone 2).
\item \textbf{Foumban} : dans la région de l'Ouest, au nord-est de Bafoussam (zone 3) ; palais du sultanat bamoun, musée, artisanat d'art.
\item \textbf{Réserve du Dja} : dans le Sud-Est, au-delà de Sangmélima en direction de l'Est (zone 4) ; classée au patrimoine mondial de l'UNESCO depuis 1987.
\end{itemize}

\emph{Conclusion.} La carte montre que les potentialités couvrent tout le territoire national, avec une spécialisation par grande zone naturelle : la mer au sud-ouest, le safari au nord, la culture à l'ouest, la forêt au sud-est. On peut compléter par les grands festivals qui ponctuent l'année : le ngondo des Sawa sur le littoral (Douala) ou le nyem-nyem des peuples de l'Adamaoua (Ngaoundéré), qui montrent que le dynamisme culturel dépasse les seuls hauts plateaux.
\end{exoresolu}

\courssub{Méthode 3 : Construire un croquis de flux touristiques}

\begin{methode}[Construire un croquis de flux]
\begin{enumerate}
\item \textbf{Tracer le fond} : contour simplifié du Cameroun (polygone) ; placer les principales villes (Douala, Yaoundé, Maroua, Kribi, Bafoussam, Garoua) par des points proportionnels ou nommés.
\item \textbf{Identifier les pôles} : pôles émetteurs (aéroports internationaux de Douala et Yaoundé, tourisme d'affaires) et pôles récepteurs (parcs, plages, sites culturels).
\item \textbf{Dessiner les flux} par des flèches dont l'\textbf{épaisseur est proportionnelle} à l'importance du flux : flèche large pour les flux internationaux vers Douala/Yaoundé, flèches fines pour les flux intérieurs vers les sites.
\item \textbf{Distinguer} par des couleurs les flux internationaux et les flux internes (tourisme domestique).
\item \textbf{Légender} chaque type de flèche, ajouter titre, nord et échelle indicative.
\end{enumerate}
\textbf{Réflexe} : un croquis n'est pas une carte topographique ; on simplifie, on schématise, mais chaque élément doit être juste et légendé.
\end{methode}

\begin{exoresolu}[Croquis des flux touristiques vers le septentrion]
\textbf{Énoncé.} Construisez un croquis intitulé « Les flux touristiques vers les parcs du septentrion camerounais ». Vous ferez apparaître : les portes d'entrée internationales, les villes-relais du Nord, les parcs de Waza et de la Bénoué, et la hiérarchie des flux. (8 points)
\tcblower
\textbf{Solution détaillée.}

\emph{Étape 1 — Fond de carte.} Contour simplifié du Cameroun ; points nommés : Douala et Yaoundé (aéroports internationaux, portes d'entrée), Ngaoundéré (gare terminus du Transcamerounais), Garoua et Maroua (villes-relais du septentrion, dotées d'aéroports).

\emph{Étape 2 — Sites.} Étoiles vertes : parc de Waza (au nord de Maroua) et parc de la Bénoué (à l'est de Garoua).

\emph{Étape 3 — Flux.}
\begin{itemize}
\item Flèche large (flux international) : de l'extérieur vers Douala et Yaoundé ;
\item flèches moyennes : Douala/Yaoundé $\to$ Ngaoundéré (par le rail, Transcamerounais) et $\to$ Maroua/Garoua (par la route nationale n°1 ou par air) ;
\item flèches fines : Maroua $\to$ Waza ; Garoua $\to$ Bénoué.
\end{itemize}

\emph{Étape 4 — Légende.} 1) portes d'entrée internationales ; 2) villes-relais ; 3) parcs nationaux ; 4) flux internationaux ; 5) flux intérieurs principaux ; 6) flux intérieurs secondaires ; plus flèche du nord et échelle indicative.

\emph{Lecture du croquis.} Le croquis met en évidence la \textbf{désarticulation spatiale} du tourisme camerounais : les parcs du septentrion, principaux attraits safari, sont à plus de \num{1000} km des portes d'entrée internationales, ce qui allonge et renchérit les circuits. D'où l'importance des aéroports de Maroua et Garoua et de la route nationale n°1 comme axes de désenclavement touristique.

\emph{Conclusion.} La hiérarchie des flèches traduit une réalité : le flux s'amenuise au fur et à mesure qu'on s'éloigne des métropoles du sud, signe d'un potentiel septentrional encore mal connecté.
\end{exoresolu}

\courssub{Méthode 4 : Inventorier et classer des documents touristiques}

\begin{methode}[Inventorier et classer des documents]
\begin{enumerate}
\item \textbf{Inventorier} : pour chaque document (carte, photographie, tableau statistique), noter la nature, le titre ou le contenu, la source et la date.
\item \textbf{Dégager les critères de classement} avant de classer : type de tourisme (balnéaire, safari, culturel, écotourisme, affaires), zone géographique, ou nature de l'information (potentialités / importance / problèmes).
\item \textbf{Classer} chaque document dans une rubrique et \textbf{justifier} le classement en une phrase (« cette photo de lions dans la savane relève du tourisme de safari car... »).
\item \textbf{Croiser} les documents : un tableau statistique (arrivées, recettes) sert à évaluer l'importance économique ; une photo d'une piste dégradée illustre les problèmes.
\item \textbf{Conclure} en dégageant ce que l'ensemble des documents révèle (richesse du potentiel, faiblesse de la valorisation).
\end{enumerate}
\textbf{Réflexe} : un classement sans critère explicite est sanctionné ; annonce toujours ton critère avant de ranger les documents.
\end{methode}

\begin{exoresolu}[Classer cinq documents sur le tourisme camerounais]
\textbf{Énoncé.} On dispose de cinq documents : (A) photo des chutes de la Lobé à Kribi ; (B) tableau des arrivées de touristes internationaux (environ 1 million de visiteurs par an à la fin des années 2010, ordre de grandeur) ; (C) photo d'une danse traditionnelle bamiléké à Bandjoun ; (D) photo d'une piste défoncée menant à un parc ; (E) photo d'éléphants dans le parc de la Bénoué. Classez ces documents selon les trois thèmes du programme (potentialités, importance, problèmes) en justifiant chaque choix. (7 points)
\tcblower
\textbf{Solution détaillée.}

\emph{Critère retenu} : le thème du programme auquel chaque document apporte une information.

\begin{center}
\begin{tabularx}{\linewidth}{|l|X|X|}
\hline
\rowcolor{popblueL} \textbf{Document} & \textbf{Thème} & \textbf{Justification} \\
\hline
A (chutes de la Lobé) & Potentialités & Site naturel remarquable du littoral : attrait de la mer et des cascades, tourisme balnéaire. \\
\hline
B (arrivées de touristes) & Importance & Le tableau statistique mesure le poids économique du secteur (devises, emplois induits). \\
\hline
C (danse bamiléké) & Potentialités & Dynamisme culturel des hauts plateaux de l'Ouest : tourisme culturel. \\
\hline
D (piste défoncée) & Problèmes & Illustration de l'insuffisance des voies de communication, frein majeur à l'accès aux sites. \\
\hline
E (éléphants, Bénoué) & Potentialités & Faune de savane : attrait du safari dans les parcs du septentrion. \\
\hline
\end{tabularx}
\end{center}

\emph{Croisement.} Le document B chiffre une activité encore modeste (environ 1 million de visiteurs, dont une forte part de tourisme d'affaires et de la diaspora, ordre de grandeur à manier avec prudence), alors que A, C et E montrent un potentiel considérable : l'écart entre les deux s'explique en partie par le problème illustré en D.

\emph{Conclusion.} L'inventaire croisé révèle le paradoxe du tourisme camerounais : un pays « Afrique en miniature » au potentiel exceptionnel, mais dont la valorisation reste limitée par les infrastructures.
\end{exoresolu}

\courssub{Méthode 5 : Exploiter un tableau statistique touristique}

\begin{methode}[Lire et calculer dans un tableau statistique]
\begin{enumerate}
\item \textbf{Lire le titre, l'unité et la source} avant tout calcul (effectifs en milliers ? en \% ? recettes en milliards de FCFA ?).
\item \textbf{Dégager les faits saillants} : valeur maximale, minimale, évolution générale (hausse, baisse, stagnation).
\item \textbf{Calculer les indicateurs utiles} :
\begin{itemize}
\item part en pourcentage : $\text{part} = \dfrac{\text{valeur de la catégorie}}{\text{total}} \times 100$ ;
\item taux de variation : $t = \dfrac{V_f - V_i}{V_i} \times 100$ (en \%) ;
\item évolution absolue : $V_f - V_i$.
\end{itemize}
\item \textbf{Interpréter géographiquement} chaque résultat (relier aux potentialités, aux infrastructures, au contexte de sécurité).
\item \textbf{Rédiger} en phrases : un chiffre non commenté ne rapporte pas de point.
\end{enumerate}
\textbf{Réflexe} : vérifie qu'une somme de parts fait bien 100 \% (aux arrondis près) et qu'un taux négatif est bien présenté comme une baisse.
\end{methode}

\begin{exoresolu}[Calculs sur les arrivées de touristes]
\textbf{Énoncé.} Les arrivées de visiteurs internationaux au Cameroun sont estimées à environ \num{650000} en 2010 et à environ \num{1000000} en 2019 (ordres de grandeur). On estime par ailleurs qu'en 2019, environ 40 \% de ces visiteurs relèvent du tourisme d'affaires et des visites familiales (diaspora), et 60 \% du tourisme d'agrément et de découverte. 1) Calculez l'évolution absolue et le taux de variation entre 2010 et 2019. 2) Calculez le nombre de visiteurs de chaque catégorie en 2019. 3) Interprétez. (8 points)
\tcblower
\textbf{Solution détaillée.}

\emph{1. Évolution entre 2010 et 2019.}
\begin{align*}
\text{Évolution absolue} &= \num{1000000} - \num{650000} = \num{350000} \text{ visiteurs.}\\
t &= \frac{\num{1000000} - \num{650000}}{\num{650000}} \times 100 = \frac{\num{350000}}{\num{650000}} \times 100 \approx 53,8\,\%.
\end{align*}
Les arrivées ont donc augmenté d'environ 54 \% en neuf ans, soit une hausse moyenne d'environ 5 \% par an (ordre de grandeur).

\emph{2. Répartition de 2019.}
\begin{align*}
\text{Affaires et diaspora} &= \num{1000000} \times \frac{40}{100} = \num{400000} \text{ visiteurs.}\\
\text{Agrément et découverte} &= \num{1000000} \times \frac{60}{100} = \num{600000} \text{ visiteurs.}
\end{align*}
Vérification : $\num{400000} + \num{600000} = \num{1000000}$, le total est respecté.

\emph{3. Interprétation.} La progression traduit l'amélioration de l'accueil (hôtellerie à Douala et Yaoundé, grands événements, stabilité relative) et l'attractivité du pays. Mais la part élevée du \cle{tourisme d'affaires} et des visites de la diaspora (environ 40 \%) montre que le tourisme de loisirs proprement dit (safaris, plages, culture) reste sous-développé par rapport au potentiel : le Cameroun attire davantage par ses affaires que par ses sites, contrairement au Kenya ou à la Tanzanie.

\emph{Conclusion.} Les calculs confirment un secteur en croissance mais déséquilibré : la valorisation touristique des potentialités naturelles et culturelles demeure un chantier majeur pour l'économie nationale.
\end{exoresolu}

\courssub{Méthode 6 : Rédiger un commentaire de document (texte ou image) sur le tourisme}

\begin{methode}[Commenter un document en géographie]
\begin{enumerate}
\item \textbf{Présenter le document} (2-3 lignes) : nature, auteur ou source, date, thème.
\item \textbf{Décrire ou dégager les idées principales} : pour une image, décrire du premier plan vers l'arrière-plan ; pour un texte, relever 2 ou 3 idées-forces avec des mots du document.
\item \textbf{Analyser} : mobiliser tes connaissances du cours pour expliquer ce que le document montre (potentialités, importance économique, problèmes). C'est l'étape qui rapporte le plus de points.
\item \textbf{Discuter les limites} : ce que le document ne dit pas (date ancienne, échelle locale, point de vue de l'auteur).
\item \textbf{Conclure} en répondant explicitement à la problématique posée par le sujet.
\end{enumerate}
\textbf{Réflexe} : respecte l'équilibre description / analyse ; un commentaire qui ne fait que paraphraser le document plafonne à la moitié des points.
\end{methode}

\begin{exoresolu}[Commentaire d'un texte sur le tourisme au Cameroun]
\textbf{Énoncé.} « Le Cameroun, c'est l'Afrique en miniature : plages, savanes peuplées de grands mammifères, forêts équatoriales, montagnes, cultures multiples. Pourtant, le pays accueille moins de touristes de loisirs que certains de ses voisins, faute d'infrastructures d'accueil et de communication suffisantes. » Commentez ce texte en montrant qu'il illustre le contraste entre potentialités et valorisation du tourisme camerounais. (10 points)
\tcblower
\textbf{Solution détaillée.}

\emph{Présentation.} Il s'agit d'un extrait de texte de vulgarisation géographique portant sur le tourisme au Cameroun ; il oppose la richesse du potentiel (« Afrique en miniature ») à la faiblesse de la fréquentation.

\emph{Idées principales.} Deux idées-forces : (1) une diversité exceptionnelle d'attraits naturels et culturels ; (2) un déficit de valorisation imputé aux infrastructures.

\emph{Analyse.} La première idée correspond exactement aux potentialités étudiées : l'attrait de la mer (plages de Kribi et Limbé), le safari dans les parcs du septentrion (Waza, Bénoué, Bouba-Ndjida), le dynamisme culturel des hauts plateaux (chefferies bamiléké, palais du sultan bamoun à Foumban), auquel s'ajoutent les grands festivals d'autres régions (ngondo des Sawa à Douala, nyem-nyem dans l'Adamaoua), et l'écotourisme dans le plateau forestier (réserve du Dja, Campo-Ma'an). La seconde idée renvoie aux problèmes du secteur : insuffisance des structures d'accueil (parc hôtelier limité hors de Douala et Yaoundé), mauvais état de nombreuses voies de communication (pistes vers les parcs), déficit de formation des acteurs (guides, hôteliers), insécurité dans l'Extrême-Nord et dans les régions du Nord-Ouest et du Sud-Ouest. On peut ajouter que le tourisme conserve pourtant une importance réelle : devises, emplois (hôtellerie, artisanat, transport), rayonnement politique et culturel du pays, et contribution à l'intégration nationale en faisant circuler les Camerounais dans leur propre pays.

\emph{Limites.} Le texte ne chiffre ni les arrivées ni les recettes, et ne mentionne pas le poids du tourisme d'affaires ; il ne dit rien non plus des politiques récentes de promotion du secteur (aménagement et classement des hôtels, actions de promotion touristique).

\emph{Conclusion.} Le texte illustre bien le paradoxe central de la leçon : le Cameroun possède des potentialités « énormes » mais une valorisation limitée ; lever les contraintes d'infrastructures, de formation et de sécurité est la condition pour transformer ce potentiel en levier de développement et d'intégration nationale.
\end{exoresolu}

\begin{attention}[Les 5 erreurs qui coûtent des points au Bac]
\begin{enumerate}
\item \textbf{Confondre les types de tourisme et leurs zones.} Le safari se situe dans les parcs du \emph{septentrion} (Waza, Bénoué), l'écotourisme dans le \emph{plateau forestier} du Sud (Dja, Campo-Ma'an) : inverser les deux est une erreur de localisation sanctionnée. Autre confusion fréquente : la chasse sportive ne se pratique pas \emph{dans} les parcs nationaux, mais dans les zones d'intérêt cynégétique voisines.
\item \textbf{Légende incomplète ou incohérente.} Utiliser un figuré sur la carte sans le mettre dans la légende (ou l'inverse), oublier le titre, le nord ou l'épaisseur proportionnelle des flèches de flux : chaque oubli coûte un point.
\item \textbf{Erreur de calcul de pourcentage.} Diviser par la valeur finale au lieu de la valeur initiale dans un taux de variation, ou oublier de multiplier par 100. Écris toujours la formule avant d'appliquer : $t = \dfrac{V_f - V_i}{V_i} \times 100$.
\item \textbf{Paraphraser le document au lieu de l'analyser.} Un commentaire qui répète le texte sans mobiliser le cours (potentialités, importance, problèmes) plafonne à la moitié des points : chaque idée du document doit être expliquée par une connaissance.
\item \textbf{Présenter des ordres de grandeur comme des chiffres exacts.} Les statistiques touristiques camerounaises sont incertaines : écris « environ 1 million de visiteurs (ordre de grandeur) » plutôt qu'un chiffre précis inventé, et cite prudemment les faits sûrs (Dja classé au patrimoine mondial de l'UNESCO en 1987, Waza réserve de biosphère depuis 1979).
\end{enumerate}
\end{attention}

\newpage

\coursec{Exercices}

\courssub{Je vérifie mes connaissances}

\begin{exercice}[QCM : les fondamentaux du tourisme camerounais]
Pour chaque question, entoure la (ou les) bonne(s) réponse(s).

\textbf{1.} Un site touristique est :
\begin{enumerate}[label=\alph*)]
\item tout espace fréquenté par des habitants ;
\item un lieu remarquable qui attire des visiteurs en raison de ses qualités naturelles ou culturelles ;
\item uniquement un monument classé par l'UNESCO ;
\item une zone réservée aux touristes étrangers.
\end{enumerate}

\textbf{2.} Le tourisme d'affaires désigne :
\begin{enumerate}[label=\alph*)]
\item les déplacements touristiques motivés par des activités professionnelles (congrès, séminaires, missions) ;
\item l'achat de souvenirs par les touristes ;
\item les safaris organisés par des entreprises privées ;
\item le tourisme pratiqué uniquement dans les capitales régionales.
\end{enumerate}

\textbf{3.} Les safaris au Cameroun se pratiquent principalement :
\begin{enumerate}[label=\alph*)]
\item dans les parcs nationaux du septentrion (Waza, Bénoué, Bouba Ndjida) ;
\item sur les plages de Kribi ;
\item dans les plantations de la CDC ;
\item sur le plateau forestier du Sud.
\end{enumerate}

\textbf{4.} Le dynamisme culturel des hauts plateaux de l'Ouest se manifeste surtout par :
\begin{enumerate}[label=\alph*)]
\item les chefferies, les palais, les musées et les danses traditionnelles ;
\item les ports de pêche industrielle ;
\item les barrages hydroélectriques ;
\item les mines de fer.
\end{enumerate}
\end{exercice}

\begin{exercice}[Vrai ou faux ? Justifie ta réponse]
Pour chaque affirmation, indique si elle est vraie ou fausse, puis justifie en une ou deux phrases.

\textbf{1.} « Le Cameroun est souvent surnommé \emph{l'Afrique en miniature}, ce qui constitue un atout majeur pour son tourisme. »

\textbf{2.} « Les chutes de la Lobé sont un site d'écotourisme situé dans les hauts plateaux de l'Ouest. »

\textbf{3.} « Le déficit de formation des acteurs du secteur est l'un des problèmes du tourisme camerounais. »

\textbf{4.} « Le tourisme ne rapporte aucune devise au Cameroun. »

\textbf{5.} « L'insécurité dans certaines régions (Extrême-Nord, Nord-Ouest, Sud-Ouest) n'a aucune incidence sur la fréquentation touristique. »
\end{exercice}

\begin{exercice}[Définitions et exemples]
\begin{enumerate}
\item Définis en deux ou trois lignes chacune des notions suivantes : \cle{site touristique} ; \cle{tourisme d'affaires} ; \cle{écotourisme}.
\item Pour chacune de ces trois notions, cite \textbf{deux exemples camerounais} précis (nom du site ou de la ville, région).
\item Classe les sites suivants selon le type de tourisme dominant (balnéaire, safari, culturel, écotourisme, affaires) : plage de Kribi ; parc national de Waza ; chefferie de Bafoussam ; réserve forestière du Dja ; palais des congrès de Yaoundé ; chutes de la Lobé.
\end{enumerate}
\end{exercice}

\begin{exercice}[Calculs directs sur les arrivées touristiques]
Le tableau ci-dessous donne les arrivées de touristes internationaux au Cameroun (ordres de grandeur, source : ministère du Tourisme et des Loisirs / OMT).

\begin{center}
\begin{tabularx}{\linewidth}{|l|X|X|X|X|}\hline
\rowcolor{popblueL} \textbf{Année} & \textbf{2015} & \textbf{2017} & \textbf{2019} & \textbf{2021} \\ \hline
Arrivées (milliers) & 900 & 1 000 & 1 100 & 500 \\ \hline
\end{tabularx}
\end{center}

\begin{enumerate}
\item Calcule le taux de variation des arrivées entre 2015 et 2019 (en \%, arrondi au dixième).
\item Calcule le taux de variation entre 2019 et 2021. À quel événement mondial attribues-tu cette évolution ?
\item Sachant qu'un touriste international dépense en moyenne environ 500 000 FCFA par séjour, estime les recettes touristiques de l'année 2019 en milliards de FCFA.
\item Quelle part représentent les arrivées de 2021 par rapport à celles de 2019 (en \%) ?
\end{enumerate}
\end{exercice}

\courssub{Je m'entraîne}

\begin{exercice}[Localisation sur carte muette]
À partir d'une carte muette du Cameroun (contour, limites régionales, villes principales) :

\begin{enumerate}
\item Localise et légende \textbf{cinq sites touristiques} : parc national de Waza, chutes de la Lobé, mont Cameroun, chefferie de Foumban, plage de Kribi.
\item Pour chaque site, précise dans la légende le \textbf{type de tourisme} associé (safari, balnéaire, culturel, écotourisme, randonnée).
\item Trace en pointillés les deux principaux axes d'accès : la route Yaoundé--Douala--Kribi et l'axe Yaoundé--Ngaoundéré--Maroua.
\item N'oublie pas les éléments obligatoires de toute carte : titre, légende, orientation (flèche du nord), échelle indicative.
\end{enumerate}
\end{exercice}

\begin{exercice}[Construction d'un diagramme en bâtons]
On te donne les arrivées de touristes internationaux au Cameroun (ordres de grandeur, en milliers) :

\begin{center}
\begin{tabularx}{\linewidth}{|l|X|X|X|X|X|X|}\hline
\rowcolor{popblueL} \textbf{Année} & 2014 & 2015 & 2016 & 2017 & 2018 & 2019 \\ \hline
Arrivées (milliers) & 850 & 900 & 950 & 1 000 & 1 050 & 1 100 \\ \hline
\end{tabularx}
\end{center}

\begin{enumerate}
\item Construis un diagramme en bâtons représentant l'évolution des arrivées (axe horizontal : années ; axe vertical : arrivées en milliers, échelle : 1 cm pour 100 000 touristes).
\item Calcule l'augmentation totale des arrivées entre 2014 et 2019, en valeur absolue et en \%.
\item Rédige un commentaire de \textbf{dix lignes} : présente le document, décris la tendance générale, compare les valeurs extrêmes et propose une explication de cette évolution (atouts du pays, promotion du \emph{Cameroun, Afrique en miniature}).
\end{enumerate}
\end{exercice}

\begin{exercice}[Analyse de photographie : les chutes de la Lobé]
On te présente une photographie des chutes de la Lobé (région du Sud, près de Kribi) : une série de cascades larges et peu hautes se jetant directement dans l'océan Atlantique, bordées d'une plage de sable et d'une forêt dense ; quelques pirogues de pêcheurs et un groupe de visiteurs sont visibles au premier plan.

\begin{enumerate}
\item Décris méthodiquement la photographie (du premier plan vers l'arrière-plan, ou du centre vers les bords).
\item Identifie le ou les type(s) de tourisme que ce site peut accueillir. Justifie à partir d'indices visibles sur la photo.
\item Dégage \textbf{deux atouts} et \textbf{deux contraintes} du site (accessibilité, équipements, encadrement).
\item Propose deux mesures concrètes pour mieux valoriser ce site.
\end{enumerate}
\end{exercice}

\begin{exercice}[Classement et croquis des potentialités]
\begin{enumerate}
\item Recopie et complète le tableau suivant en classant les potentialités touristiques du Cameroun selon les quatre grands ensembles du programme :

\begin{center}
\begin{tabularx}{\linewidth}{|l|X|X|}\hline
\rowcolor{popblueL} \textbf{Ensemble géographique} & \textbf{Type de tourisme dominant} & \textbf{Exemples de sites} \\ \hline
Littoral atlantique & & \\ \hline
Septentrion (Nord, Extrême-Nord, Adamaoua) & & \\ \hline
Hauts plateaux de l'Ouest et du Nord-Ouest & & \\ \hline
Plateau forestier du Sud et de l'Est & & \\ \hline
\end{tabularx}
\end{center}

\item Construis un croquis simple du Cameroun (contour simplifié en polygone) figurant ces quatre ensembles par quatre couleurs différentes, avec titre, légende et flèche du nord.
\item Explique en cinq lignes pourquoi on dit que le Cameroun est \emph{l'Afrique en miniature} et en quoi cela constitue un atout touristique.
\end{enumerate}
\end{exercice}

\courssub{Je me prépare au Bac}

\begin{exercice}[Dissertation type Bac : potentialités, importance et problèmes]
\textbf{Sujet :} « Le tourisme au Cameroun : potentialités, importance et problèmes. »

Tu rédigeras un devoir complet et structuré.

\begin{enumerate}
\item \textbf{Introduction} (environ 10 lignes) : amène le sujet à partir d'une situation concrète (par exemple l'arrivée d'un groupe de touristes à l'aéroport de Douala), définis les termes clés (\emph{site touristique}, \emph{tourisme}), présente la problématique et annonce le plan.
\item \textbf{Développement} en trois parties :
\begin{enumerate}
\item Les potentialités touristiques énormes et variées du Cameroun (mer, safari, culture, écotourisme) ;
\item L'importance économique, politique et sociale du tourisme (devises, emplois, rayonnement international, intégration nationale) ;
\item Les problèmes qui freinent le secteur (déficit de formation, insuffisance des structures d'accueil et des voies de communication, insécurité).
\end{enumerate}
\item \textbf{Conclusion} (environ 8 lignes) : bilan nuancé et ouverture sur les perspectives (tourisme durable, promotion internationale).
\end{enumerate}

\emph{Critères d'évaluation : exactitude des connaissances, richesse des exemples camerounais précis, cohérence du plan, qualité de la langue.}
\end{exercice}

\begin{exercice}[Étude de cas : le parc national de Waza et les difficultés du septentrion]
\textbf{Document :} le parc national de Waza (Extrême-Nord), créé en 1934 et classé réserve de biosphère par l'UNESCO en 1979, couvre environ 170 000 hectares. Il abrite éléphants, lions, girafes, antilopes et une avifaune exceptionnelle. Dans les années 2000, il accueillait plusieurs dizaines de milliers de visiteurs par an. Depuis le milieu des années 2010, les exactions du groupe Boko Haram dans l'Extrême-Nord ont provoqué l'effondrement de la fréquentation : les arrivées ont chuté de plus de 80 \% (ordre de grandeur), de nombreux hôtels de Maroua ont fermé et des centaines d'emplois (guides, chauffeurs, artisans, restaurateurs) ont disparu.

\begin{enumerate}
\item Présente le document (nature, sujet, chiffres clés) en trois ou quatre lignes.
\item Identifie le type de tourisme pratiqué à Waza et montre en quoi ce parc illustre les potentialités du septentrion camerounais.
\item À partir du document, dégage \textbf{trois conséquences} de l'insécurité sur l'économie touristique locale (fréquentation, hébergement, emploi).
\item Calcule le nombre approximatif de visiteurs restants si le parc recevait 30 000 visiteurs par an avant la crise et que la fréquentation a chuté de 80 \%.
\item L'insécurité est-elle le \emph{seul} obstacle au développement du tourisme dans le septentrion ? Mobilise au moins deux autres problèmes du secteur (voies de communication, structures d'accueil, formation des acteurs).
\item Propose deux mesures que l'État et les acteurs locaux pourraient prendre pour relancer le tourisme à Waza.
\end{enumerate}
\end{exercice}

\begin{exercice}[Sujet de réflexion argumentée : l'insécurité, principal obstacle ?]
\textbf{Sujet :} « L'insécurité est-elle le principal obstacle au développement du tourisme au Cameroun ? »

Tu rédigeras un texte argumenté de 25 à 30 lignes, structuré selon la démarche suivante :

\begin{enumerate}
\item \textbf{Introduction} : pose le problème à partir d'un fait d'actualité (baisse de la fréquentation de l'Extrême-Nord, crises dans le Nord-Ouest et le Sud-Ouest) et formule clairement la problématique.
\item \textbf{Premier moment} : montre que l'insécurité est un obstacle réel et grave (annulation de séjours, fermeture d'hôtels, image du pays dégradée à l'étranger).
\item \textbf{Deuxième moment} : nuance en mobilisant les \emph{autres} problèmes du secteur : déficit de formation des acteurs, insuffisance des structures d'accueil (hôtels, restaurants classés), mauvais état ou absence de voies de communication vers certains sites (routes vers l'Est, pistes des parcs), faiblesse de la promotion internationale.
\item \textbf{Conclusion} : prends position de manière argumentée et propose une hiérarchie des solutions (sécurisation, investissements dans les infrastructures, formation professionnelle, marketing touristique).
\end{enumerate}

\emph{Attendus : au moins quatre exemples camerounais précis (sites, villes, régions), un vocabulaire géographique correct, une argumentation en deux temps (thèse / nuance).}
\end{exercice}

\newpage

\coursec{Corrigés des exercices}

\coursec{Le tourisme au Cameroun}

\courssub{Je vérifie mes connaissances}

\begin{corrige}[Exercice 1 — QCM : les fondamentaux du tourisme camerounais]
\textbf{1.} \textbf{b)} Un site touristique est un lieu remarquable qui attire des visiteurs en raison de ses qualités naturelles (paysage, faune, flore) ou culturelles (patrimoine historique, traditions, architecture). Les autres propositions sont fausses : un espace fréquenté par des habitants n'est pas nécessairement touristique (a) ; un site n'a pas besoin d'être classé UNESCO (c) ; il n'est pas réservé aux étrangers (d).

\textbf{2.} \textbf{a)} Le tourisme d'affaires (ou tourisme professionnel) concerne les déplacements motivés par des activités professionnelles : congrès, séminaires, conférences, missions commerciales, salons. Les autres propositions correspondent au tourisme de loisir (achats, safaris) ou sont trop restrictives (capitales régionales).

\textbf{3.} \textbf{a)} Les safaris (observation de la faune sauvage en véhicule) se pratiquent principalement dans les parcs nationaux du septentrion : Waza (Extrême-Nord), Bénoué (Nord), Bouba Ndjida (Nord), Faro (Nord). Kribi est balnéaire (b), la CDC est agro-industrielle (c), le plateau forestier est voué à l'écotourisme (d).

\textbf{4.} \textbf{a)} Les hauts plateaux de l'Ouest (Bamileke, Bamoun) sont réputés pour leur dynamisme culturel : chefferies traditionnelles (Bafoussam, Bandjoun, Foumban), palais des sultans/rois, musées (musée des civilisations de Dschang, musée du palais de Foumban), danses, masques, architecture traditionnelle. Les ports (b), barrages (c) et mines (d) relèvent d'autres secteurs.
\end{corrige}

\begin{corrige}[Exercice 2 — Vrai ou faux ? Justifie ta réponse]
\textbf{1.} \textbf{Vrai.} Le Cameroun est surnommé « l'Afrique en miniature » car il concentre sur son territoire la quasi-totalité des climats, reliefs, végétations et cultures du continent (océan, montagne, savane, forêt, désert, diversité ethnique). Cette diversité est un atout majeur car elle permet d'offrir une palette touristique très large sur un seul territoire.

\textbf{2.} \textbf{Faux.} Les chutes de la Lobé se situent sur le \textbf{littoral atlantique}, dans la région du Sud, près de Kribi (océan). Elles ne sont pas dans les hauts plateaux de l'Ouest (régions de l'Ouest et du Nord-Ouest).

\textbf{3.} \textbf{Vrai.} Le programme cite explicitement le « déficit de formation des acteurs » parmi les problèmes du secteur. Il manque de personnel qualifié (guides, hôteliers, managers, agents de promotion) formés aux standards internationaux.

\textbf{4.} \textbf{Faux.} Le tourisme est une source importante de devises pour le Cameroun (recettes estimées à plusieurs centaines de milliards de FCFA avant la crise Covid-19). Il contribue à la balance des paiements.

\textbf{5.} \textbf{Faux.} L'insécurité (exactions de Boko Haram dans l'Extrême-Nord, crise anglophone au Nord-Ouest et Sud-Ouest) a un impact dramatique : annulation de séjours, fermeture d'hôtels et de camps de brousse, perte d'emplois, dégradation de l'image du pays à l'étranger, baisse drastique des arrivées internationales dans ces régions.
\end{corrige}

\begin{corrige}[Exercice 3 — Définitions et exemples]
\textbf{1. Définitions}
\begin{itemize}
\item \cle{Site touristique} : Lieu géographique (naturel ou aménagé) présentant un intérêt particulier (esthétique, scientifique, historique, culturel, récréatif) qui attire des visiteurs extérieurs au lieu de résidence habituelle pour une durée limitée.
\item \cle{Tourisme d'affaires} : Forme de tourisme dont la motivation principale est professionnelle : participation à des congrès, séminaires, conférences, salons, missions commerciales, formations. Il génère des retombées économiques fortes (hôtellerie haut de gamme, restauration, transport, location de salles).
\item \cle{Écotourisme} : Tourisme responsable dans les zones naturelles qui contribue à la conservation de l'environnement et au bien-être des populations locales. Il implique l'observation de la nature (faune, flore, paysages), l'éducation à l'environnement et une empreinte écologique minimale.
\end{itemize}

\textbf{2. Exemples camerounais (deux par notion)}
\begin{itemize}
\item \textbf{Site touristique} : Mont Cameroun (région du Sud-Ouest) ; Chutes de la Lobé (région du Sud).
\item \textbf{Tourisme d'affaires} : Palais des congrès de Yaoundé (région du Centre) ; Hôtel Sawa / Akwa Palace à Douala pour salons professionnels (région du Littoral).
\item \textbf{Écotourisme} : Réserve de biosphère du Dja (régions du Sud et de l'Est) ; Parc national de Korup (région du Sud-Ouest) ; Sanctuaire des gorilles de Mengamé (région du Sud).
\end{itemize}

\textbf{3. Classement des sites}
\begin{center}
\begin{tabularx}{\linewidth}{|l|X|}\hline
\rowcolor{popblueL} \textbf{Site} & \textbf{Type de tourisme dominant} \\ \hline
Plage de Kribi & Balnéaire \\ \hline
Parc national de Waza & Safari (faune sauvage) \\ \hline
Chefferie de Bafoussam & Culturel \\ \hline
Réserve forestière du Dja & Écotourisme \\ \hline
Palais des congrès de Yaoundé & Tourisme d'affaires \\ \hline
Chutes de la Lobé & Écotourisme / Balnéaire (site mixte) \\ \hline
\end{tabularx}
\end{center}
\end{corrige}

\begin{corrige}[Exercice 4 — Calculs directs sur les arrivées touristiques]
\textbf{Données :} 2015 = 900 000 ; 2017 = 1 000 000 ; 2019 = 1 100 000 ; 2021 = 500 000.

\textbf{1. Taux de variation 2015 $\to$ 2019 :}
\[
TV = \frac{1\,100 - 900}{900} \times 100 = \frac{200}{900} \times 100 \approx 22,2\%
\]
Les arrivées ont augmenté de \textbf{22,2 \%} entre 2015 et 2019.

\textbf{2. Taux de variation 2019 $\to$ 2021 :}
\[
TV = \frac{500 - 1\,100}{1\,100} \times 100 = \frac{-600}{1\,100} \times 100 \approx -54,5\%
\]
Les arrivées ont chuté de \textbf{54,5 \%}. Cette évolution s'explique par la \textbf{pandémie de Covid-19} (fermeture des frontières, restrictions de voyage, peur sanitaire) qui a paralysé le tourisme mondial en 2020-2021.

\textbf{3. Recettes touristiques 2019 :}
\[
1\,100\,000 \text{ touristes} \times 500\,000 \text{ FCFA} = 550\,000\,000\,000 \text{ FCFA} = \textbf{550 milliards de FCFA}.
\]

\textbf{4. Part des arrivées 2021 par rapport à 2019 :}
\[
\frac{500}{1\,100} \times 100 \approx 45,5\%
\]
En 2021, le Cameroun n'a récupéré que \textbf{45,5 \%} de son niveau de fréquentation de 2019.
\end{corrige}

\courssub{Je m'entraîne}

\begin{corrige}[Exercice 5 — Localisation sur carte muette]
\textbf{Correction attendue (description de la carte à rendre) :}

\begin{popfigure}
\begin{center}
\includegraphics[width=\linewidth,height=9.5cm,keepaspectratio]{N01/cmr_map.jpg}
\end{center}
\legende{Fond de carte du Cameroun (régions, villes, routes, voie ferrée) servant de modèle : sur ce fond, l'élève localise cinq sites touristiques types (par exemple Waza, Kribi, Foumban, le mont Cameroun, Dja) et les axes d'accès principaux, déjà visibles en rouge (routes) sur la carte. \\ Source : Wikimedia Commons, CIA, domaine public.}
\end{popfigure}

\textbf{Points obligatoires vérifiés :}
\begin{itemize}
\item Titre précis.
\item Légende complète (symboles villes, sites, axes).
\item Orientation (flèche Nord).
\item Échelle indicative (ex: 1 cm $\approx$ 200 km).
\item 5 sites localisés correctement : Waza (Extrême-Nord), Lobé/Kribi (Sud), Mont Cameroun (Sud-Ouest), Foumban (Ouest).
\item Types de tourisme indiqués dans la légende.
\item Axes Yaoundé-Douala-Kribi et Yaoundé-Ngaoundéré-Maroua en pointillés.
\end{itemize}
\end{corrige}

\begin{corrige}[Exercice 6 — Construction d'un diagramme en bâtons]
\textbf{1. Diagramme en bâtons (pgfplots)}

\begin{popfigure}
\begin{tikzpicture}
\begin{axis}[
    ybar,
    bar width=0.6cm,
    width=\linewidth,
    height=8cm,
    xlabel={Année},
    ylabel={Arrivées (en milliers)},
    xtick={2014,2015,2016,2017,2018,2019},
    xticklabels={2014,2015,2016,2017,2018,2019},
    ymin=0, ymax=1200,
    ymajorgrids=true,
    grid style={popmuted},
    axis lines=left,
    enlarge x limits=0.15,
    title={Évolution des arrivées touristiques internationales au Cameroun (2014-2019)},
    title style={font=\bfseries\small},
    nodes near coords,
    nodes near coords align={vertical},
    every node near coord/.append style={font=\tiny, popdark}
]
\addplot[popblue, fill=popblueL] coordinates {
    (2014,850) (2015,900) (2016,950) (2017,1000) (2018,1050) (2019,1100)
};
\end{axis}
\end{tikzpicture}
\legende{Diagramme en bâtons de l'évolution des arrivées touristiques (en milliers) de 2014 à 2019. Source : MINTOUL / OMT. Échelle : 1 cm pour 100 000 touristes sur l'axe vertical.}
\end{popfigure}

\textbf{2. Calculs}
\begin{itemize}
\item \textbf{Augmentation absolue :} $1\,100 - 850 = \textbf{250 000 touristes}$ (soit 250 milliers).
\item \textbf{Taux de variation :} $\frac{250}{850} \times 100 \approx \textbf{29,4 \%}$.
\end{itemize}

\textbf{3. Commentaire (environ 10 lignes)}
Le diagramme en bâtons représente l'évolution annuelle des arrivées de touristes internationaux au Cameroun entre 2014 et 2019, en milliers. La tendance générale est à une \textbf{hausse continue et régulière} : les arrivées passent de 850 000 en 2014 à 1 100 000 en 2019, soit une progression totale de 29,4 \%. L'augmentation est quasi linéaire (+50 000 par an en moyenne). Cette croissance s'explique par les atouts majeurs du pays (« Afrique en miniature ») : diversité des paysages (mer, montagne, savane, forêt), richesse culturelle (chefferies, danses), et faune exceptionnelle (parcs du Nord). Elle témoigne aussi des efforts de promotion (campagne « Cameroun, Afrique en miniature »), de la préparation d'événements internationaux (infrastructures pour la CAN 2021 organisée en 2022) et d'une relative stabilité politique avant les crises sécuritaires récentes.
\end{corrige}

\begin{corrige}[Exercice 7 — Analyse de photographie : les chutes de la Lobé]
\textbf{1. Description méthodique}
\begin{itemize}
\item \textbf{Premier plan :} Une plage de sable fin clair, quelques pirogues traditionnelles en bois coloré échouées ou à l'eau, un petit groupe de visiteurs (touristes) observant les chutes, vêtus de vêtements légers.
\item \textbf{Plan moyen :} La série de cascades de la Lobé : un rideau d'eau large (environ 100 m), peu haut (quelques mètres), se déversant directement dans l'océan Atlantique. L'eau est trouble (chargée en alluvions). La rive opposée est bordée d'une forêt dense, verte, tropicale.
\item \textbf{Arrière-plan :} L'horizon marin (océan Atlantique) sous un ciel clair ou légèrement nuageux. La végétation forestière semble descendre jusqu'au bord de l'eau.
\end{itemize}

\textbf{2. Types de tourisme identifiés}
\begin{itemize}
\item \textbf{Écotourisme / Tourisme nature :} Présence de la forêt dense, de l'eau douce se jetant dans la mer (phénomène rare), biodiversité visible (oiseaux, végétation). Site naturel préservé.
\item \textbf{Balnéaire :} Plage de sable, océan, baignade possible (prudence aux courants).
\item \textbf{Tourisme culturel / de découverte :} Pirogues de pêcheurs (activité traditionnelle), possibilité de visiter les villages Bagyeli (Pygmées) proches.
\end{itemize}
\textit{Justification :} La photo montre simultanément le milieu forestier, le littoral, l'eau douce/salée et la présence humaine (pêcheurs, visiteurs).

\textbf{3. Atouts et contraintes}
\begin{center}
\begin{tabularx}{\linewidth}{|l|X|}\hline
\rowcolor{popblueL} \textbf{Atouts} & \textbf{Contraintes} \\ \hline
1. Site unique au monde : chutes se jetant directement dans l'océan. & 1. Accessibilité difficile : piste depuis Kribi (environ 10-15 km), dégradée en saison des pluies. \\ \hline
2. Cadre naturel exceptionnel (forêt, plage, mer, chutes) propice à l'écotourisme haut de gamme. & 2. Infrastructures d'accueil insuffisantes : peu d'hébergements de standing sur site, absence de centre d'interprétation, gestion des déchets perfectible. \\ \hline
\end{tabularx}
\end{center}

\textbf{4. Mesures de valorisation}
\begin{enumerate}
\item Aménager et bitumer la piste d'accès depuis Kribi (RN17 ou piste locale) et créer un parking organisé avec navettes fluviales ou pédestres pour limiter l'érosion.
\item Construire un \textbf{centre d'interprétation de l'écosystème} (panneaux éducatifs, guides locaux formés Bagyeli/Bakoko) et développer des écolodges certifiés (normes environnementales) gérés en partenariat avec les communautés locales (revenus partagés).
\end{enumerate}
\end{corrige}

\begin{corrige}[Exercice 8 — Classement et croquis des potentialités]
\textbf{1. Tableau complété}

\begin{center}
\begin{tabularx}{\linewidth}{|l|X|X|}\hline
\rowcolor{popblueL} \textbf{Ensemble géographique} & \textbf{Type de tourisme dominant} & \textbf{Exemples de sites} \\ \hline
Littoral atlantique & Balnéaire, tourisme d'affaires, écotourisme côtier & Plages de Kribi, Limbé, Douala (plages, hôtels, palais des congrès), Chutes de la Lobé, Réserve de la Douala-Edéa \\ \hline
Septentrion (Nord, Extrême-Nord, Adamaoua) & Safari (faune), tourisme culturel (lamidats), randonnée (monts Mandara) & Parcs nationaux de Waza, Bénoué, Bouba Ndjida, Faro ; Lamidats de Rey Bouba, Maroua ; Monts Mandara (Rhumsiki) ; Sites d'art rupestre de Bidzar \\ \hline
Hauts plateaux de l'Ouest et du Nord-Ouest & Tourisme culturel (chefferies, palais, musées), randonnée, agrotourisme & Chefferies de Bafoussam, Bandjoun, Baham, Foumban (palais des sultans, musée) ; Dschang (musée des civilisations, lac) ; Monts Bamboutos ; Plantations de thé de Njombe \\ \hline
Plateau forestier du Sud et de l'Est & Écotourisme (forêt dense, faune, pygmées), tourisme fluvial & Réserve de biosphère du Dja (UNESCO), Parc national de Lobéké, Parc de Campo Ma'an, Chutes de la Lobé, Villages Baka/Bagyeli \\ \hline
\end{tabularx}
\end{center}

\textbf{2. Croquis schématique du Cameroun (quatre ensembles)}

\begin{popfigure}
\begin{center}
\includegraphics[width=\linewidth,height=11cm,keepaspectratio]{N01/cmr_physique.png}
\end{center}
\legende{Fond de carte physique du Cameroun sur lequel figurer les quatre ensembles porteurs de potentialités touristiques : le Littoral (plaine côtière, bas de la carte à gauche), le Septentrion (savanes et plaines du nord), les Hauts plateaux (Ouest et Nord-Ouest, zone de reliefs de plus de 1000 m) et le Plateau forestier du Sud. Le croquis attendu les distingue par quatre couleurs (Littoral : bleu ; Septentrion : orange ; Hauts plateaux : vert ; Plateau forestier : violet). \\ Source : Wikimedia Commons, Urutseg, CC0.}
\end{popfigure}

\textbf{3. Pourquoi « l'Afrique en miniature » ? (5 lignes)}
Le Cameroun concentre sur 475 442 km$^2$ les grands types de milieux africains : littoral atlantique (plages, mangroves), montagne volcanique active (Mont Cameroun, 4 095 m), savanes soudano-sahéliennes au Nord (faune de brousse), forêt dense humide guinéo-congolaise au Sud/Est (biodiversité, pygmées), et hauts plateaux tempérés à l'Ouest (cultures, chefferies). Cette \textbf{diversité biophysique et culturelle} unique permet d'observer en un seul pays l'essentiel des paysages et écosystèmes du continent, offrant une palette touristique complète (mer, safari, culture, nature, montagne) sans changer de destination.
\end{corrige}

\courssub{Je me prépare au Bac}

\begin{corrige}[Exercice 9 — Dissertation type Bac : potentialités, importance et problèmes]

\textbf{Barème indicatif (sur 20) :}
\begin{itemize}
\item Introduction (3 pts) : Accroche, définitions, problématique, plan.
\item Partie I - Potentialités (5 pts) : 4 ensembles bien décrits, exemples précis, vocabulaire.
\item Partie II - Importance (5 pts) : 3 dimensions (éco, politico, socio), chiffres/ordres de grandeur, lien intégration.
\item Partie III - Problèmes (5 pts) : 3 problèmes majeurs développés, exemples concrets, lien cause-effet.
\item Conclusion (2 pts) : Bilan, ouverture pertinente.
\item Qualité de la langue, structure, propreté (bonus/malus).
\end{itemize}

\textbf{Points de vigilance (pièges à éviter) :}
\begin{itemize}
\item Ne pas confondre « potentialités » (ce qui \emph{existe}) et « réalisations » (ce qui \emph{est exploité}).
\item Citer des \textbf{noms propres} (parcs, villes, chefferies, régions) pour chaque exemple. Éviter les généralités (« il y a des animaux au Nord » $\to$ « éléphants, lions à Waza »).
\item Lier explicitement le tourisme à l'\textbf{intégration nationale} (famille de situations) : brassage des populations, connaissance mutuelle des régions, fierté nationale, aménagement du territoire.
\item Ne pas oublier le \textbf{tourisme d'affaires} (Yaoundé, Douala) dans l'importance économique.
\item Chiffres : utiliser des ordres de grandeur (ex: « environ 1 million d'arrivées avant Covid », « plusieurs centaines de milliards FCFA »).
\end{itemize}

\textbf{Proposition de rédaction structurée}

\textbf{Introduction}
[Accroche] En décembre 2023, un vol charter en provenance d'Europe atterrit à l'aéroport international de Douala, déchargeant un groupe de touristes venus découvrir les « chutes de la Lobé » et le « parc de Waza » en deux semaines. Cette scène illustre la capacité du Cameroun à offrir une synthèse de l'Afrique.
[Définitions] Un \emph{site touristique} est un lieu remarquable attirant des visiteurs pour ses qualités naturelles ou culturelles. Le \emph{tourisme} désigne l'ensemble des activités des personnes voyageant hors de leur environnement habituel pour une durée inférieure à un an (loisir, affaires, etc.).
[Problématique] Malgré des potentialités exceptionnelles, le secteur peine à devenir un pilier majeur de l'économie. \textbf{Comment expliquer le décalage entre les atouts touristiques du Cameroun et la performance réelle de la filière ?}
[Annonce du plan] Nous analyserons d'abord les potentialités énormes et variées (I), puis l'importance multidimensionnelle du secteur (II), avant d'examiner les problèmes structurels qui le freinent (III).

\textbf{Développement}

\textbf{I. Des potentialités touristiques énormes et variées : « l'Afrique en miniature »}
Le Cameroun possède une diversité biophysique unique en Afrique, structurée en quatre grands ensembles.
\begin{enumerate}
\item \textbf{Le littoral atlantique (Sud-Ouest, Littoral, Sud) :} 400 km de côtes offrent le tourisme balnéaire (plages de Kribi, Limbé, sable noir volcanique ou blanc), le tourisme d'affaires (Douala, Yaoundé : hôtels 5*, palais des congrès) et l'écotourisme côtier (chutes de la Lobé se jetant dans l'océan, réserve de la Douala-Edéa, mangroves).
\item \textbf{Le Septentrion (Adamaoua, Nord, Extrême-Nord) :} C'est le domaine du \textbf{safari}. Les parcs nationaux (Waza, Bénoué, Bouba Ndjida, Faro) abritent la « Big Five » africaine (éléphant, lion, léopard, buffle, rhinocéros - ce dernier rare) et une avifaune exceptionnelle (plus de 300 espèces à Waza). Le tourisme culturel y complète l'offre : lamidats traditionnels (Rey Bouba, Maroua), architecture soudanaise, monts Mandara (Rhumsiki, Kapsiki) pour la randonnée.
\item \textbf{Les Hauts Plateaux de l'Ouest et du Nord-Ouest :} Cœur du \textbf{tourisme culturel}. Densité de chefferies traditionnelles de 1er et 2e degré (Bafoussam, Bandjoun, Baham, Foumban), palais des sultans Bamoun (Foumban, musée des arts et traditions), musées (Dschang), danses, masques, artisanat d'art (pipes, tissus ndop, bronze). Le relief (monts Bamboutos, lacs de cratère) et les plantations (thé de Njombe, café) offrent un cadre pour l'agrotourisme et la randonnée.
\item \textbf{Le Plateau forestier du Sud et de l'Est :} Réservé à l'\textbf{écotourisme} de haute valeur. Forêt dense du bassin du Congo : réserve de biosphère du Dja (patrimoine UNESCO), parcs de Lobéké, Campo Ma'an, Boumba Bek. Observation des grands singes (gorilles, chimpanzés), éléphants de forêt, bongos. Rencontre avec les populations autochtones Baka/Bagyeli (chasse, cueillette, pharmacopée). Chutes d'eau (Lobé, Mbi).
\end{enumerate}
Cette variété permet le \textbf{circuit combiné} (ex: Douala $\to$ Kribi $\to$ Dja $\to$ Yaoundé $\to$ Ouest $\to$ Nord), atout majeur pour les tours-opérateurs.

\textbf{II. Une importance économique, politique et sociale majeure}
Le tourisme est un levier de développement structurel.
\begin{enumerate}
\item \textbf{Économique :} Source de \textbf{devises} (recettes estimées \textasciitilde 550 milliards FCFA en 2019 avant Covid). Contribution au PIB (\textasciitilde 2-3 \% direct, plus en induit). Création d'\textbf{emplois} directs (hôtellerie, restauration, transport, guides, artisanat) et indirects (agriculture vivrière pour hôtels, BTP). Développement des PME locales.
\item \textbf{Politique et diplomatique :} Outil de \textbf{rayonnement international} (image positive, « brand Cameroon »). Support à la diplomatie culturelle (festivals Ngondo, Nyem-Nyem, Festi-Bamiléké). Facteur de \textbf{stabilité} par l'emploi des jeunes.
\item \textbf{Sociale et intégration nationale :} \textbf{Brassage des populations} : les Camerounais du Sud découvrent le Nord (safari) et inversement (plages, culture). Renforcement de l'unité nationale par la connaissance mutuelle des patrimoines. Valorisation des cultures locales (langues, danses, architecture) face à l'uniformisation. Aménagement du territoire : ouverture de pistes, électrification des zones touristiques (ex: Waza, Rhumsiki).
\end{enumerate}

\textbf{III. Des problèmes structurels qui freinent l'essor du secteur}
Malgré les atouts, la part du Cameroun dans le tourisme africain reste faible (\textless 2 \% des arrivées continentales).
\begin{enumerate}
\item \textbf{Déficit de formation et de professionnalisme :} Pénurie de guides certifiés, de managers d'hôtels, de chefs de cuisine aux normes internationales. Formation initiale (École hôtelière de Yaoundé, IFT de Douala) insuffisante et déconnectée du terrain. Turn-over élevé, faible maîtrise des langues (anglais, chinois, espagnol).
\item \textbf{Insuffisance des structures d'accueil et des infrastructures :} Parc hôtelier vieillissant, peu d'établissements classés 4-5* hors Douala/Yaoundé. Camps de brousse (Waza, Bénoué) souvent vétustes, sans électricité stable ni eau courante. Manque de centres de congrès régionaux. \textbf{Voies de communication} : routes d'accès aux sites dégradées (pistes vers parcs du Nord, route vers l'Est/Dja), aéroports régionaux (Maroua, Ngaoundéré, Bertoua) sous-équipés, vols intérieurs chers et irréguliers (Camair-Co).
\item \textbf{Insécurité et image dégradée :} Crise Boko Haram (Extrême-Nord) depuis 2014 : effondrement de Waza (-80\% fréquentation). Crise anglophone (Nord-Ouest, Sud-Ouest) depuis 2016 : fermetures hôtels à Bamenda, Buea, Limbé ; Mont Cameroun inaccessible par moments. Enlèvements, barrages, couvre-feu. Conséquence : \textbf{avis négatifs} des ministères des Affaires étrangères occidentaux (zone rouge/orange), annulation de contrats par tours-opérateurs européens.
\end{enumerate}
\textit{Autres freins :} Promotion internationale timide (budget MINTOUL faible), visa encore contraignant (bien que e-visa lancé), saisonnalité marquée (juillet-août, décembre), faible implication des communautés locales dans la gouvernance des parcs.

\textbf{Conclusion}
Bilan : Le Cameroun possède un « produit » touristique de classe mondiale (diversité, authenticité), mais sa « mise sur le marché » est défaillante. L'insécurité actuelle agit comme un révélateur brutal de la fragilité structurelle (infrastructures, formation, promotion).
Ouverture : La relance passe par une \textbf{stratégie intégrée} : 1) \textbf{Sécurisation prioritaire} des corridors touristiques (Nord, Ouest) ; 2) \textbf{Investissements massifs} dans les infrastructures (routes, aéroports, énergie, hôtels) via PPP ; 3) \textbf{Professionnalisation} (écoles hôtelières, certification guides, langues) ; 4) \textbf{Marketing ciblé} (niche écotourisme haut de gamme, diaspora, tourisme intra-africain ZLECAf) ; 5) \textbf{Tourisme durable} (revenus pour communautés, conservation biodiversité). Le tourisme peut devenir le « pétrole vert » du Cameroun, vecteur d'intégration et de prospérité partagée, à condition de lever les verrous structurels.
\end{corrige}

\begin{corrige}[Exercice 10 — Étude de cas : le parc national de Waza et les difficultés du septentrion]

\textbf{1. Présentation du document}
Ce document est un \textbf{texte de géographie / étude de cas} portant sur le \textbf{parc national de Waza} (Extrême-Nord, Cameroun). Il fournit des chiffres clés : création 1934, classement UNESCO 1979 (Réserve de biosphère), superficie 170 000 ha, fréquentation passée « plusieurs dizaines de milliers » (années 2000), chute de 80 \% depuis milieu des années 2010 due à Boko Haram, fermeture d'hôtels à Maroua, perte de centaines d'emplois.

\textbf{2. Type de tourisme et potentialités du septentrion}
Le type dominant est le \textbf{safari photographique} (observation de la faune sauvage en véhicule 4x4). Waza illustre parfaitement les potentialités du septentrion : écosystème soudano-sahélien (plaine d'inondation du Logone, savane arbustive), concentration exceptionnelle de mégafaune (éléphants, lions, girafes, kob, damalisques, avifaune paléarctique migratrice), accessibilité relative depuis Maroua (aéroport, hôtels). C'est la « vitrine » du tourisme de nature au Cameroun.

\textbf{3. Trois conséquences de l'insécurité sur l'économie touristique locale}
\begin{enumerate}
\item \textbf{Effondrement de la fréquentation :} Chute de plus de 80 \% des visiteurs (passage de $\sim$30 000 à $\sim$6 000), perte directe des droits d'entrée (recettes parcs) et des dépenses sur site (guides, pourboires).
\item \textbf{Fermeture des structures d'hébergement :} Hôtels et camps de brousse de Maroua et aux abords du parc (Campement de Waza, hôtels de Maroua) ont fermé faute de clients, entraînant perte d'actifs et dettes bancaires.
\item \textbf{Destruction massive d'emplois :} Disparition de centaines d'emplois directs (guides pisteurs, chauffeurs 4x4, cuisiniers, serveurs, artisans vendant souvenirs, restaurateurs) et indirects (agriculteurs fournisseurs, transporteurs). Précarisation des populations riveraines (perte revenus, braconnage de survie).
\end{enumerate}

\textbf{4. Calcul du nombre de visiteurs restants}
Avant crise : 30 000 visiteurs/an. Chute de 80 \%.
Visiteurs restants = $30\,000 \times (1 - 0,80) = 30\,000 \times 0,20 = \textbf{6 000 visiteurs/an}$ (ordre de grandeur).

\textbf{5. L'insécurité est-elle le seul obstacle ? Non.}
Deux autres problèmes majeurs (programme) :
\begin{itemize}
\item \textbf{Voies de communication :} La route Maroua-Waza (RN1 puis piste) est dégradée, impraticable en saison des pluies (juillet-octobre), isolant le parc 4 mois/an. L'aéroport de Maroua a peu de vols intérieurs réguliers et abordables. L'axe Yaoundé-Ngaoundéré-Maroua (route/rail) est long (2 jours) et fatiguant.
\item \textbf{Structures d'accueil insuffisantes / vétustes :} Le campement de Waza (gestion publique) est souvent en mauvais état (électricité par groupe électrogène capricieuse, plomberie défaillante, literie usée). Peu d'alternatives privées de standing. Manque de centres d'interprétation, de signalétique, de points d'eau d'observation aménagés.
\item \textbf{(Bonus) Déficit de formation :} Guides-pisteurs souvent formés « sur le tas », connaissances naturalistes limitées, faible maîtrise langues étrangères, pas de certification officielle renouvelée.
\end{itemize}

\textbf{6. Deux mesures pour relancer le tourisme à Waza}
\begin{enumerate}
\item \textbf{Sécurisation et reconquête du territoire :} Déploiement renforcé du BIR (Bataillon d'Intervention Rapide) et de la garde forestière sur les axes d'accès (Maroua-Waza-Kousseri) et à l'intérieur du parc (postes avancés). Opérations de « nettoyage » des abords. Mise en place d'un « corridor sécurisé » garanti par l'État pour les tours-opérateurs (convois escortés si nécessaire). Communication officielle sur la levée des zones rouges par les ambassades partenaires.
\item \textbf{Programme d'investissement et de professionnalisation « Waza Relance » :} Réhabilitation complète du campement de Waza (normes 3* écologiques : solaire, eau traitée, déchets) via PPP (Partenariat Public-Privé). Formation certifiante des guides (faune, flore, secourisme, anglais) par l'École de la Faune de Garoua. Création d'un fonds de garantie pour les tours-opérateurs (remboursement si annulation force majeure). Promotion ciblée « Waza Safe » auprès de la diaspora, des expatriés (Yaoundé/Douala) et du marché régional (Tchad, Nigeria, Afrique centrale).
\end{enumerate}
\end{corrige}

\begin{corrige}[Exercice 11 — Sujet de réflexion argumentée : l'insécurité, principal obstacle ?]

\textbf{Barème indicatif :}
\begin{itemize}
\item Introduction (3 pts) : Actualité, problématique claire.
\item Moment 1 - Thèse : Insécurité obstacle grave (4 pts) : exemples précis (Waza, Mont Cameroun, Bamenda), mécanismes (image, annulations, emplois).
\item Moment 2 - Nuance : Autres obstacles structurels (4 pts) : formation, infrastructures, promotion, visa. Hiérarchisation.
\item Conclusion (3 pts) : Position argumentée, hiérarchie des solutions.
\item Vocabulaire géo, exemples précis, structure (4 pts).
\end{itemize}

\textbf{Proposition de rédaction (25-30 lignes)}

\textbf{Introduction}
Depuis 2014, l'Extrême-Nord subit les exactions de Boko Haram, vidant le parc de Waza de ses touristes (-80 \%). Depuis 2016, la crise anglophone ferme les hôtels de Bamenda, Buea, Limbé et rend l'ascension du Mont Cameroun hasardeuse. \textbf{Problématique :} L'insécurité est-elle l'unique ou le principal verrou du tourisme camerounais, ou ne fait-elle qu'aggraver des faiblesses structurelles préexistantes ?

\textbf{Premier moment : L'insécurité, obstacle réel et immédiat}
L'insécurité est indéniablement un \textbf{frein majeur, voire bloquant}, pour les zones concernées. Elle agit par trois canaux : 1) \textbf{Annulation directe} : les tours-opérateurs (ex: Voyageurs du Monde, Atalante) retirent le Cameroun de leurs catalogues ; les ministères étrangers (France, USA, Allemagne) classent l'Extrême-Nord, le Nord-Ouest, le Sud-Ouest en zone rouge (déconseillé formel). 2) \textbf{Destruction de l'offre} : à Waza, le campement est pillé ; à Bamenda, l'hôtel Ayaba est fermé ; les guides de Mont Cameroun (Buea) ont fui. 3) \textbf{Effet d'image « tâche d'huile »} : le touriste international ne distingue pas les régions ; « Cameroun = danger » s'impose, pénalisant même les zones sûres (Kribi, Douala, Yaoundé, Ouest). C'est un \textbf{choc de demande} brutal.

\textbf{Deuxième moment : Des obstacles structurels profonds, antérieurs et persistants}
Pourtant, \textbf{l'insécurité ne saurait masquer les carences chroniques} qui limitent le secteur même en temps de paix relative (ex: 2010-2013, arrivées stagnantes \textasciitilde 600-900k/an).
\begin{itemize}
\item \textbf{Formation :} L'École hôtelière de Yaoundé forme trop peu ; les guides de Waza ou Korup manquent de certification internationale (FGASA) ; l'anglais, langue du tourisme mondial, est mal maîtrisé hors zone anglophone.
\item \textbf{Infrastructures :} La route vers l'Est (Bertoua-Yokadouma-Dja) est une piste ; l'accès à Lobéké/Campo Ma'an demande 2 jours de 4x4. L'aéroport de Maroua n'a pas de vol direct international. Le parc hôtelier 4/5* se compte sur les doigts d'une main (Hilton Yaoundé, Sawa Douala, Mont Fébé).
\item \textbf{Promotion / Visa :} Le budget de l'Office National du Tourisme est dérisoire face à la Côte d'Ivoire ou au Kenya. Le visa (même e-visa) reste un frein psychologique et administratif comparé aux exemptions (Maroc, Sénégal, Rwanda).
\end{itemize}
Ces facteurs expliquent pourquoi le Cameroun capte \textless 2 \% du tourisme africain alors qu'il a plus de diversité que le Kenya ou la Tanzanie.

\textbf{Conclusion}
\textbf{Position :} L'insécurité est l'\textbf{obstacle conjoncturel le plus violent} (arrêt brutal), mais les obstacles \textbf{structurels (infrastructures, formation, promotion, gouvernance)} sont la cause profonde de la sous-performance chronique. \textbf{Hiérarchie des solutions :} 1) \textbf{Urgence :} Sécurisation des corridors touristiques (Waza, Mont Cameroun, Rhumsiki) par forces armées + garde forestière + comités de vigilance locaux. 2) \textbf{Investissement :} Programme pluriannuel routes/aéroports/énergie vers les parcs (Waza, Dja, Lobéké, Korup) + réhabilitation camps/hôtels via PPP. 3) \textbf{Humain :} Refonte formation (écoles hôtelières, guides certifiés, langues). 4) \textbf{Marketing :} Cible « niche haute valeur » (écotourisme, ornithologie, culturel, diaspora) + facilitation visa (exemption CEMAC, visa à l'arrivée aéroports). Sans cette chaîne complète, la paix revenue ne remplira pas les hôtels.
\end{corrige}

\newpage

\coursec{Fiche bilan}

\courssub{Carte mentale de la leçon}

\begin{popfigure}
\begin{center}\tcbox[colback=popink,colframe=popink,coltext=white,arc=9pt,boxsep=4pt,left=6pt,right=6pt]{\bfseries\small Le tourisme au Cameroun}\end{center}
\vspace{-2pt}
\begin{tcbraster}[raster columns=3,raster equal height=rows,raster column skip=6pt,raster row skip=6pt,enhanced,blankest]
\begin{tcolorbox}[enhanced,colback=popblue,colframe=popblue,coltext=white,boxrule=0.9pt,arc=3pt,left=4pt,right=4pt,top=3pt,bottom=3pt,boxsep=1pt,halign=left]\bfseries\scriptsize Définitions \footnotesize site touristique, tourisme d'affaires\end{tcolorbox}
\begin{tcolorbox}[enhanced,colback=popgreen,colframe=popgreen,coltext=white,boxrule=0.9pt,arc=3pt,left=4pt,right=4pt,top=3pt,bottom=3pt,boxsep=1pt,halign=left]\bfseries\scriptsize Potentialités \footnotesize mer, safari, hauts plateaux, écotourisme\end{tcolorbox}
\begin{tcolorbox}[enhanced,colback=poporange,colframe=poporange,coltext=white,boxrule=0.9pt,arc=3pt,left=4pt,right=4pt,top=3pt,bottom=3pt,boxsep=1pt,halign=left]\bfseries\scriptsize Importance \footnotesize économique, politique, sociale\end{tcolorbox}
\begin{tcolorbox}[enhanced,colback=poppink,colframe=poppink,coltext=white,boxrule=0.9pt,arc=3pt,left=4pt,right=4pt,top=3pt,bottom=3pt,boxsep=1pt,halign=left]\bfseries\scriptsize Problèmes \footnotesize formation, accueil, voies, insécurité\end{tcolorbox}
\begin{tcolorbox}[enhanced,colback=popgold,colframe=popgold,coltext=white,boxrule=0.9pt,arc=3pt,left=4pt,right=4pt,top=3pt,bottom=3pt,boxsep=1pt,halign=left]\bfseries\scriptsize Perspectives \footnotesize valorisation, intégration nationale\end{tcolorbox}
\end{tcbraster}

\legende{Carte mentale de synthèse : les cinq branches du chapitre « Le tourisme au Cameroun ». Le nœud central rappelle l'objet d'étude ; chaque branche colorée correspond à une partie du plan de la leçon.}
\end{popfigure}

\courssub{L'essentiel à retenir}

\begin{aretenir}[Le tourisme au Cameroun : l'essentiel]
\textbf{1. Définitions clés}
\begin{itemize}
  \item \cle{Tourisme} : ensemble des activités de déplacement et de séjour des personnes hors de leur lieu de résidence habituel, pour les loisirs, les affaires ou d'autres motifs.
  \item \cle{Site touristique} : lieu (naturel ou culturel) qui, par son attrait particulier, attire les visiteurs et fait l'objet d'une fréquentation touristique (ex. : chutes de la Lobé à Kribi, mont Cameroun, chefferies de l'Ouest).
  \item \cle{Tourisme d'affaires} : tourisme lié aux déplacements professionnels (congrès, séminaires, salons, missions) ; il est très développé à Yaoundé et Douala, pôles politique et économique du pays.
  \item \cle{Écotourisme} : forme de tourisme responsable centrée sur la découverte de la nature et le respect de l'environnement et des populations locales.
\end{itemize}

\textbf{2. Les potentialités touristiques (à connaître par cœur)}
\begin{center}
\begin{tabularx}{\linewidth}{|l|X|}
\hline
\rowcolor{popblueL} \textbf{Atout} & \textbf{Exemples camerounais} \\
\hline
Attrait de la mer & Plages de Kribi et de Limbé, chutes de la Lobé se jetant dans l'océan Atlantique \\
\hline
Safari (septentrion) & Parcs nationaux de Waza, de la Bénoué, de Bouba Ndjida \\
\hline
Dynamisme culturel des hauts plateaux & Chefferies bamiléké et bamoun (Foumban, Bafoussam, Bandjoun), danses, artisanat, cases traditionnelles \\
\hline
Écotourisme (plateau forestier) & Forêts du Sud et de l'Est, réserve de faune du Dja (patrimoine mondial de l'UNESCO), parc national de Korup \\
\hline
Relief et paysages & Mont Cameroun (\SI{4040}{m}, point culminant du pays), lacs de cratère du mont Manengouba, monts Mandara \\
\hline
\end{tabularx}
\end{center}

\textbf{3. L'importance du tourisme}
\begin{itemize}
  \item \textbf{Économique} : devises étrangères, recettes fiscales, emplois (hôtellerie, restauration, transport, artisanat, guidage), stimulation des filières locales.
  \item \textbf{Politique} : rayonnement international du pays (« l'Afrique en miniature »), vitrine de la stabilité, outil de diplomatie culturelle.
  \item \textbf{Sociale} : valorisation des cultures et des identités locales, renforcement de l'\cle{intégration nationale} par la découverte mutuelle des régions, création d'infrastructures profitant aux populations.
\end{itemize}

\textbf{4. Les problèmes du secteur}
\begin{itemize}
  \item Déficit de \cle{formation} des acteurs (guides, hôteliers, agents d'accueil) ;
  \item Insuffisance des \cle{structures d'accueil} (hôtels de standing inégalement répartis, concentrés à Yaoundé, Douala, Kribi) ;
  \item État dégradé des \cle{voies de communication} (routes d'accès aux parcs et sites difficiles, surtout en saison des pluies) ;
  \item \cle{Insécurité} dans certaines zones (Extrême-Nord, régions du Nord-Ouest et du Sud-Ouest), braconnage dans les parcs ;
  \item Promotion touristique encore faible à l'international.
\end{itemize}

\textbf{5. Méthodes express}
\begin{itemize}
  \item \emph{Croquis} : contour simplifié du Cameroun, localiser les 4 grandes zones d'atouts (littoral, septentrion, hauts plateaux de l'Ouest, plateau forestier du Sud-Est), légende numérotée, titre, nord, échelle indicative.
  \item \emph{Commentaire de document} : identifier la nature du document, dégager l'idée principale, mobiliser les quatre atouts et les quatre problèmes, conclure sur l'intégration nationale.
\end{itemize}
\end{aretenir}

\begin{attention}[Précisions à ne pas négliger]
\begin{itemize}
  \item Le point culminant du Cameroun est le mont Cameroun, à environ \SI{4040}{m} (et non \SI{4095}{m}) : retiens l'ordre de grandeur « un peu plus de \SI{4000}{m} ».
  \item La réserve du Dja est une \emph{réserve de faune} inscrite au patrimoine mondial de l'UNESCO ; Korup est un \emph{parc national} : ne confonds pas les statuts.
  \item On écrit « monts Mandara » (massif de l'Extrême-Nord) et non « mandara » seul.
  \item Le tourisme d'affaires ne se limite pas aux loisirs : congrès, séminaires et missions professionnelles en font partie ; c'est une définition exigible au Bac.
\end{itemize}
\end{attention}

\begin{aretenir}[Checklist avant le Bac]
\begin{itemize}
  \item[$\square$] Je sais définir précisément un \cle{site touristique} et le \cle{tourisme d'affaires}.
  \item[$\square$] Je connais les quatre grandes potentialités touristiques du Cameroun et un exemple précis pour chacune.
  \item[$\square$] Je sais citer les parcs nationaux du septentrion : Waza, Bénoué, Bouba Ndjida.
  \item[$\square$] Je sais expliquer pourquoi on appelle le Cameroun « l'Afrique en miniature » (diversité des reliefs, des climats, des paysages et des cultures réunie en un seul pays).
  \item[$\square$] Je distingue les trois dimensions de l'importance du tourisme : économique, politique, sociale.
  \item[$\square$] Je connais au moins quatre problèmes du secteur touristique camerounais.
  \item[$\square$] Je sais relier le tourisme à la famille de situations « l'intégration nationale ».
  \item[$\square$] Je sais construire un croquis légendé des grandes zones touristiques du Cameroun.
  \item[$\square$] Je sais localiser Kribi, Limbé, Foumban, Waza et la réserve du Dja sur une carte muette.
  \item[$\square$] Je sais proposer des solutions concrètes pour valoriser le tourisme camerounais (formation des acteurs, aménagement des routes d'accès, promotion internationale, sécurisation des sites).
\end{itemize}
\end{aretenir}

\findecours

\end{document}
